\documentclass[10pt]{article}
\usepackage[utf8]{inputenc}
\usepackage[T5]{fontenc}
\usepackage[english]{babel}
\usepackage{amsmath}
\usepackage{amsfonts}
\usepackage{amssymb}
\usepackage[version=4]{mhchem}
\usepackage{stmaryrd}
\usepackage{bbold}
\usepackage{lmodern}
\usepackage[a4paper,top=2.2cm,bottom=2.5cm,left=2.2cm,right=2.2cm]{geometry}

% Nới lỏng quy tắc ngắt dòng để tránh tràn lề
\tolerance=2000
\emergencystretch=2em
\sloppy

% Hỗ trợ các ký tự Unicode đặc biệt nếu xuất hiện
\providecommand{\DeclareUnicodeCharacter}[2]{}
\DeclareUnicodeCharacter{2021}{\ifmmode\ddagger\else{$\ddagger$}\fi}
\DeclareUnicodeCharacter{25A1}{\ifmmode\square\else{$\square$}\fi}

% Lệnh chú thích chân trang như bản gốc
\let\svthefootnote\thefootnote
\newcommand\blfootnotetext[1]{%
  \let\thefootnote\relax\footnote{#1}%
  \addtocounter{footnote}{-1}%
  \let\thefootnote\svthefootnote%
}

\let\svfootnotetext\footnotetext
\renewcommand\footnotetext[2][?]{%
  \if\relax#1\relax%
    \ifnum\value{footnote}=0\blfootnotetext{#2}\else\svfootnotetext{#2}\fi%
  \else%
    \if?#1\ifnum\value{footnote}=0\blfootnotetext{#2}\else\svfootnotetext{#2}\fi%
    \else\svfootnotetext[#1]{#2}\fi%
  \fi
}

\title{Định nghĩa sơ đồ về Tính toán Đa thức Lượng tử}
\author{Tomoyuki Yamakam ‡}
\date{}

\begin{document}
\maketitle

% ---- Nội dung từ page_01_07_vi.tex ----
\begin{abstract}
Trong bốn thập kỷ qua, khái niệm về tính toán đa thức lượng tử đã được mô hình hóa toán học bởi các máy Turing lượng tử cũng như các mạch lượng tử. Bài báo này tìm kiếm mô hình thứ ba, đó là một tương tự lượng tử của định nghĩa sơ đồ (quy nạp hoặc xây dựng) của các hàm (nguyên thủy) đệ quy. Đối với các hàm lượng tử ánh xạ không gian Hilbert hữu hạn chiều vào chính nó, chúng tôi trình bày một định nghĩa sơ đồ như vậy, bao gồm một tập hợp nhỏ các hàm lượng tử ban đầu và một số quy tắc xây dựng quy định cách xây dựng một hàm lượng tử mới từ các hàm hiện có. Chúng tôi chứng minh rằng định nghĩa sơ đồ của chúng tôi mô tả chính xác tất cả các hàm có thể tính toán được với xác suất thành công cao trên các máy Turing lượng tử được hình thành tốt trong thời gian đa thức, hoặc tương đương là các họ đồng nhất của các mạch lượng tử kích thước đa thức. Định nghĩa sơ đồ mới của chúng tôi khá đơn giản và trực quan và, quan trọng hơn, nó tránh được việc giới thiệu phức tạp về điều kiện hình thành tốt áp đặt lên mô hình máy Turing lượng tử cũng như điều kiện đồng nhất cần thiết cho mô hình mạch lượng tử. Cách tiếp cận mới của chúng tôi có thể mở ra cánh cửa cho sự phức tạp mô tả của các hàm lượng tử, lý thuyết về các hàm lượng tử loại cao hơn, phát triển các lý thuyết bậc một mới cho tính toán lượng tử và thiết kế các ngôn ngữ lập trình cho máy tính lượng tử thực tế.
\end{abstract}

Từ khóa: tính toán lượng tử, hàm lượng tử, máy Turing lượng tử, mạch lượng tử, định nghĩa sơ đồ, độ phức tạp mô tả, tính toán đa thức, định lý dạng chuẩn

\section*{1 Bối cảnh, Động lực và Kết quả Chính}
Vào đầu những năm 1980, một ý tưởng đột phá đã xuất hiện về việc khai thác vật lý lượng tử để xây dựng các thiết bị tính toán cơ học, được gọi là máy tính lượng tử, đã hoàn toàn thay đổi cách chúng ta hình dung về "máy tính." Những khám phá tiếp theo về các tính toán lượng tử hiệu quả hơn để phân tích các số nguyên dương [30] và tìm kiếm cơ sở dữ liệu không có cấu trúc [14, 15] so với các tính toán cổ điển đã thúc đẩy chúng ta tìm kiếm thêm các vấn đề toán học và thực tiễn có thể được giải quyết hiệu quả trên các máy tính lượng tử. Tính hiệu quả trong tính toán lượng tử kể từ đó đã nhanh chóng trở thành một chủ đề nghiên cứu quan trọng của khoa học máy tính cũng như vật lý.

Như một mô hình toán học để hiện thực hóa tính toán lượng tử, Deutsch 11 đã giới thiệu khái niệm về máy Turing lượng tử (hoặc QTM, viết tắt) , sau này đã được Yao 40 thảo luận và tiếp tục được Bernstein và Vazirani [5] tinh chỉnh. Mô hình cơ học này mở rộng lớn mô hình Turing cổ điển (xác suất) bằng cách cho phép một hiện tượng vật lý, được gọi là giao thoa lượng tử, xảy ra trong quá trình tính toán của nó. Một hình thức Hamiltonian khác của các máy Turing cũng đã được Benioff [3] đề xuất. Một QTM có khả năng tính toán một hàm lượng tử ánh xạ một không gian Hilbert hữu hạn chiều vào chính nó bằng cách tiến hóa một cách đơn vị một sự superposition của các cấu hình (cổ điển) của máy, bắt đầu với một chuỗi đầu vào nhất định và một trạng thái bên trong ban đầu. Một cách sử dụng hạn chế hơn của thuật ngữ "hàm lượng tử" được tìm thấy trong, ví dụ, [38], trong đó các hàm lượng tử nhận các chuỗi đầu vào cổ điển và tạo ra hoặc là các chuỗi đầu ra cổ điển của các QTM hoặc là các xác suất chấp nhận của các QTM. Trong suốt bài báo này, tuy nhiên, các hàm lượng tử chỉ đề cập đến các hàm hoạt động trên các không gian Hilbert có kích thước tùy ý.

Để đảm bảo tính đơn vị của quá trình tính toán lượng tử, một QTM yêu cầu cơ chế của nó phải đáp ứng các điều kiện được gọi là điều kiện hình thành tốt trên một mô hình QTM băng đơn 5 và một mô hình đa băng 36, 38 cũng như [27]. Tham khảo phần 2.2 để biết định nghĩa chính xác của chúng.

Bernstein và Vazirani đã hình thành một lớp độ phức tạp mới, được ký hiệu là BQP, như là tập hợp của tất cả các ngôn ngữ được công nhận bởi các QTM được hình thành tốt chạy trong thời gian đa thức với xác suất lỗi bị giới hạn

\footnotetext{*Công việc này được thực hiện trong khi tác giả còn ở Đại học Ottawa giữa năm 1999 và 2003, và nó được hỗ trợ tài chính bởi Hội đồng Khoa học Tự nhiên và Kỹ thuật Canada.\\
${ }^{\dagger}$ Một tóm tắt mở rộng đã xuất hiện dưới tiêu đề "Một định nghĩa đệ quy về tính toán đa thức lượng tử (tóm tắt mở rộng)" trong Kỷ yếu của Hội thảo lần thứ 9 về Các Mô hình Tự động và Ứng dụng Không Cổ điển (NCMA 2017), Prague, Cộng hòa Séc, 17-18 tháng 8 năm 2017, Österreichische Computer Gesellschaft 2017, Hội máy tính Áo, trang 243-258, 2017. Bài báo hiện tại nhằm mục đích sửa chữa các mô tả sai trong tóm tắt mở rộng và cung cấp các giải thích chi tiết hơn cho tất cả các chứng minh bị bỏ qua do giới hạn trang.\\
${ }^{\ddagger}$ Đơn vị hiện tại: Khoa Kỹ thuật, Đại học Fukui, 3-9-1 Bunkyo, Fukui, 910-8507 Nhật Bản
}
từ trên xuống dưới bởi $1 / 3$. Hơn nữa, các QTM được trang bị băng xuất có thể tính toán các hàm có giá trị chuỗi thay vì các ngôn ngữ, và các hàm đó tạo thành một lớp hàm, được gọi là FBQP .

Từ một góc độ khác, Yao 40 đã mở rộng khái niệm mạng lượng tử của Deutsch [12] và hình thành một khái niệm về mạch lượng tử, đó là một tương tự lượng tử của mạch Boolean cổ điển. Khác với mô hình mạch Boolean cổ điển, một mạch lượng tử được cấu thành từ các cổng lượng tử, mỗi cổng đại diện cho một phép biến đổi đơn vị tác động lên một không gian Hilbert có kích thước nhỏ, cố định. Để hoạt động như một "bộ lập trình" phép toán đơn vị, một họ các mạch lượng tử yêu cầu điều kiện đồng nhất, đảm bảo rằng bản thiết kế của mỗi mạch lượng tử có thể được thực hiện dễ dàng. Yao đã chỉ ra rằng một họ đồng nhất của các mạch lượng tử đủ mạnh để mô phỏng một máy Turing lượng tử được hình thành tốt. Như Nishimura và Ozawa 26 đã chỉ ra, điều kiện đồng nhất của một họ mạch lượng tử là cần thiết để nắm bắt chính xác tính toán đa thức lượng tử. Với điều kiện đồng nhất này, BQP và FBQP được đặc trưng chính xác bởi các họ mạch lượng tử đồng nhất được tạo thành từ nhiều cổng lượng tử có kích thước đa thức.

Bài báo hiện tại táo bạo tiếp cận theo cách thứ ba để đặc trưng hóa tính toán đa thức lượng tử. Khác với các mô hình thiết bị cơ học đã đề cập, cách tiếp cận của chúng tôi là mở rộng định nghĩa sơ đồ (quy nạp hoặc xây dựng) của các hàm (nguyên thủy) đệ quy trên các số tự nhiên. Định nghĩa sơ đồ như vậy đã được Peano nghĩ ra vào thế kỷ 19 [28], trái ngược với định nghĩa được đưa ra bởi mô hình máy Turing của Turing 32. Sơ đồ cổ điển này bao gồm một tập hợp nhỏ các hàm ban đầu và một tập hợp nhỏ các quy tắc, quy định cách xây dựng một hàm mới từ các hàm hiện có. Ví dụ, mỗi hàm đệ quy nguyên thủy được xây dựng từ các hàm hằng số, kế tiếp, và chiếu bởi việc áp dụng hữu hạn các quy tắc hợp thành và quy tắc đệ quy nguyên thủy. Đặc biệt, quy tắc đệ quy nguyên thủy giới thiệu một hàm mới có giá trị được xác định bởi quy tắc suy diễn. Các hàm đệ quy (dưới dạng các hàm đệ quy $\mu$ [19, 20]) yêu cầu thêm một sơ đồ, được gọi là phép toán tối thiểu (hoặc số nhỏ nhất). Các hàm này trùng với hình thức chính quy của các hàm đệ quy tổng quát của Gödel-Herbrand (xem [10]). Để có cái nhìn lịch sử về những khái niệm này, xem, ví dụ, 29. Những cách tiếp cận sơ đồ tương tự để nắm bắt tính toán đa thức cổ điển đã được tìm kiếm trong tài liệu $[7,8,9,24,35$. Những cách tiếp cận đó đã dẫn đến những chủ đề nghiên cứu khá khác biệt so với những gì mô hình máy Turing cung cấp.

Mục đích của chúng tôi trong bài báo này là đưa ra một định nghĩa sơ đồ của các hàm lượng tử để nắm bắt khái niệm về tính toán đa thức lượng tử và, quan trọng hơn, để làm cho một định nghĩa như vậy đơn giản hơn và trực quan hơn cho lợi ích thực tiễn của chính chúng tôi. Định nghĩa sơ đồ của chúng tôi (Định nghĩa 3.1) bao gồm một tập hợp các hàm lượng tử ban đầu, $I$ (đơn vị), $N O T$ (phủ định của một qubit), $P H A S E_{\theta}$ (dịch pha bởi $e^{i \theta}$ ), $R O T_{\theta}$ (xoay quanh trục $x y$ bởi góc $\theta$ ), $S W A P$ (hoán đổi giữa hai qubit), và MEAS (đo lường dự đoán từng phần), cũng như các quy tắc xây dựng, bao gồm hợp thành (Compo $[\cdot, \cdot]$ ), phân nhánh (Branch $[\cdot, \cdot]$ ), và đệ quy lượng tử nhiều qubit ( $k Q R e c[\cdot, \cdot \mid \cdot]$ ). Lựa chọn của chúng tôi về các hàm lượng tử ban đầu và các quy tắc xây dựng này chủ yếu xuất phát từ một tập hợp các cổng lượng tử phổ quát đã được sử dụng trong tài liệu trước đây. Ngược lại, đệ quy lượng tử của chúng tôi thì hoàn toàn khác về bản chất so với quy tắc đệ quy được sử dụng để xây dựng các hàm đệ quy nguyên thủy. Thay vì sử dụng hàm kế tiếp để đếm ngược số lần lặp quy nạp trong quy tắc đệ quy nguyên thủy, đệ quy lượng tử sử dụng một chiến lược chia để trị để giảm số lượng qubit có thể truy cập cần thiết cho việc thực hiện một hàm lượng tử xác định. Trong khuôn khổ mới của chúng tôi, chúng tôi có thể thực hiện các phép toán đơn vị điển hình, chẳng hạn như biến đổi Walsh-Hadamard (WH), biến đổi có điều kiện-NOT (CNOT), và dịch pha toàn cục (GPS).

Một lợi ích ngay lập tức của định nghĩa sơ đồ của chúng tôi là chúng tôi có thể tránh được việc giới thiệu phức tạp về điều kiện hình thành tốt áp đặt lên mô hình QTM và điều kiện đồng nhất trên mô hình mạch lượng tử. Một lợi thế khác của các sơ đồ của chúng tôi là mỗi sơ đồ có một phép đảo ngược riêng; nghĩa là, đối với bất kỳ hàm lượng tử nào $g$ được định nghĩa bởi một trong các sơ đồ, phép đảo ngược của nó $g^{-1}$ cũng được định nghĩa bởi cùng một loại sơ đồ. Ví dụ, các phép đảo ngược của các hàm lượng tử $R O T_{\theta}$ và $k Q R e c_{t}\left[g, h, p \mid\left\{f_{s}\right\}_{s \in\{0,1\}^{k}}\right]$ được giới thiệu trong Định nghĩa 3.1 chính xác là $R O T_{-\theta}$ và $k Q R e c_{t}\left[g^{-1}, p^{-1}, h^{-1} \mid\left\{f_{s}^{-1}\right\}_{s \in\{0,1\}^{k}}\right]$, tương ứng (Định lý 3.5).

Để giải thích thêm về những đóng góp chính của chúng tôi, đã đến lúc giới thiệu một ký hiệu ngắn gọn của $\square_{1}^{\mathrm{QP}}$ (nơi □ được phát âm là "hình vuông") để chỉ tập hợp tất cả các hàm lượng tử được xây dựng từ các hàm lượng tử ban đầu và bằng một chuỗi hữu hạn các ứng dụng tuần tự của các quy tắc xây dựng. Vì phép đo từng phần (MEAS) không phải là một phép toán đơn vị, chúng tôi ký hiệu lớp thu được từ $\square_{1}^{\mathrm{QP}}$ mà không sử dụng MEAS là $\widehat{\square_{1}^{\mathrm{QP}}}$. Ngắn gọn, hãy để chúng tôi thảo luận về những khác biệt rõ ràng giữa định nghĩa sơ đồ của chúng tôi và hai hình thức đã đề cập của các hàm có thể tính toán đa thức lượng tử về mặt các QTMs và các mạch lượng tử. Hai sự khác biệt chính được liệt kê dưới đây.

\begin{enumerate}
  \item Trong khi một mạch lượng tử đơn lẻ nhận một số qubit đầu vào cố định, hàm lượng tử của chúng tôi nhận một số "tùy ý" qubit làm đầu vào. Tình huống này tương tự như các QTMs vì một QTM có một băng vô hạn và sử dụng một số ô băng tùy ý trong quá trình tính toán của nó như là không gian lưu trữ bổ sung. Ngược lại với các QTMs, một hàm $\widehat{\square_{1}^{\mathrm{QP}}}$-function được xây dựng bằng cách sử dụng cùng một số qubit như đầu vào ban đầu của nó\\
theo cách mà một mạch lượng tử có cùng số qubit đầu vào và qubit đầu ra.
  \item Hai mô hình máy tính này yêu cầu một mô tả thuật toán để chỉ định hành vi của mỗi máy; cụ thể hơn, một QTM sử dụng một hàm chuyển tiếp, mô tả thuật toán cách mỗi bước của máy tác động lên một số qubit nhất định, và một họ các mạch lượng tử sử dụng điều kiện đồng nhất của nó để thể hiện thiết kế của các cổng lượng tử trong mỗi mạch lượng tử. Không giống như hai mô hình này, không có hàm $\square_{1}^{\mathrm{QP}}$-function nào có bất kỳ cơ chế nào để lưu trữ thông tin về chính mô tả của hàm đó mà quá trình xây dựng tự nó chỉ định hành vi của hàm.
\end{enumerate}

Như một hệ quả, những khác biệt đã đề cập ở trên giúp các hàm $\square_{1}^{\mathrm{QP}}$-functions có một vị trí đặc biệt trong số tất cả các mô hình tính toán có thể đặc trưng cho tính toán đa thức lượng tử, và do đó chúng tôi mong đợi chúng sẽ đóng một vai trò quan trọng trong việc phân tích các đặc điểm của tính toán đa thức lượng tử từ một góc độ hoàn toàn khác.

Trong Phần 3.1, chúng tôi sẽ trình bày chính thức định nghĩa sơ đồ của chúng tôi về các hàm $\square_{1}^{\mathrm{QP}}$ (cũng như các hàm $\widehat{\square_{1}^{\mathrm{QP}}}$ ) và chỉ ra trong Phần 4.1 rằng $\square_{1}^{\mathrm{QP}}$ (cũng như $\widehat{\square_{1}^{\mathrm{QP}}}$ ) có thể đặc trưng cho tất cả các hàm trong FBQP. Chính xác hơn, chúng tôi khẳng định trong định lý chính (Định lý 4.1) rằng bất kỳ hàm nào từ $\{0,1\}^{*}$ đến $\{0,1\}^{*}$ trong FBQP đều có thể được đặc trưng bởi một đa thức nhất định $p$ và một hàm lượng tử nhất định $g \in \square_{1}^{\mathrm{QP}}$ theo cách mà, bằng cách sử dụng một sơ đồ mã hóa thích hợp, trong trạng thái lượng tử cuối cùng của $g$ trên các thể hiện $x$ và giới hạn thời gian $p(|x|)$, chúng tôi quan sát được một giá trị đầu ra $f(x)$ với xác suất cao. Định lý này sẽ được chia thành hai định lý phụ, các Định lý 4.2 và 4.3 Định lý phụ trước sẽ được chứng minh trong Phần 4.1 tuy nhiên, chứng minh của định lý phụ sau thì dài đến mức sẽ được hoãn lại cho đến Phần 5. Trong chứng minh này, chúng tôi sẽ xây dựng một hàm $\square_{1}^{\mathrm{QP}}$-function có thể mô phỏng hành vi của một QTM nhất định.

Lưu ý rằng, vì BQP là một trường hợp đặc biệt của FBQP, BQP cũng được đặc trưng bởi mô hình của chúng tôi. Trong chứng minh của chúng tôi về định lý đặc trưng (Định lý 4.1), chúng tôi sẽ sử dụng một kết quả chính của Bernstein và Vazirani [5] và của Yao 40 một cách rộng rãi. Trong Phần 4.2, chúng tôi sẽ áp dụng sự đặc trưng của chúng tôi, với sự trợ giúp của một QTM phổ quát [5, 26], để thu được một phiên bản lượng tử của định lý dạng chuẩn của Kleene [19, 20, trong đó có một cặp hàm và đại số đệ quy nguyên thủy phổ quát có thể mô tả hành vi của mọi hàm đệ quy.

Khác với tính toán cổ điển trên các số tự nhiên (tương đương, các chuỗi trên các bảng chữ cái hữu hạn bằng các sơ đồ mã hóa thích hợp), tính toán lượng tử là một chuỗi nhất định của các phép biến đổi của một vectơ trong một không gian Hilbert hữu hạn chiều và chúng tôi chỉ cần độ chính xác cao để xấp xỉ mỗi hàm trong FBQP bằng một vectơ như vậy. Thực tế này cho phép chúng tôi chọn một tập hợp các sơ đồ (các hàm lượng tử ban đầu và các quy tắc xây dựng) khác nhau để nắm bắt bản chất của tính toán lượng tử. Trong Phần 6.1, chúng tôi sẽ thảo luận về vấn đề này bằng một ví dụ về dạng tổng quát của biến đổi Fourier lượng tử (QFT). Biến đổi này có thể không được "tính toán chính xác" trong khuôn khổ hiện tại của chúng tôi về $\square_{1}^{\mathrm{QP}}$ nhưng chúng tôi có thể dễ dàng mở rộng $\square_{1}^{\mathrm{QP}}$ để tính toán chính xác QFT tổng quát nếu chúng tôi bao gồm một hàm lượng tử ban đầu bổ sung, chẳng hạn như CROT (xoay có điều kiện).

Liên quan đến nghiên cứu trong tương lai về chủ đề hiện tại, chúng tôi sẽ thảo luận trong Phần 6 bốn hướng nghiên cứu mới của chủ đề này. Định nghĩa sơ đồ của chúng tôi không chỉ cung cấp một cách khác để mô tả các ngôn ngữ và hàm có thể tính toán lượng tử trong thời gian đa thức mà còn một cách đơn giản để đo lường độ phức tạp "mô tả" của một ngôn ngữ và một hàm bị giới hạn cho các thể hiện có độ dài xác định. Đo lường độ phức tạp mới này sẽ hữu ích để chứng minh các tính chất cơ bản của các hàm $\widehat{\square_{1}^{\mathrm{QP}}}$-functions trong Phần 3. Ứng dụng trong tương lai của nó sẽ được thảo luận ngắn gọn trong Phần 6.2

Kleene 21, 22 đã định nghĩa các hàm đệ quy của các loại cao hơn bằng cách mở rộng các hàm đệ quy nguyên thủy đã đề cập ở trên. Một nghiên cứu tổng quát hơn về các hàm loại cao hơn đã được tiến hành trong lý thuyết độ phức tạp tính toán trong nhiều thập kỷ [8, 9, 24, 31, 35]. Theo một tinh thần tương tự, định nghĩa sơ đồ của chúng tôi cho phép chúng tôi nghiên cứu các hàm lượng tử loại cao hơn. Trong Phần 6.3, bằng cách sử dụng các hàm oracle (các oracle hàm hoặc oracle), chúng tôi sẽ định nghĩa các hàm lượng tử loại 2, có thể dẫn dắt chúng tôi đến một lĩnh vực nghiên cứu phong phú trong tương lai.

Một định nghĩa sơ đồ về cách xây dựng một hàm mục tiêu $\square_{1}^{\mathrm{QP}}$-function có thể được coi là một "chương trình" mô tả một chuỗi các hướng dẫn về các sơ đồ nào sẽ sử dụng. Do đó, việc hình thành sơ đồ của chúng tôi mở ra một cánh cửa cho một ứng dụng thực tiễn mới trong việc thiết kế các ngôn ngữ lập trình ngắn gọn để điều khiển các hoạt động của các máy tính lượng tử thực tế trong tương lai. Trong Phần 6.4, chúng tôi sẽ tranh luận ngắn gọn về một ứng dụng của định nghĩa sơ đồ đối với sự phát triển trong tương lai của "các ngôn ngữ lập trình lượng tử." Như một ứng dụng tiếp theo của định nghĩa sơ đồ của chúng tôi, chúng tôi cũng có thể xem xét một khía cạnh mới của các lý thuyết bậc nhất và các tiểu lý thuyết của chúng. Trước đây, một tương tự lượng tử của NP (lớp thời gian đa thức không xác định) và hơn thế nữa đã được tìm kiếm trong [37] với việc sử dụng các định lượng bị giới hạn trên các trạng thái lượng tử trong các không gian Hilbert hữu hạn chiều. Theo một cách tương tự, chúng tôi mong đợi rằng $\square_{1}^{\mathrm{QP}}$ sẽ phục vụ như một nền tảng cho việc giới thiệu các lý thuyết bậc nhất và các tiểu lý thuyết của chúng trên các trạng thái lượng tử trong các không gian Hilbert.

\section*{2 Các Khái niệm và Ký hiệu Cơ bản}
Chúng tôi bắt đầu bằng cách giải thích các khái niệm và ký hiệu cơ bản cần thiết để đọc các phần tiếp theo. Giả sử độc giả đã quen thuộc với các máy Turing cổ điển (xem, ví dụ, [16]). Đối với nền tảng của thông tin và tính toán lượng tử, ngược lại, độc giả tham khảo các sách giáo khoa cơ bản, ví dụ, [18, 25].

\subsection*{2.1 Số, Ngôn ngữ và Qustrings}
Ký hiệu $\mathbb{Z}$ chỉ tập hợp tất cả các số nguyên và $\mathbb{N}$ biểu thị tập hợp tất cả các số tự nhiên (tức là, các số nguyên không âm). Để tiện lợi, chúng tôi đặt $\mathbb{N}^{+}=\mathbb{N}-\{0\}$. Hơn nữa, $\mathbb{Q}$ biểu thị tập hợp tất cả các số hữu tỷ và $\mathbb{R}$ chỉ ra tập hợp tất cả các số thực. Đối với hai số $m, n \in \mathbb{Z}$ với $m \leq n$, ký hiệu $[m, n]_{\mathbb{Z}}$ biểu thị một khoảng số nguyên $\{m, m+1, m+2, \ldots, n\}$, so với một khoảng thực $[\alpha, \beta]$ cho $\alpha, \beta \in \mathbb{R}$ với $\alpha \leq \beta$. Đặc biệt, $[n]$ là viết tắt cho $[1, n]_{\mathbb{Z}}$ cho bất kỳ $n \in \mathbb{N}^{+}$. Bằng $\mathbb{C}$, chúng tôi biểu thị tập hợp tất cả các số phức. Cho $\alpha \in \mathbb{C}, \alpha^{*}$ biểu thị liên hợp phức của $\alpha$. Các đa thức được giả định có các hệ số là số tự nhiên và do đó chúng tạo ra các giá trị không âm từ các đầu vào không âm. Một số thực $\alpha$ được gọi là có thể xấp xỉ trong thời gian đa thức ${ }^{§}$ nếu tồn tại một máy Turing xác định đa băng trong thời gian đa thức $M$ (được trang bị một băng xuất chỉ ghi) mà, trên mỗi đầu vào có dạng $1^{n}$ cho một số tự nhiên $n$, sản xuất một phân số nhị phân hữu hạn, $M\left(1^{n}\right)$, trên băng xuất được chỉ định của nó với $\left|M\left(1^{n}\right)-\alpha\right| \leq 2^{-n}$. Gọi $\tilde{\mathbb{C}}$ là tập hợp các số phức mà phần thực và phần ảo của chúng đều có thể xấp xỉ trong thời gian đa thức. Đối với một bit $a \in\{0,1\}, \bar{a}$ chỉ ra $1-a$. Cho một ma trận $A, A^{T}$ biểu thị chuyển vị của nó và $A^{\dagger}$ biểu thị chuyển vị liên hợp của $A$.

Một bảng chữ cái là một tập hợp hữu hạn không rỗng các "ký hiệu" hoặc "chữ cái." Cho một bảng chữ cái như vậy $\Sigma$, một chuỗi trên $\Sigma$ là một chuỗi hữu hạn các ký hiệu lấy từ $\Sigma$. Phép nối của hai chuỗi $u$ và $w$ được biểu thị là $u \cdot w$ hoặc đơn giản hơn là $u w$. Độ dài của một chuỗi $x$, được ký hiệu là $|x|$, là số lần xuất hiện của các ký hiệu trong $x$. Đặc biệt, chuỗi rỗng có độ dài 0 và được ký hiệu là $\lambda$. Chúng tôi viết $\Sigma^{n}$ cho tập hợp con của $\Sigma^{*}$ chỉ bao gồm tất cả các chuỗi có độ dài $n$ và chúng tôi đặt $\Sigma^{*}=\bigcup_{n \in \mathbb{N}} \Sigma^{n}$ (tập hợp tất cả các chuỗi trên $\Sigma$ ). Một ngôn ngữ trên $\Sigma$ là một tập con của $\Sigma^{*}$. Cho một ngôn ngữ $S$, hàm đặc trưng của nó cũng được biểu thị bởi $S$; tức là, $S(x)=1$ cho tất cả $x \in S$ và $S(x)=0$ cho tất cả $x \notin S$. Một hàm trên $\Sigma^{*}$ (tức là, từ $\Sigma^{*}$ đến $\Sigma^{*}$ ) được giới hạn bởi đa thức nếu tồn tại một đa thức $p$ thỏa mãn $|f(x)| \leq p(|x|)$ cho tất cả các chuỗi $x \in \Sigma^{*}$.

Đối với mỗi số tự nhiên $k \geq 1, \mathcal{H}_{k}$ biểu thị một không gian Hilbert có kích thước $k$ và mỗi phần tử của $\mathcal{H}_{k}$ được biểu thị dưới dạng $|\phi\rangle$ bằng cách sử dụng ký hiệu "ket" của Dirac. Trong bài báo này, chúng tôi chỉ quan tâm đến trường hợp mà $k$ là một lũy thừa của 2 và chúng tôi ngầm giả định rằng $k$ có dạng $2^{n}$ cho một số nhất định $n \in \mathbb{N}$. Bất kỳ phần tử nào của $\mathcal{H}_{2}$ có chuẩn đơn vị được gọi là một bit lượng tử hoặc một qubit. Bằng cách chọn một cơ sở tính toán chuẩn $B_{1}=\{|0\rangle,|1\rangle\}$, mỗi qubit $|\phi\rangle$ có thể được biểu thị dưới dạng $\alpha_{0}|0\rangle+\alpha_{1}|1\rangle$ cho một sự lựa chọn thích hợp của hai giá trị $\alpha_{0}, \alpha_{1} \in \mathbb{C}$ (được gọi là biên độ) thỏa mãn $\left|\alpha_{0}\right|^{2}+\left|\alpha_{1}\right|^{2}=1$. Chúng tôi cũng biểu thị $|\phi\rangle$ dưới dạng một vectơ cột có dạng $\binom{\alpha_{0}}{\alpha_{1}}$; đặc biệt, $|0\rangle=\binom{1}{0}$ và $|1\rangle=\binom{0}{1}$. Trong trường hợp tổng quát hơn của $n \geq 1$, chúng tôi sử dụng $B_{n}=\left\{|s\rangle \mid s \in\{0,1\}^{n}\right\}$ làm cơ sở tính toán của $\mathcal{H}_{2^{n}}$ với $\left|B_{n}\right|=2^{n}$. Cho bất kỳ số nào $n \in \mathbb{N}^{+}$, một qustring có độ dài $n$ là một vectơ $|\phi\rangle$ của $\mathcal{H}_{2^{n}}$ có chuẩn đơn vị; nghĩa là, nó có dạng $\sum_{s \in\{0,1\}^{n}} \alpha_{s}|s\rangle$, trong đó mỗi biên độ $\alpha_{s}$ nằm trong $\mathbb{C}$ với $\sum_{s \in\{0,1\}^{n}}\left|\alpha_{s}\right|^{2}=1$. Lưu ý rằng một qubit là một qustring có độ dài 1 . Ngoại lệ là vectơ null, được ký hiệu đơn giản là $\mathbf{0}$, có chuẩn 0 . Mặc dù vectơ null có thể là một qustring của độ dài "tùy ý" $n$, chúng tôi thay vào đó gọi nó là qustring có độ dài 0 để tiện lợi. Chúng tôi sử dụng ký hiệu $\Phi_{n}$ cho mỗi $n \in \mathbb{N}$ để chỉ tập hợp tất cả các qustring có độ dài $n$. Cuối cùng, chúng tôi đặt $\Phi_{\infty}=\bigcup_{n \in \mathbb{N}} \Phi_{n}$ (tập hợp tất cả các qustring).

Khi $s=s_{1} s_{2} \cdots s_{n}$ với $s_{i} \in\{0,1\}$ cho bất kỳ chỉ số nào $i \in[n]$, qustring $|s\rangle$ trùng với $\left|s_{1}\right\rangle \otimes \left|s_{2}\right\rangle \otimes \cdots \otimes\left|s_{n}\right\rangle$, trong đó $\otimes$ biểu thị tích tensor và được biểu thị như, ví dụ, $|00\rangle=\left(\begin{array}{lll}1 & 0 & 0\end{array}\right)^{T}$, $|01\rangle=\left(\begin{array}{llll}0 & 1 & 0 & 0\end{array}\right)^{T}$, và $|11\rangle=\left(\begin{array}{llll}0 & 0 & 0 & 1\end{array}\right)^{T}$. Liên hợp transposed của $|s\rangle$ được ký hiệu là $\langle s|$ (với ký hiệu "bra"). Ví dụ, nếu $|\phi\rangle=\alpha|01\rangle+\beta|10\rangle$, thì $\langle\phi|=\alpha^{*}\langle 01|+\beta^{*}\langle 10|$. Tích vô hướng của $|\phi\rangle$ và $|\psi\rangle$ được biểu thị là $\langle\phi \mid \psi\rangle$ và chuẩn của $|\phi\rangle$ do đó là $\sqrt{\langle\phi \mid \psi\rangle}$. Khi chúng tôi quan sát hoặc đo lường $|\phi\rangle$ trong cơ sở tính toán $B_{n}$, chúng tôi thu được mỗi chuỗi $s \in\{0,1\}^{n}$ với xác suất $|\langle s \mid \phi\rangle|^{2}$.

Gọi $\mathcal{H}_{\infty}=\bigcup_{n \in \mathbb{N}^{+}} \mathcal{H}_{2^{n}}$. Chúng tôi mở rộng khái niệm "độ dài" cho các trạng thái lượng tử tùy ý trong $\mathcal{H}_{\infty}$. Đối với mỗi vectơ không null $|\phi\rangle$ trong $\mathcal{H}_{\infty}$, độ dài của $|\phi\rangle$, được ký hiệu là $\ell(|\phi\rangle)$, là số nguyên tối thiểu $n \in \mathbb{N}$ thỏa mãn $|\phi\rangle \in \mathcal{H}_{2^{n}}$; nói cách khác, $\ell(|\phi\rangle)$ là logarit của kích thước của vectơ $|\phi\rangle$. Theo quy ước, chúng tôi cũng đặt $\ell(\mathbf{0})=\ell(\alpha)=0$ cho vectơ null $\mathbf{0}$ và bất kỳ số vô hướng nào $\alpha \in \mathbb{C}$. Với quy ước này, nếu $\ell(|\phi\rangle)=0$ cho một trạng thái lượng tử $|\phi\rangle$, thì $|\phi\rangle$ nhất định phải là qustring có độ dài 0 . Một qustring $|\phi\rangle$ có độ dài $n$ được gọi là cơ bản nếu $|\phi\rangle=|s\rangle$ cho một chuỗi nhị phân nhất định $s$ và chúng tôi thường xác định một qustring cơ bản như vậy $|s\rangle$ với chuỗi nhị phân cổ điển tương ứng $s$ vì sự tiện lợi. Vì tất cả các qustring cơ bản trong $\Phi_{n}$ tạo thành $B_{n}, \mathcal{H}_{2^{n}}$ được tạo thành bởi tất cả các phần tử trong $\Phi_{n}$.

\footnotetext{${ }^{§}$ Ko và Friedman 23 đã lần đầu tiên giới thiệu khái niệm này dưới cái tên "có thể tính toán trong thời gian đa thức." Để tránh sự nhầm lẫn của độc giả trong bài báo này, chúng tôi thích sử dụng thuật ngữ "xấp xỉ trong thời gian đa thức."
}Phép dấu chấm phẩy trên một hệ thống $B$ của một hệ thống tổ hợp $A B$, được ký hiệu là $t r_{B}$, là một toán tử lượng tử mà $\operatorname{tr}_{B}(|\phi\rangle\langle\phi|)$ là một vectơ thu được bằng cách loại bỏ $B$ từ sản phẩm ngoài $|\phi\rangle\langle\phi|$ của một trạng thái lượng tử $|\phi\rangle$. Liên quan đến một trạng thái lượng tử $|\phi\rangle$ của $n$ qubit, chúng tôi sử dụng ký hiệu tiện lợi $\operatorname{tr}_{k}(|\phi\rangle\langle\phi|)$ để chỉ trạng thái lượng tử thu được từ $|\phi\rangle$ bằng cách loại bỏ tất cả các qubit ngoại trừ $k$ qubit đầu tiên. Ví dụ, nó theo cho $\sigma_{1}, \sigma_{2}, \tau_{1}, \tau_{2} \in\{0,1\}$ rằng $\operatorname{tr}_{1}\left(\left|\sigma_{1}\right\rangle\left\langle\sigma_{2}\right| \otimes\left|\tau_{1}\right\rangle\left\langle\tau_{2}\right|\right)=\left|\sigma_{1}\right\rangle\left\langle\sigma_{2}\right| \cdot \operatorname{tr}\left(\left|\tau_{1}\right\rangle\left\langle\tau_{2}\right|\right)$, trong đó $\operatorname{tr}(B)$ biểu thị dấu chấm của một ma trận $B$. Chuẩn dấu chấm $\|A\|_{\mathrm{tr}}$ của một ma trận vuông $A$ được định nghĩa bởi $\|A\|_{\mathrm{tr}}=\operatorname{tr}\left(\sqrt{A A^{\dagger}}\right)$. Khoảng cách biến thiên tổng thể giữa hai tập hợp $p=\left\{p_{i}\right\}_{i \in A}$ và $q=\left\{q_{i}\right\}_{i \in A}$ của các số thực trên một tập hợp chỉ số hữu hạn $A$ là $\frac{1}{2}\|p-q\|_{1}=\frac{1}{2} \sum_{i \in A}| | p_{i}\left|-\left|q_{i}\right|\right|$.

Trong suốt bài báo này, chúng tôi có những quy ước đặc biệt liên quan đến ba ký hiệu, $|\cdot\rangle, \otimes$, và $\|\cdot\|$, lần lượt biểu thị các trạng thái lượng tử, tích tensor, và chuẩn $\ell_{2}$-norm. Những quy ước này hơi khác so với các quy ước tiêu chuẩn được sử dụng trong, ví dụ, 25, nhưng chúng làm cho các mô tả toán học của chúng tôi trong các phần sau đơn giản và ngắn gọn hơn.\\
Các Quy ước Ký hiệu: Chúng tôi viết tắt $|\phi\rangle \otimes|\psi\rangle$ là $|\phi\rangle|\psi\rangle$ cho bất kỳ hai vectơ $|\phi\rangle$ và $|\psi\rangle$. Đối với hai chuỗi nhị phân $s$ và $t,|s t\rangle$ có nghĩa là $|s\rangle \otimes|t\rangle$ hoặc $|s\rangle|t\rangle$. Giả sử $k$ và $n$ là hai số nguyên với $0<k<n$. Bất kỳ qustring nào $|\phi\rangle$ có độ dài $n$ được biểu thị chung là $|\phi\rangle=\sum_{s:|s|=k}|s\rangle\left|\phi_{s}\right\rangle$, trong đó mỗi $\left|\phi_{s}\right\rangle$ là một qustring có độ dài $n-k$. Qustring $\left|\phi_{s}\right\rangle$ này có thể được coi là kết quả của việc áp dụng một phép đo dự đoán từng phần lên $k$ qubit đầu tiên của $|\phi\rangle$, và do đó có thể biểu thị ngắn gọn là $\langle s \mid \phi\rangle$. Với ký hiệu mới, thuận tiện này, $|\phi\rangle$ trùng với $\sum_{s:|s|=k}|s\rangle \otimes\langle s \mid \phi\rangle$, được đơn giản hóa thành $\sum_{s:|s|=k}|s\rangle\langle s \mid \phi\rangle$. Lưu ý rằng, khi $k=n,\langle s \mid \phi\rangle$ là một số vô hướng, giả sử, $\alpha$ trong $\mathbb{C}$. Do đó, $|s\rangle \otimes\langle s \mid \phi\rangle$ là $|s\rangle \otimes \alpha$ và nó được coi là một vectơ cột $\alpha|s\rangle$; tương tự, chúng tôi xác định $\alpha \otimes|s\rangle$ với $\alpha|s\rangle$. Trong những trường hợp này, $\otimes$ chỉ được coi là phép nhân vô hướng. Do đó, sự bình đẳng $|\phi\rangle=\sum_{s:|s|=k}|s\rangle\langle s \mid \phi\rangle$ vẫn đúng ngay cả khi $k=n$. Liên quan đến vectơ null $\mathbf{0}$, chúng tôi cũng có quy định đặc biệt sau đây: đối với bất kỳ vectơ nào $|\phi\rangle \in \mathcal{H}_{\infty}$, (i) $0 \otimes|\phi\rangle=|\phi\rangle \otimes 0=\mathbf{0}$, (ii) $|\phi\rangle \otimes \mathbf{0}=\mathbf{0} \otimes|\phi\rangle=\mathbf{0}$, và (iii) khi $|\psi\rangle$ là vectơ null, $\langle\phi \mid \psi\rangle=\langle\psi \mid \phi\rangle=0$. Liên quan đến những quy ước về phép đo dự đoán từng phần $\langle\phi \mid \psi\rangle$, chúng tôi cũng mở rộng việc sử dụng quy ước về chuẩn $\|\cdot\|$ cho các số vô hướng. Khi $\ell(|\phi\rangle)=\ell(|\psi\rangle),\|\langle\phi \mid \psi\rangle\|$ biểu thị giá trị tuyệt đối $|\langle\phi \mid \psi\rangle|$; nói chung hơn, đối với bất kỳ số nào $\alpha \in \mathbb{C},\|\alpha\|$ có nghĩa là $|\alpha|$. Với những quy ước bổ sung này, khi $|\phi\rangle$ có dạng $\sum_{s:|s|=k}|s\rangle\langle s \mid \phi\rangle$, phương trình $\|||\phi\rangle\left\|^{2}=\sum_{s:|s|=k}\right\|\langle s \mid \phi\rangle \|^{2}$ luôn đúng cho bất kỳ chỉ số $k \in[n]$.

\subsection*{2.2 Máy Turing Lượng Tử}
Chúng tôi giả định rằng độc giả đã có kiến thức cơ bản về khái niệm máy Turing lượng tử (hoặc QTM) được định nghĩa trong 5. Như đã làm trong 36, chúng tôi cho phép một QTM được trang bị nhiều băng và di chuyển các đầu băng của nó không đồng thời sang bên phải hoặc bên trái, hoặc làm cho các đầu băng đứng yên. Một QTM như vậy cũng đã được thảo luận ở nơi khác (ví dụ, 27) và được biết là tương đương với mô hình được đề xuất trong [5].

Để tính toán các hàm từ $\Sigma^{*}$ đến $\Sigma^{*}$ trên một bảng chữ cái cho trước $\Sigma$, chúng tôi thường giới thiệu các QTM như là những máy được trang bị các băng xuất trên đó các chuỗi xuất được ghi theo một cách nhất định bởi thời gian mà các máy dừng lại. Bằng cách xác định các ngôn ngữ với các hàm đặc trưng của chúng, các QTM như vậy có thể được coi là các bộ chấp nhận ngôn ngữ.

Một cách chính thức, một máy Turing lượng tử $k$ băng (được gọi là QTM $k$ băng), với $k \in \mathbb{N}^{+}$, là một bộ sáu $\left(Q, \Sigma, \Gamma_{1} \times \cdots \times \Gamma_{k}, \delta, q_{0}, Q_{f}\right)$, trong đó $Q$ là một tập hợp hữu hạn các trạng thái bên trong bao gồm trạng thái ban đầu $q_{0}$ và một tập hợp các trạng thái cuối $Q_{f}$ với $Q_{f} \subseteq Q$, mỗi $\Gamma_{i}$ là một bảng chữ cái được sử dụng cho băng $i$ với một ký hiệu trống được phân biệt \# thỏa mãn $\Sigma \subseteq \Gamma_{1}$, và $\delta$ là một hàm chuyển tiếp lượng tử từ $Q \times \tilde{\Gamma}^{(k)} \times Q \times \tilde{\Gamma}^{(k)} \times\{L, N, R\}^{k}$ đến $\mathbb{C}$, trong đó $\tilde{\Gamma}^{(k)}$ đại diện cho $\Gamma_{1} \times \cdots \times \Gamma_{k}$. Để tiện lợi, chúng tôi xác định $L, N$, và $R$ với $-1,0$, và +1 , tương ứng, và chúng tôi đặt $D=\{0, \pm 1\}$. Để biết thêm thông tin, tham khảo [36].

Tất cả các ô băng của mỗi băng được đánh số tuần tự bởi các số nguyên. Ô được đánh số 0 trên mỗi băng được gọi là ô bắt đầu. Vào đầu quá trình tính toán, $M$ ở trong trạng thái bên trong $q_{0}$, tất cả các băng ngoại trừ băng đầu vào đều trống, và tất cả các đầu băng đang quét các ô bắt đầu. Một chuỗi đầu vào nhất định $x_{1} x_{2} \cdots x_{n}$ ban đầu được ghi trên băng đầu vào theo cách mà, đối với mỗi chỉ số $i \in[n], x_{i}$ nằm trong ô $i$ (không phải ô $i-1$ ). Khi $M$ vào một trạng thái cuối, đầu ra của $M$ là nội dung của chuỗi được ghi trên một băng xuất (nếu $M$ chỉ có một băng, thì băng xuất là cùng một băng được sử dụng để giữ các đầu vào) từ ô bắt đầu, kéo dài sang bên phải cho đến ký hiệu trống đầu tiên. Một cấu hình của $M$ được biểu thị dưới dạng một bộ ba $\left(p,\left(h_{i}\right)_{i \in[k]},\left(z_{i}\right)_{i \in[k]}\right)$, chỉ ra rằng $M$ hiện đang ở trong trạng thái bên trong $p$ với $k$ đầu băng tại các ô được đánh số bởi $h_{1}, \ldots, h_{k}$ với nội dung băng $z_{1}, \ldots, z_{k}$, tương ứng. Khái niệm cấu hình sẽ được sửa đổi một chút trong các Phần 45 để đơn giản hóa chứng minh cho định lý chính của chúng tôi. Một cấu hình ban đầu có dạng ( $q_{0}, 0, x$ ) và một cấu hình cuối là một cấu hình có trạng thái cuối. Không gian cấu hình được tạo thành bởi các vectơ cơ sở trong $\left\{|q, h, z\rangle \mid q \in Q, h \in \mathbb{Z}^{k}, z \in \Gamma_{1}^{*} \times \cdots \Gamma_{k}^{*}\right\}$. Đối với một chuỗi không rỗng $z_{i}$ và một chỉ số $h \in\left[\left|z_{i}\right|\right], z_{i}[h]$ biểu thị ký hiệu $h$ trong $z_{i}$. Ví dụ, nếu $z_{i}=01101$, thì $z_{i}[1]=0, z_{i}[2]=1$, và $z_{i}[5]=1$. Toán tử tiến hóa theo thời gian $U_{\delta}$ của $M$ tác động lên không gian cấu hình được sinh ra từ $\delta$ như sau

$$
U_{\delta}|p, h, z\rangle=\sum_{q, w, d} \delta\left(p, z_{h}, q, z_{h}^{\prime}, d\right)\left|q, h_{d}, z^{\prime}\right\rangle,
$$

trong đó $p \in Q$, $h=\left(h_{i}\right)_{i \in[k]} \in \mathbb{Z}^{k}, z=\left(z_{i}\right)_{i \in[k]} \in \Gamma_{1}^{*} \times \cdots \times \Gamma_{k}^{*}, z_{h}=\left(z_{i}\left[h_{i}\right]\right)_{i \in[k]}, h_{d}=\left(h_{i}+d_{i}\right)_{i \in[k]}$, và $z^{\prime}=\left(z_{i}^{\prime}\right)_{i \in[k]}$, trong đó mỗi $z_{i}^{\prime}$ giống như $z$ ngoại trừ ký hiệu ở vị trí $h_{i}$. Hơn nữa, trong tổng, các biến $q, z_{h}^{\prime}=\left(z_{i}^{\prime}\left[h_{i}\right]\right)_{i \in[k]}$, và $d=\left(d_{i}\right)_{i \in[k]}$ lần lượt thay đổi trên $Q, \tilde{\Gamma}^{(k)}$, và $D^{k}$. Bất kỳ mục nào của $U_{\delta}$ được gọi là một biên độ. Cơ học lượng tử yêu cầu toán tử tiến hóa theo thời gian $U_{\delta}$ của QTM phải là đơn vị.

Mỗi bước của $M$ bao gồm hai giai đoạn: đầu tiên áp dụng $\delta$ và sau đó thực hiện một phép đo dự đoán từng phần, trong đó chúng tôi kiểm tra xem $M$ có đang ở trong trạng thái cuối hay không (tức là, một trạng thái bên trong thuộc $Q_{f}$ ). Một cách hình thức, một quá trình tính toán của $M$ trên đầu vào $x$ là một chuỗi các superposition của các cấu hình được tạo ra bởi các ứng dụng tuần tự của $U_{\delta}$, bắt đầu từ một cấu hình ban đầu của $M$ trên $x$. Nếu $M$ vào một trạng thái cuối dọc theo một đường dẫn tính toán, đường dẫn tính toán này sẽ kết thúc; nếu không, quá trình tính toán của nó phải tiếp tục.

Một QTM $k$ băng $M=\left(Q, \Sigma, \tilde{\Gamma}^{(k)}, \delta, q_{0}, Q_{f}\right)$ được hình thành tốt nếu $\delta$ thỏa mãn ba điều kiện cục bộ: độ dài đơn vị, tách biệt, và trực giao. Để giải thích các điều kiện này, như đã trình bày trong [36, Lemma 1], trước tiên chúng tôi giới thiệu các ký hiệu sau. Để tiện lợi, chúng tôi đặt $E=\{0, \pm 1, \pm 2\}$ và $H=\{0, \pm 1, \natural\}$. Cho các phần tử $(p, \sigma, \tau) \in Q \times\left(\tilde{\Gamma}^{(k)}\right)^{2}, \epsilon=\left(\varepsilon_{i}\right)_{i \in[k]} \in E^{k}$, và $d=\left(d_{i}\right)_{i \in[k]} \in D^{k}$, chúng tôi định nghĩa $D_{\epsilon}=\left\{d \in D^{k} \mid \forall i \in\right. \left.[k]\left(\left|2 d_{i}-\varepsilon_{i}\right| \leq 1\right)\right\}$ và $E_{d}=\left\{\varepsilon \in E^{k} \mid d \in D_{\epsilon}\right\}$. Hơn nữa, đặt $h_{d, \epsilon}=\left(h_{d_{i}, \varepsilon_{i}}\right)_{i \in[k]}$, trong đó $h_{d_{i}, \varepsilon_{i}}=2 d_{i}-\varepsilon_{i}$ nếu $\varepsilon_{i} \neq 0$ và $h_{d_{i}, \varepsilon_{i}}=\natural$ trong trường hợp ngược lại. Cuối cùng, chúng tôi định nghĩa $\delta(p, \sigma)=\sum_{q, \tau, d} \delta(p, \sigma, q, \tau, d)|q, \tau, d\rangle$ và $\delta[p, \sigma, \tau \mid \epsilon]= \sum_{q \in Q} \sum_{d \in D_{\epsilon}} \delta(p, \sigma, q, \tau, d)\left|E_{d}\right|^{-1 / 2}|q\rangle\left|h_{d, \epsilon}\right\rangle$, trong đó $\sigma, \tau \in \tilde{\Gamma}^{(k)}$ và $d \in D^{k}$.

\begin{enumerate}
  \item (độ dài đơn vị) $\|\delta(p, \sigma)\|=1$ cho tất cả $(p, \sigma) \in Q \times \tilde{\Gamma}^{(k)}$.
  \item (trực giao) $\delta\left(p_{1}, \sigma_{1}\right) \cdot \delta\left(p_{2}, \sigma_{2}\right)=0$ cho bất kỳ cặp khác biệt nào $\left(p_{1}, \sigma_{1}\right),\left(p_{2}, \sigma_{2}\right) \in Q \times \tilde{\Gamma}^{(k)}$.
  \item (tách biệt) $\delta\left[p_{1}, \sigma_{1}, \tau_{1} \mid \epsilon\right] \cdot \delta\left[p_{2}, \sigma_{2}, \tau_{2} \mid \epsilon^{\prime}\right]=0$ cho bất kỳ cặp khác biệt nào $\epsilon, \epsilon^{\prime} \in E^{k}$ và cho bất kỳ cặp $\left(p_{1}, \sigma_{1}, \tau_{1}\right),\left(p_{2}, \sigma_{2}, \tau_{2}\right) \in Q \times\left(\tilde{\Gamma}^{(k)}\right)^{2}$.\\
Điều kiện hình thành tốt của một QTM nắm bắt tính đơn vị của toán tử tiến hóa theo thời gian của nó.\\[0pt]
Định lý 2.1 (Định lý Hình thành tốt của [36]) Một QTM $k$ băng với một hàm chuyển tiếp $\delta$ được hình thành tốt nếu và chỉ nếu toán tử tiến hóa theo thời gian của $M$ bảo toàn chuẩn $\ell_{2}$.
\end{enumerate}

Cho một tập hợp không rỗng $K$ của $\mathbb{C}$, chúng tôi nói rằng một QTM có biên độ trong $K$ nếu tất cả các giá trị của hàm chuyển tiếp lượng tử của nó thuộc về $K$. Việc giới hạn lựa chọn biên độ trong một tập hợp hợp lý $K$ là rất quan trọng. Với việc sử dụng một tập hợp như vậy $K$, chúng tôi giới thiệu hai lớp độ phức tạp quan trọng $\mathrm{BQP}_{K}$ và $\mathrm{FBQP}_{K}$.

Định nghĩa 2.2 Giả sử $K$ là bất kỳ tập hợp không rỗng nào của $\mathbb{C}$ và $\Sigma$ là bất kỳ bảng chữ cái nào.

\begin{enumerate}
  \item Một tập hợp con $S$ của $\Sigma^{*}$ thuộc về $\mathrm{BQP}_{K}$ nếu tồn tại một QTM đa băng, thời gian đa thức, được hình thành tốt $M$ với biên độ trong $K$ sao cho, đối với mọi chuỗi $x, M$ xuất ra $S(x)$ với xác suất ít nhất $2 / 3[5]$.
  \item Một hàm đơn giá trị $f$ từ $\Sigma^{*}$ đến $\Sigma^{*}$ được gọi là có thể tính toán trong thời gian đa thức với lỗi giới hạn nếu tồn tại một QTM đa băng, thời gian đa thức, được hình thành tốt $M$ với biên độ trong $K$ sao cho, trên mọi đầu vào $x, M$ xuất ra $f(x)$ với xác suất ít nhất $2 / 3$. Gọi $\mathrm{FBQP}_{K}$ là tập hợp tất cả các hàm như vậy 38.
\end{enumerate}

Việc sử dụng các biên độ phức tạp tùy ý hóa ra làm cho $\mathrm{BQP}_{K}$ trở nên mạnh mẽ. Như Adleman, DeMarrais, và Huang 1 đã chứng minh, $\mathrm{BQP}_{\mathbb{C}}$ chứa tất cả các ngôn ngữ có thể có, và do đó $\mathrm{BQP}_{\mathbb{C}}$ không còn là đệ quy nữa. Do đó, chúng tôi thường chỉ chú ý đến các biên độ có thể xấp xỉ trong thời gian đa thức và, vì lý do này, khi $K=\tilde{\mathbb{C}}$, chúng tôi luôn bỏ qua chỉ số $K$ và viết tắt BQP và FBQP thay cho $\mathrm{BQP}_{K}$ và $\mathrm{FBQP}_{K}$, tương ứng. Cũng có thể giới hạn tập hợp biên độ $K$ hơn nữa thành $\left\{0, \pm 1, \pm \frac{3}{5}, \pm \frac{4}{5}\right\}$ vì $\mathrm{BQP}=\mathrm{BQP}_{\left\{0, \pm 1, \pm \frac{3}{5}, \pm \frac{4}{5}\right\}}$ đúng 1 .

\subsection*{2.3 Mạch Lượng Tử}
Một cổng lượng tử $k$-qubit, với $k \in \mathbb{N}^{+}$, là một phép toán đơn vị tác động lên một không gian Hilbert có kích thước $2^{k}$. Vì bất kỳ trạng thái lượng tử nào cũng là một vectơ trong một không gian Hilbert nhất định, nên mỗi mục của trạng thái lượng tử như vậy thường được gọi là một biên độ. Các phép toán đơn vị, chẳng hạn như biến đổi Walsh-Hadamard (WH) và biến đổi có điều kiện-NOT (CNOT) được định nghĩa như sau

$$
W H=\frac{1}{\sqrt{2}}\left(\begin{array}{cc}
1 & 1 \\
1 & -1
\end{array}\right) \text { và } C N O T=\left(\begin{array}{cccc}
1 & 0 & 0 & 0 \\
0 & 1 & 0 & 0 \\
0 & 0 & 0 & 1 \\
0 & 0 & 1 & 0
\end{array}\right)
$$

là các cổng lượng tử điển hình tác động lên 1 qubit và 2 qubit, tương ứng. Nếu một cổng lượng tử $U$ tác động lên $k$ qubit được áp dụng cho một trạng thái lượng tử $k$-qubit $|\phi\rangle$, thì chúng tôi thu được một trạng thái lượng tử mới $U|\phi\rangle$. Lưu ý rằng mỗi cổng lượng tử đều bảo toàn chuẩn của bất kỳ trạng thái lượng tử nào được cung cấp dưới dạng đầu vào. Một mạch lượng tử là một tích của một số hữu hạn các lớp, trong đó mỗi lớp là một tích Kronecker của các cổng lượng tử được phép. Chúng tôi thường tập trung vào một tập hợp cụ thể các cổng lượng tử để xây dựng các mạch lượng tử. Hãy xem xét các cổng lượng tử cụ thể: cổng CNOT và ba cổng một qubit có dạng

$$
Z_{1, \theta}=\left(\begin{array}{cc}
e^{i \theta} & 0 \\
0 & 1
\end{array}\right), \quad Z_{2, \theta}=\left(\begin{array}{cc}
1 & 0 \\
0 & e^{i \theta}
\end{array}\right), \quad \text { và } R_{\theta}=\left(\begin{array}{rr}
\cos \theta & -\sin \theta \\
\sin \theta & \cos \theta
\end{array}\right)
$$

trong đó $\theta$ là một số thực với $0 \leq \theta \leq 2 \pi$. Lưu ý rằng $W H$ bằng $R_{\frac{\pi}{4}}$. Những cổng này tạo thành một tập hợp cổng lượng tử phổ quát [2] vì $W H$ và $Z_{2, \frac{\pi}{4}}$ (được gọi là cổng $\pi / 8$) có thể xấp xỉ bất kỳ phép toán đơn vị một qubit nào với độ chính xác tùy ý. Vì sự tiện lợi, chúng tôi gọi chúng là các cổng cơ bản. Tập hợp các cổng CNOT, WH, và $Z_{2, \frac{\pi}{4}}$ cũng được biết là phổ quát 6.

Cho một tập hợp biên độ $K$, một mạch lượng tử $C$ được cho là có biên độ trong $K$ nếu tất cả các mục của mỗi cổng lượng tử được sử dụng bên trong $C$ đều được lấy từ $K$. Đối với bất kỳ cổng lượng tử $k$-qubit và bất kỳ số nguyên nào $n>k, G^{(n)}$ biểu thị $G \otimes I^{\otimes n-k}$, sự mở rộng $n$-qubit của $G$. Một mạch lượng tử $n$-qubit được định nghĩa một cách chính thức là một chuỗi hữu hạn $\left(G_{m}, \pi_{m}\right),\left(G_{m-1}, \pi_{m-1}\right), \cdots,\left(G_{1}, \pi_{1}\right)$ sao cho mỗi $G_{i}$ là một cổng lượng tử $n_{i}$-qubit với $n_{i} \leq n$ và $\pi_{i}$ là một hoán vị trên $\{1,2, \ldots, n\}$. Mạch lượng tử này đại diện cho phép toán đơn vị $U=U_{m} U_{m-1} \cdots U_{1}$, trong đó $U_{i}$ có dạng $V_{\pi_{i}}^{\dagger} G_{i}^{(n)} V_{\pi_{i}}$ và $V_{\pi_{i}}\left(\left|x_{1} \cdots x_{n}\right\rangle\right)=\left|x_{\pi_{i}(1)} \cdots x_{\pi_{i}(n)}\right\rangle$ cho mỗi $i \in[m]$. Kích thước của một mạch lượng tử là tổng số cổng lượng tử trong đó. Yao 40 và sau này Nishimura và Ozawa 26 đã chỉ ra rằng, đối với bất kỳ QTM $k$ băng và một đa thức $p$, tồn tại một họ các mạch lượng tử có kích thước $O\left(p(n)^{k+1}\right)$ chính xác mô phỏng $M$.

Một họ $\left\{C_{n}\right\}_{n \in \mathbb{N}}$ của một mạch lượng tử được gọi là đồng nhất P nếu tồn tại một máy Turing (cổ điển) xác định mà, trên đầu vào $1^{n}$, tạo ra một mã của $C_{n}$ trong thời gian đa thức kích thước của $C_{n}$, với điều kiện chúng tôi sử dụng một sơ đồ mã hóa cố định, hiệu quả để mô tả mỗi mạch lượng tử.

Định lý 2.3 [40 (xem thêm [26]) Đối với bất kỳ ngôn ngữ nào $L$ trên một bảng chữ cái $\{0,1\}, L$ thuộc về BQP nếu và chỉ nếu tồn tại một đa thức $p$ và một họ đồng nhất P $\left\{C_{n}\right\}_{n \in \mathbb{N}}$ của các mạch lượng tử có biên độ $\tilde{\mathbb{C}}$-amplitudes sao cho $\left\|\langle L(x)| C_{|x|}\left|x 10^{p(|x|)}\right\rangle\right\|^{2} \geq \frac{2}{3}$ đúng cho tất cả $x \in\{0,1\}^{*}$, trong đó $L$ được coi là hàm đặc trưng của $L$.

Chứng minh đầy cảm hứng của Yao 40 về Định lý 2.3 cung cấp nền tảng cho chứng minh của chúng tôi về Định lý 4.3 , điều này cung cấp trong Phần 3.1 một sự mô phỏng của một QTM được hình thành tốt bởi một hàm $\square_{1}^{\text {QP }}$-function được chọn một cách thích hợp.

\section*{3 A New, Simple Schematic Definition}
Như đã lưu ý trong Phần 1 , định nghĩa "sơ đồ" của hàm đệ quy có nghĩa là một cách quy nạp (hoặc xây dựng) để định nghĩa tập hợp các hàm có thể tính toán và nó liên quan đến một tập hợp nhỏ các hàm được gọi là các hàm ban đầu cũng như một tập hợp nhỏ các quy tắc xây dựng, được áp dụng tuần tự hữu hạn nhiều lần để xây dựng các hàm phức tạp hơn từ một số hàm đã được xây dựng trước đó. Một sự đặc trưng sơ đồ tương tự đã được biết đến đối với các hàm có thể tính toán trong thời gian đa thức (cũng như các ngôn ngữ) [7, 8, 9, 24, 35. Theo hướng đi này, chúng tôi mong muốn trình bày một định nghĩa sơ đồ mới, đơn giản được cấu thành từ một tập hợp nhỏ các hàm lượng tử ban đầu và một tập hợp nhỏ các quy tắc xây dựng, nhằm mục đích làm cho định nghĩa sơ đồ này nắm bắt một cách thích hợp các hàm lượng tử có thể tính toán trong thời gian đa thức, nơi mà một hàm lượng tử là một hàm ánh xạ $\mathcal{H}_{\infty}$ đến $\mathcal{H}_{\infty}$. Như đã được đề cập ngắn gọn trong Phần 1, điều quan trọng cần lưu ý là thuật ngữ "hàm lượng tử" của chúng tôi hoàn toàn khác với thuật ngữ được sử dụng trong, ví dụ, [38], trong đó "hàm lượng tử" đề cập đến các hàm nhận các chuỗi đầu vào cổ điển và tạo ra hoặc là các chuỗi đầu ra cổ điển hoặc là các xác suất chấp nhận của các QTM đa băng có thời gian đa thức được hình thành tốt và do đó nó ánh xạ $\Sigma^{*}$ đến hoặc là $\Sigma^{*}$ hoặc là khoảng thực $[0,1]$, nơi mà $\Sigma$ là một bảng chữ cái thích hợp.

% ---- Nội dung từ page_08_14_vi.tex ----
\subsection*{3.1 Định nghĩa các hàm $\square_{1}^{\mathrm{QP}}$}
Định nghĩa theo lược đồ của chúng tôi dẫn đến một lớp hàm đặc biệt, gọi là $\square_{1}^{\mathrm{QP}}$ (trong đó $\square$ được đọc là "square"), nắm bắt các hàm lượng tử tính được trong thời gian đa thức ánh xạ từ $\mathcal{H}_{\infty}$ đến $\mathcal{H}_{\infty}$; lớp này gồm một tập nhỏ các hàm lượng tử khởi đầu và bốn quy tắc xây dựng đặc biệt: hợp thành, hoán đổi, phân nhánh và đệ quy lượng tử đa qubit. Định nghĩa 3.1 trình bày hình thức chính thức cho định nghĩa theo lược đồ của chúng tôi.

Sau đây, ta nói rằng một hàm lượng tử $f$ từ $\mathcal{H}_{\infty}$ đến $\mathcal{H}_{\infty}$ là bảo toàn số chiều nếu, với mọi trạng thái lượng tử $|\phi\rangle \in \mathcal{H}_{\infty}$ và mọi số $n \in \mathbb{N}^{+},|\phi\rangle \in \mathcal{H}_{2^{n}}$ kéo theo $f(|\phi\rangle) \in \mathcal{H}_{2^{n}}$ (tức là $\ell(|\phi\rangle)=\ell(f(|\phi\rangle))$ ).

Định nghĩa 3.1 Gọi $\square_{1}^{\text {QP }}$ là tập hợp tất cả các hàm lượng tử thu được từ các hàm lượng tử khởi đầu trong Lược đồ I bằng một số hữu hạn lần (kể cả 0 lần) áp dụng các quy tắc xây dựng II-IV lên các hàm lượng tử đã được xây dựng trước đó, trong đó các Lược đồ I-IV được cho như sau. Gọi $|\phi\rangle$ là một trạng thái lượng tử bất kỳ trong $\mathcal{H}_{\infty}$.\\
I. Các hàm lượng tử khởi đầu. Cho $\theta \in[0,2 \pi) \cap \tilde{\mathbb{C}}$ và $a \in\{0,1\}$.

\begin{enumerate}
  \item $I(|\phi\rangle)=|\phi\rangle$. (đồng nhất)
  \item $P H A S E_{\theta}(|\phi\rangle)=|0\rangle\langle 0 \mid \phi\rangle+e^{i \theta}|1\rangle\langle 1 \mid \phi\rangle$. (dịch pha)
  \item $R O T_{\theta}(|\phi\rangle)=\cos \theta|\phi\rangle+\sin \theta(|1\rangle\langle 0 \mid \phi\rangle-|0\rangle\langle 1 \mid \phi\rangle)$. (quay quanh trục $x y$ với góc $\theta$ )
  \item $N O T(|\phi\rangle)=|0\rangle\langle 1 \mid \phi\rangle+|1\rangle\langle 0 \mid \phi\rangle$. (phủ định)
  \item $S W A P(|\phi\rangle)=\left\{\begin{array}{ll}|\phi\rangle & \text { if } \ell(|\phi\rangle) \leq 1, \\ \sum_{a, b \in\{0,1\}}|a b\rangle\langle b a \mid \phi\rangle & \text { otherwise. }\end{array}\right.$ (hoán đổi 2 qubit)
  \item $M E A S[a](|\phi\rangle)=|a\rangle\langle a \mid \phi\rangle$. (đo chiếu từng phần)\\
II. Quy tắc hợp thành. Từ $g$ và $h$, ta định nghĩa Compo $[g, h]$ như sau:
\end{enumerate}

Compo $[g, h](|\phi\rangle)=g \circ h(|\phi\rangle)(=g(h(|\phi\rangle)))$.\\
III. Quy tắc phân nhánh. Từ $g$ và $h$, ta định nghĩa Branch $[g, h]$ như sau:\\
(i) Branch $[g, h](|\phi\rangle)=|\phi\rangle \quad$ nếu $\ell(|\phi\rangle) \leq 1$,\\
(ii) Branch $[g, h](|\phi\rangle)=|0\rangle \otimes g(\langle 0 \mid \phi\rangle)+|1\rangle \otimes h(\langle 1 \mid \phi\rangle) \quad$ ngược lại.\\
IV. Quy tắc đệ quy lượng tử đa qubit. Từ $g, h$, $p$ bảo toàn số chiều, và $k, t \in \mathbb{N}^{+}$, ta định nghĩa $k Q \operatorname{Rec}_{t}\left[g, h, p \mid \mathcal{F}_{k}\right]$ như sau:\\
(i) $k Q \operatorname{Rec}_{t}\left[g, h, p \mid \mathcal{F}_{k}\right](|\phi\rangle)=g(|\phi\rangle) \quad$ nếu $\ell(|\phi\rangle) \leq t$,\\
(ii) $k Q R e c_{t}\left[g, h, p \mid \mathcal{F}_{k}\right](|\phi\rangle)=h\left(\sum_{s:|s|=k}|s\rangle \otimes f_{s}\left(\left\langle s \mid \psi_{p, \phi}\right\rangle\right)\right) \quad$ ngược lại,\\
trong đó $\left|\psi_{p, \phi}\right\rangle=p(|\phi\rangle)$ và $\mathcal{F}_{k}=\left\{f_{s}\right\}_{s \in\{0,1\}^{k}} \subseteq\left\{k Q \operatorname{Rec}_{t}\left[g, h, p \mid \mathcal{F}_{k}\right], I\right\}$. Để nhấn mạnh " $k$," chúng tôi gọi quy tắc này là đệ quy lượng tử $k$-qubit. Trong trường hợp $k=1$, để ngắn gọn ta viết $Q \operatorname{Rec}_{t}\left[g, h, p \mid f_{0}, f_{1}\right]$ thay cho $1 Q R e c_{t}\left[g, h, p \mid\left\{f_{0}, f_{1}\right\}\right]$.

Trong Lược đồ I, $P H A S E_{\theta}$ và $R O T_{\theta}$ lần lượt tương ứng với các ma trận $Z_{2, \theta}$ và $R_{\theta}$ đã cho ở Mục 2.3. Sau này, chúng tôi sẽ lập luận rằng góc $\theta$ trong $P H A S E_{\theta}$ và $R O T_{\theta}$ có thể cố định thành $\pi / 4$. Hàm lượng tử MEAS gắn với phép đo chiếu từng phần, trong cơ sở tính toán $\{0,1\}$, áp dụng lên qubit đầu tiên của $|\phi\rangle$ khi $\ell(|\phi\rangle) \geq 1$, và hiển nhiên suy ra rằng $\ell($ MEAS $[i](|\phi\rangle)) \leq \ell(|\phi\rangle)$.

Trước khi đi xa hơn, để giúp người đọc hiểu hành vi của các hàm lượng tử khởi đầu liệt kê trong Lược đồ I, chúng tôi minh họa ngắn gọn cách các hàm này biến đổi các qustring cơ bản độ dài 3. Với các bit $a, b, c, d \in\{0,1\}$ với $d \neq a$, ta có $I(|a b c\rangle)=|a b c\rangle, P H A S E_{\theta}\left(|a b c\rangle=e^{i \theta a}|a b c\rangle, R O T_{\theta}(|a b c\rangle)=\right. \cos \theta|a b c\rangle+(-1)^{a} \sin \theta|\bar{a} b c\rangle, N O T(|a b c\rangle)=|\bar{a} b c\rangle, S W A P(|a b c\rangle)=|b a c\rangle, M E A S[a](|a b c\rangle)=|a b c\rangle$, và $M E A S[d](|a b c\rangle)=\mathbf{0}$, trong đó $\bar{a}=1-a$.

Lược đồ IV là lõi của định nghĩa $\square_{1}^{\mathrm{QP}}$. Quy tắc đệ quy chuẩn dùng để định nghĩa một hàm đệ quy nguyên thủy $f$ từ hai hàm $g$ và $h$ có dạng: $f(0, x)=g(x)$ và $f(n+1, x)=h(n, x, f(n, x))$ với mọi $n \in \mathbb{N}$. Quy tắc này cần một bộ đếm nội bộ (ở vị trí đối số thứ nhất của $f$) để điều khiển số lần lặp áp dụng $h$. Tuy nhiên, trong Lược đồ IV, chúng tôi không dùng bộ đếm như vậy. Thay vào đó, chúng tôi dùng chiến lược chia để trị để chia một trạng thái lượng tử đã cho theo từng qubit. Ở mỗi bước quy nạp, chúng tôi xử lý các trạng thái lượng tử ngắn hơn $k$-qubit cho đến khi trạng thái lượng tử có độ dài không quá $t$. Chúng tôi muốn đưa ra hai ví dụ cụ thể về cách Lược đồ IV hoạt động trong các trường hợp $k=1,2$. Lưu ý rằng trong trường hợp $k=1$, Lược đồ IV trở thành quy tắc đệ quy lượng tử 1-qubit (hay đơn qubit) được mô tả như sau:\\
(i') $Q R e c_{t}\left[g, h, p \mid f_{0}, f_{1}\right](|\phi\rangle)=g(|\phi\rangle) \quad$ nếu $\ell(|\phi\rangle) \leq t$,\\
(ii') $Q R e c_{t}\left[g, h, p \mid f_{0}, f_{1}\right](|\phi\rangle)=h\left(|0\rangle \otimes f_{0}\left(\left\langle 0 \mid \psi_{p, \phi}\right\rangle\right)+|1\rangle \otimes f_{1}\left(\left\langle 1 \mid \psi_{p, \phi}\right\rangle\right)\right) \quad$ ngược lại, trong đó mỗi hàm $f_{0}$ và $f_{1}$ phải là $I$ hoặc $Q \operatorname{Rec}_{t}\left[g, h, p \mid f_{0}, f_{1}\right]$.

\footnotetext{${ }^{\text {I }}$ Hình thức hiện tại của các Lược đồ I-IV hiệu chỉnh những sai khác do hình thức ban đầu trong bản tóm tắt mở rộng 39 gây ra.
}Ví dụ 1. Trước hết xét trường hợp $k=1$. Trong ví dụ này, ta đặt $t=1, g=N O T, h=S W A P$, và $p=N O T$, đồng thời viết ngắn gọn $F$ cho $Q R e c_{1}\left[N O T, S W A P, N O T \mid f_{0}, f_{1}\right]$ với $f_{0}=I$ và $f_{1}=F$. Cần lưu ý rằng $\mathbf{0}$ là một đối tượng đặc biệt và theo Bổ đề 3.4 (1), ta có $g(\mathbf{0})=\mathbf{0}$ với mọi hàm lượng tử $g$. Gọi $|\phi\rangle$ là một trạng thái lượng tử bất kỳ trong $\mathcal{H}_{\infty}$ được đưa vào $F$.\\
(1) Giả sử $|\phi\rangle$ có độ dài 1 và có dạng tổng quát $\alpha|0\rangle+\beta|1\rangle$. Theo Bổ đề 3.4(2), chỉ cần xét cơ sở tính toán $B_{1}=\{|0\rangle,|1\rangle\}$. Vì $\ell(|\phi\rangle) \leq t$, kết quả của $f$ trên các qubit cơ sở $|0\rangle$ và $|1\rangle$ được tính như sau.\\
(i) $F(|0\rangle)=g(|0\rangle)=N O T(|0\rangle)=|1\rangle$.\\
(ii) $F(|1\rangle)=g(|1\rangle)=N O T(|1\rangle)=|0\rangle$.

Hiển nhiên $F(\mathbf{0})=\mathbf{0}$. Vì $|\phi\rangle$ là chồng chập dạng $\alpha|0\rangle+\beta|1\rangle$, các tính toán trên lập tức suy ra

$$
F(|\phi\rangle)=F(\alpha|0\rangle+\beta|1\rangle)=N O T(\alpha|0\rangle+\beta|1\rangle)=\alpha N O T(|0\rangle)+\beta N O T(|1\rangle)=\alpha|1\rangle+\beta|0\rangle .
$$

(2) Tiếp theo, giả sử $|\phi\rangle$ có độ dài 2. Trong trường hợp $|\phi\rangle$ là trạng thái lượng tử cơ sở $|00\rangle$ trong $B_{2}=\{|00\rangle,|01\rangle,|10\rangle,|11\rangle\}$, ta có $\left|\psi_{p, \phi}\right\rangle=\left|\psi_{N O T, 00}\right\rangle=N O T(|00\rangle)=|10\rangle$. Tương tự, nếu $|\phi\rangle$ lần lượt là $|01\rangle,|10\rangle$, và $|11\rangle$, thì trạng thái lượng tử $\left|\psi_{p, \phi}\right\rangle$ tương ứng là $|11\rangle,|00\rangle$, và $|01\rangle$. Liên quan đến phép đo chiếu từng phần, suy ra rằng với mọi bit $a \in\{0,1\},\langle 1 \mid 1 a\rangle=|a\rangle,\langle 0 \mid 0 a\rangle=|a\rangle,\langle 1 \mid 0 a\rangle=\mathbf{0}$, và $\langle 0 \mid 1 a\rangle=\mathbf{0}$. Cũng lưu ý rằng $I(\mathbf{0})=\mathbf{0}$. Dùng các đẳng thức này cùng với $F(|0\rangle)=|1\rangle$ và $F(|1\rangle)=|0\rangle$ thu được ở (1), ta có thể tính kết quả $F(|\phi\rangle)$ như sau.\\
(i) $F(|00\rangle)=h\left(|0\rangle \otimes f_{0}\left(\left\langle 0 \mid \psi_{N O T, 00}\right\rangle\right)+|1\rangle \otimes f_{1}\left(\left\langle 1 \mid \psi_{N O T, 00}\right\rangle\right)=h(|0\rangle \otimes I(\langle 0 \mid 10\rangle)+|1\rangle \otimes F(\langle 1 \mid 10\rangle))=\right. h(|0\rangle \otimes I(\mathbf{0})+|1\rangle \otimes F(|0\rangle))=h(|1\rangle \otimes F(|0\rangle))=S W A P(|1\rangle \otimes|1\rangle)=S W A P(|11\rangle)=|11\rangle$.\\
(ii) $F(|10\rangle)=h\left(|0\rangle \otimes f_{0}\left(\left\langle 0 \mid \psi_{N O T, 10}\right\rangle\right)+|1\rangle \otimes f_{1}\left(\left\langle 1 \mid \psi_{N O T, 10}\right\rangle\right)=h(|0\rangle \otimes I(\langle 0 \mid 00\rangle)+|1\rangle \otimes F(\langle 1 \mid 00\rangle))=\right. h(|0\rangle \otimes I(|0\rangle)+|1\rangle \otimes F(\mathbf{0}))=h(|0\rangle \otimes I(|0\rangle))=S W A P(|0\rangle \otimes|0\rangle)=S W A P(|00\rangle)=|00\rangle$.\\
(iii) $F(|01\rangle)=h\left(|0\rangle \otimes f_{0}\left(\left\langle 0 \mid \psi_{N O T, 01}\right\rangle\right)+|1\rangle \otimes f_{1}\left(\left\langle 1 \mid \psi_{N O T, 01}\right\rangle\right)=h(|0\rangle \otimes I(\langle 0 \mid 11\rangle)+|1\rangle \otimes F(\langle 1 \mid 11\rangle))=\right. h(|0\rangle \otimes I(\mathbf{0})+|1\rangle \otimes F(|1\rangle))=h(|1\rangle \otimes F(|1\rangle))=S W A P(|1\rangle \otimes|0\rangle)=S W A P(|10\rangle)=|01\rangle$.\\
(iv) $F(|11\rangle)=h\left(|0\rangle \otimes f_{0}\left(\left\langle 0 \mid \psi_{N O T, 11}\right\rangle\right)+|1\rangle \otimes f_{1}\left(\left\langle 1 \mid \psi_{N O T, 11}\right\rangle\right)=h(|0\rangle \otimes I(\langle 0 \mid 01\rangle)+|1\rangle \otimes F(\langle 1 \mid 01\rangle))=\right. h(|0\rangle \otimes I(|1\rangle)+|1\rangle \otimes F(\mathbf{0}))=h(|0\rangle \otimes I(|1\rangle))=S W A P(|0\rangle \otimes|1\rangle)=S A W P(|01\rangle)=|10\rangle$.

Theo Bổ đề 3.4 (2), chẳng hạn khi $|\phi\rangle$ có dạng $\alpha|00\rangle+\beta|10\rangle$, ta được

$$
F(|\phi\rangle)=F(\alpha|00\rangle+\beta|10\rangle)=\alpha F(|00\rangle)+\beta F(|10\rangle)=\alpha|11\rangle+\beta|00\rangle .
$$

(3) Xét trường hợp độ dài của $|\phi\rangle$ là 3. Với cơ sở tính toán $B_{3}= \{|000\rangle,|001\rangle,|010\rangle, \ldots,|111\rangle\}$, ta chỉ tính các kết quả $F(|\phi\rangle)$ cho đầu vào $|\phi\rangle$ thuộc $\{|001\rangle,|101\rangle\}$. Lưu ý rằng $\left|\psi_{N O T, 001}\right\rangle=|101\rangle$ và $\left|\psi_{N O T, 101}\right\rangle=|001\rangle$. Nhắc lại từ (1)-(2) rằng $F(|01\rangle)=|01\rangle$ và $F(\mathbf{0})=\mathbf{0}$.\\
(i) $F(|001\rangle)=h\left(|0\rangle \otimes f_{0}\left(\left\langle 0 \mid \psi_{N O T, 001}\right\rangle\right)+|1\rangle \otimes f_{1}\left(\left\langle 1 \mid \psi_{N O T, 001}\right\rangle\right)=h(|0\rangle \otimes I(\langle 0 \mid 101\rangle)+|1\rangle \otimes F(\langle 1 \mid 101\rangle))=\right. h(|0\rangle \otimes I(\mathbf{0})+|1\rangle \otimes F(|01\rangle))=h(|1\rangle \otimes F(|01\rangle))=S W A P(|1\rangle \otimes|01\rangle)=S A W P(|101\rangle)=|011\rangle$.\\
(ii) $F(|101\rangle)=h\left(|0\rangle \otimes f_{0}\left(\left\langle 0 \mid \psi_{N O T, 101}\right\rangle\right)+|1\rangle \otimes f_{1}\left(\left\langle 1 \mid \psi_{N O T, 101}\right\rangle\right)=h(|0\rangle \otimes I(\langle 0 \mid 001\rangle)+|1\rangle \otimes F(\langle 1 \mid 001\rangle))=\right. h(|0\rangle \otimes I(|01\rangle)+|1\rangle \otimes F(\mathbf{0}))=h(|0\rangle \otimes I(|01\rangle))=S W A P(|0\rangle \otimes|01\rangle)=S A W P(|001\rangle)=|001\rangle$.

Chẳng hạn nếu $|\phi\rangle$ có dạng $\alpha|001\rangle+\beta|101\rangle$, Bổ đề 3.4(2) kéo theo $F(|\phi\rangle)=F(\alpha|001\rangle+\beta|101\rangle)= \alpha F(|001\rangle)+\beta F(|101\rangle)=\alpha|011\rangle+\beta|001\rangle$.\\
Ví dụ 2. Ta xét một ví dụ khác cho trường hợp $k=2$. Chẳng hạn đặt $t=3, g=N O T, h=S W A P$, và $p=N O T$, ta viết $F$ cho $2 Q R e c_{3}\left[N O T, S W A P, N O T \mid f_{00}, f_{01}, f_{10}, f_{11}\right]$ với $f_{00}=I$, $f_{01}=N O T, f_{10}=F$, và $f_{11}=F$. Gọi $|\phi\rangle$ là một trạng thái lượng tử bất kỳ trong $\mathcal{H}_{\infty}$.\\
(1) Khi độ dài của $|\phi\rangle$ không quá 3, ta lập tức có $F(|\phi\rangle)=g(|\phi\rangle)=N O T(|\phi\rangle)$. Ví dụ, $F(|0\rangle)=|1\rangle, F(|10\rangle)=|00\rangle$, và $F(|101\rangle)=|001\rangle$.\\
(2) Giả sử độ dài của $|\phi\rangle$ là 4, xét các trạng thái lượng tử cơ sở trong $B_{4}= \{|0000\rangle,|0001\rangle,|0010\rangle, \ldots,|1111\rangle\}$. Ở đây, ta chỉ tập trung vào hai trường hợp: $|\phi\rangle \in\{|0101\rangle,|1011\rangle\}$. Ta lưu ý rằng $\left|\psi_{N O T, 0010}\right\rangle=N O T(|0010\rangle)=|1010\rangle$ và $\left|\psi_{N O T, 1011}\right\rangle=N O T(|1011\rangle)=|0011\rangle$.\\
(i) $F(|0010\rangle)=h\left(|00\rangle \otimes f_{00}\left(\left\langle 00 \mid \psi_{N O T, 0010}\right\rangle\right)+|01\rangle \otimes f_{01}\left(\left\langle 01 \mid \psi_{N O T, 0010}\right\rangle\right)+|10\rangle \otimes f_{10}\left(\left\langle 10 \mid \psi_{N O T, 0010}\right\rangle\right)+\right. \left.|11\rangle \otimes f_{11}\left(\left\langle 11 \mid \psi_{N O T, 0010}\right\rangle\right)\right)=h(|00\rangle \otimes I(\langle 00 \mid 1010\rangle)+|01\rangle \otimes N O T(\langle 01 \mid 1010\rangle)+|10\rangle \otimes F(\langle 10 \mid 1010\rangle)+|11\rangle \otimes F(\langle 11 \mid 1010\rangle))=h(|00\rangle \otimes I(\mathbf{0})+|01\rangle \otimes N O T(\mathbf{0})+|10\rangle \otimes F(|10\rangle)+|11\rangle \otimes F(\mathbf{0}))=h(|10\rangle \otimes F(|10\rangle))= S W A P(|10\rangle \otimes|00\rangle)=S W A P(|1000\rangle)=|0100\rangle$.\\
(ii) $F(|1101\rangle)=h\left(|00\rangle \otimes f_{00}\left(\left\langle 00 \mid \psi_{N O T, 1101}\right\rangle\right)+|01\rangle \otimes f_{01}\left(\left\langle 01 \mid \psi_{N O T, 1101}\right\rangle\right)+|10\rangle \otimes f_{10}\left(\left\langle 10 \mid \psi_{N O T, 1101}\right\rangle\right)+\right. \left.|11\rangle \otimes f_{11}\left(\left\langle 11 \mid \psi_{N O T, 1101}\right\rangle\right)\right)=h(|00\rangle \otimes I(\langle 00 \mid 0101\rangle)+|01\rangle \otimes N O T(\langle 01 \mid 0101\rangle)+|10\rangle \otimes F(\langle 10 \mid 0101\rangle)+|11\rangle \otimes F(\langle 11 \mid 0101\rangle))=h(|00\rangle \otimes I(\mathbf{0})+|01\rangle \otimes N O T(|01\rangle)+|10\rangle \otimes F(\mathbf{0})+|11\rangle \otimes F(\mathbf{0}))=h(|01\rangle \otimes N O T(|01\rangle))= S W A P(|01\rangle \otimes|11\rangle)=S W A P(|0111\rangle)=|1011\rangle$.

\subsection*{3.2 Một lớp con của $\square_{1}^{\mathrm{QP}}$ và các tính chất cơ bản}
Trong sáu hàm lượng tử khởi đầu của Lược đồ I, MEAS có hành vi khá khác biệt. Nó có thể thay đổi chuẩn cũng như số chiều của các trạng thái lượng tử, và điều đó khiến nó về bản chất là không khả nghịch. Vì lý do này, trong thực hành thường có lợi khi giới hạn sự chú ý vào một lớp con của $\square_{1}^{\mathrm{QP}}$, gọi là $\widehat{\square_{1}^{\mathrm{QP}}}$, lớp này cấm hoàn toàn việc sử dụng MEAS.

Định nghĩa 3.2 Ký hiệu $\widehat{\square_{1}^{\mathrm{QP}}}$ chỉ lớp con của $\square_{1}^{\mathrm{QP}}$ được định nghĩa bởi các Lược đồ I-IV ngoại trừ MEAS trong Lược đồ I.

Không dùng MEAS trong Lược đồ I, mọi hàm lượng tử $f$ trong $\widehat{\square_{1}^{\mathrm{QP}}}$ đều bảo toàn số chiều của đầu vào; nói cách khác, $f$ thỏa mãn $\ell(f(|\phi\rangle))=\ell(|\phi\rangle)$ với mọi đầu vào $|\phi\rangle \in \mathcal{H}_{\infty}$.

Tiếp theo, ta minh họa cách xây dựng các cổng đơn vị điển hình bằng định nghĩa theo lược đồ của chúng tôi.\\
Bổ đề 3.3 Các hàm sau thuộc $\widehat{\square_{1}^{\mathrm{QP}}}$. Gọi $|\phi\rangle$ là một phần tử bất kỳ trong $\mathcal{H}_{\infty}$.

\begin{enumerate}
  \item $\operatorname{CNOT}(|\phi\rangle)=\left\{\begin{array}{ll}|\phi\rangle & \text { if } \ell(|\phi\rangle) \leq 1, \\ |0\rangle\langle 0 \mid \phi\rangle+|1\rangle \otimes N O T(\langle 1 \mid \phi\rangle) & \text { otherwise. }\end{array} \quad\right.$ (controlled-NOT)
  \item $Z_{1, \theta}(|\theta\rangle)=e^{i \theta}|0\rangle\langle 0 \mid \phi\rangle+|1\rangle\langle 1 \mid \phi\rangle$.
  \item $z R O T_{\phi}(|\phi\rangle)=e^{i \theta}|0\rangle\langle 0 \mid \phi\rangle+e^{-i \theta}|1\rangle\langle 1 \mid \phi\rangle$. (quay quanh trục $z$)
  \item $G P S_{\theta}(|\phi\rangle)=e^{i \theta}|\phi\rangle$. (dịch pha toàn cục)
  \item $W H(|\phi\rangle)=\frac{1}{\sqrt{2}}|0\rangle \otimes(\langle 0 \mid \phi\rangle+\langle 1 \mid \phi\rangle)+\frac{1}{\sqrt{2}}|1\rangle \otimes(\langle 0 \mid \phi\rangle-\langle 1 \mid \phi\rangle)$. (biến đổi Walsh-Hadamard)
  \item $\operatorname{CPHASE}_{\theta}(|\phi\rangle)= \begin{cases}|\phi\rangle \\ \frac{1}{\sqrt{2}} \sum_{b \in\{0,1\}}\left(|0\rangle\langle b \mid \phi\rangle+e^{i \theta b}|1\rangle\langle b \mid \phi\rangle\right) & \text { if } \ell(|\phi\rangle) \leq 1, \quad \text { (controlled-PHASE) }\end{cases}$
\end{enumerate}

Chứng minh. Chỉ cần xây dựng từng hàm lượng tử trong bổ đề từ các hàm lượng tử khởi đầu và bằng cách áp dụng các quy tắc xây dựng. Các hàm đích này được xây dựng như sau. (1) CNOT $=$ Branch $[I, N O T]$. Lưu ý rằng khi $\ell(|\phi\rangle) \leq 1$, Lược đồ III $(\mathrm{i})$ suy ra $\operatorname{CNOT}(|\phi\rangle)=|\phi\rangle$, khớp với Mục 1 của $C N O T$. (2) $Z_{1, \theta}=N O T \circ P H A S E_{\theta} \circ N O T$. (3) $z R O T_{\theta}=Z_{1, \theta} \circ P H A S E_{-\theta}$. (4) $G P S_{\theta}= Z_{1, \theta} \circ P H A S E_{\theta}$. (5) $W H=\operatorname{ROT}_{\frac{\pi}{4}} \circ N O T$. (6) $C P H A S E_{\theta}=\operatorname{Branch}[W H, f]$, trong đó $f=\operatorname{Branch}\left[I, G P S_{\theta}\right] \circ W H \circ N O T$.

Vì mọi hàm $\square_{1}^{\mathrm{QP}}$-function $f$ được xây dựng bằng cách áp dụng các Lược đồ I-IV, một lần áp dụng một trong các lược đồ được xem là một bước cơ bản của quá trình xây dựng để sinh ra $f$. Điều này giúp ta định nghĩa độ phức tạp mô tả của $f$ là số lần tối thiểu dùng các lược đồ đó để xây dựng $f$. Ví dụ, mọi hàm khởi đầu đều có độ phức tạp mô tả bằng 1 vì chỉ dùng Lược đồ I một lần. Như đã minh họa trong chứng minh Bổ đề 3.3, các hàm lượng tử $C N O T, Z_{1, \theta}$, và $W H$ có độ phức tạp mô tả nhiều nhất là 3, còn $z R O T_{\theta}$ và $G P S_{\theta}$ nhiều nhất là 4, trong khi $C P H A S E_{\theta}$ có độ phức tạp mô tả nhiều nhất là 15. Thước đo độ phức tạp này là thiết yếu khi chứng minh, chẳng hạn, Bổ đề 3.4 vì bổ đề sẽ được chứng minh bằng quy nạp theo độ phức tạp mô tả của hàm lượng tử đích. Ở Mục 6.2, chúng tôi sẽ thảo luận ngắn về thước đo độ phức tạp này cho nghiên cứu tương lai.

Các tính chất nền tảng của các hàm $\widehat{\square_{1}^{\mathrm{QP}}}$ được nêu trong bổ đề sau. Một hàm lượng tử từ $\mathcal{H}_{\infty}$ đến $\mathcal{H}_{\infty}$ được gọi là bảo toàn chuẩn nếu $\| f(|\phi\rangle)\|=\||\phi\rangle \|$ đúng với mọi trạng thái lượng tử $|\phi\rangle$ trong $\mathcal{H}_{\infty}$.

Bổ đề 3.4 Cho $f$ là một hàm lượng tử bất kỳ trong $\widehat{\square_{1}^{\mathrm{QP}}}$ và cho $|\phi\rangle,|\psi\rangle \in \mathcal{H}_{\infty}$ cùng $\alpha \in \mathbb{C}$.

\begin{enumerate}
  \item $f(\mathbf{0})=\mathbf{0}$, trong đó $\mathbf{0}$ là vectơ không.
  \item $f(|\phi\rangle+|\psi\rangle)=f(|\phi\rangle)+f(|\psi\rangle)$.
  \item $f(\alpha|\phi\rangle)=\alpha \cdot f(|\phi\rangle)$.
  \item $f$ bảo toàn số chiều và bảo toàn chuẩn.
\end{enumerate}

Chứng minh. Gọi $f$ là một hàm $\widehat{\square_{1}^{\mathrm{QP}}}$ bất kỳ, gọi $|\phi\rangle,|\psi\rangle \in \mathcal{H}_{\infty}$, và gọi $\alpha \in \mathbb{C}$ cùng $\theta \in[0,2 \pi)$ là các hằng số. Như đã nói trước đó, ta chứng minh bổ đề bằng quy nạp theo độ phức tạp mô tả của $f$. Nếu $f$ là một trong các hàm lượng tử khởi đầu của Lược đồ I, thì dễ kiểm tra rằng nó thỏa các Điều kiện $1-4$ của bổ đề. Đặc biệt, khi $|\phi\rangle$ là vectơ không $\mathbf{0}$, tất cả các biểu thức $e^{i \theta}|\phi\rangle, \cos \theta|\phi\rangle, \sin \theta|\phi\rangle,\langle 0 \mid \phi\rangle,\langle 1 \mid \phi\rangle$, và $\langle b a \mid \phi\rangle$ dùng trong Lược đồ I đều là $\mathbf{0}$; do đó, $|b\rangle \otimes\langle 1 \mid \phi\rangle$ và $|b\rangle \otimes\langle 0 \mid \phi\rangle$ cũng là $\mathbf{0}$ với mỗi bit $b \in\{0,1\}$. Vì vậy Điều kiện 1 được chứng minh.

Trong các Lược đồ II-IV, ta xét Lược đồ IV vì các lược đồ còn lại dễ chứng minh thỏa Điều kiện 1-4. Gọi $g, h$, và $p$ là các hàm lượng tử trong $\widehat{\square_{1}^{\mathrm{QP}}}$ và giả sử $p$ bảo toàn số chiều. Theo giả thuyết quy nạp, giả sử $g, h$, và $p$ thỏa Điều kiện $1-4$. Để dễ đọc, ta viết $f$ cho $k Q R e c_{t}\left[g, h, p \mid \mathcal{F}_{k}\right]$. Trong phần sau, ta chỉ tập trung vào Điều kiện 2 và 4 vì các điều kiện còn lại dễ kiểm tra. Lập luận của ta sẽ dùng quy nạp theo độ dài đầu vào $|\phi\rangle$ đưa vào $f$.\\
(i) Mục tiêu là chứng minh $f$ thỏa Điều kiện 2. Trước hết, xét trường hợp $\ell(|\phi\rangle) \leq t$. Khi đó suy ra $f(|\phi\rangle+|\xi\rangle)=g(|\phi\rangle+|\xi\rangle)=g(|\phi\rangle)+g(|\xi\rangle)$ vì giả sử $g$ thỏa Điều kiện 2. Tiếp theo xét trường hợp $\ell(|\phi\rangle)>t$. Từ định nghĩa của $f$, ta có $f(|\phi\rangle+|\xi\rangle)=h\left(\sum_{s:|s|=k}|s\rangle \otimes f_{s}\left(\left\langle s \mid \psi_{p, \phi, \xi}\right\rangle\right)\right)$, trong đó $\left|\psi_{p, \phi, \xi}\right\rangle=p(|\phi\rangle+|\xi\rangle)$. Do $p(|\phi\rangle+|\xi\rangle)=p(|\phi\rangle)+p(|\xi\rangle)$ theo giả thuyết quy nạp, ta kết luận $\left\langle s \mid \psi_{p, \phi, \xi}\right\rangle=\left\langle s \mid \psi_{p, \phi}\right\rangle+\left\langle s \mid \psi_{p, \xi}\right\rangle$ với mỗi xâu $s \in\{0,1\}^{k}$. Theo giả thuyết quy nạp, suy ra $f_{s}\left(\left\langle s \mid \psi_{p, \phi, \xi}\right\rangle\right)=f_{s}\left(\left\langle s \mid \psi_{p, \phi}\right\rangle\right)+f_{s}\left(\left\langle s \mid \psi_{p, \xi}\right\rangle\right)$, dẫn đến đẳng thức $|s\rangle \otimes f_{s}\left(\left\langle s \mid \psi_{p, \phi, \xi}\right\rangle\right)=|s\rangle \otimes f_{s}\left(\left\langle s \mid \psi_{p, \phi}\right\rangle\right)+ |s\rangle \otimes f_{s}\left(\left\langle s \mid \psi_{p, \xi}\right\rangle\right)$. Dùng Điều kiện 2 cho $h$, ta thu được $h\left(\sum_{s:|s|=k}|s\rangle \otimes f_{s}\left(\left\langle s \mid \psi_{p, \phi, \xi}\right\rangle\right)\right)=\sum_{s:|s|=k} h(|s\rangle \otimes \left.f_{s}\left(\left\langle s \mid \psi_{p, \phi}\right\rangle\right)\right)+\sum_{s:|s|=k} h\left(|s\rangle \otimes f_{s}\left(\left\langle s \mid \psi_{p, \xi}\right\rangle\right)\right)$. Từ đó kết luận $f(|\phi\rangle+|\xi\rangle)=f(|\phi\rangle)+f(|\xi\rangle)$.\\
(ii) Ta muốn chứng minh $f$ thỏa Điều kiện 4. Theo giả thuyết quy nạp, với mọi $s \in\{0,1\}^{k}$, ta có $\left\|f_{s}\left(\left\langle s \mid \psi_{p, \phi}\right\rangle\right)\right\|=\left\|\left\langle s \mid \psi_{p, \phi}\right\rangle\right\|$ và $\| h(|\phi\rangle)\|=\||\phi\rangle \|$. Các đẳng thức này suy ra $\| f(|\phi\rangle) \|^{2}= \| h\left(\sum_{s:|s|=k}|s\rangle \otimes f_{s}\left(\left\langle s \mid \psi_{p, \phi}\right\rangle\right)\right)\left\|^{2}=\right\| \sum_{s:|s|=k}|s\rangle \otimes f_{s}\left(\left\langle s \mid \psi_{p, \phi}\right\rangle\right)\left\|^{2}=\sum_{s:|s|=k}\right\| f_{s}\left(\left\langle s \mid \psi_{p, \phi}\right\rangle\right) \|^{2}$, bằng $\sum_{s:|s|=k}\left\|\left\langle s \mid \psi_{p, \phi}\right\rangle\right\|^{2}$. Hạng cuối cùng trùng với $\|\left|\psi_{p, \phi}\right\rangle \|^{2}$, mà bằng $\| p(|\phi\rangle) \|^{2}$. Điều này suy ra Điều kiện 4 vì $\| p(|\phi\rangle)\|=\||\phi\rangle \|$.

Bổ đề 3.4(4) chỉ ra rằng mọi hàm trong $\widehat{\square_{1}^{\mathrm{QP}}}$ cũng đóng vai trò là các hàm ánh xạ từ $\Phi_{\infty}$ đến $\Phi_{\infty}$ thay vì từ $\mathcal{H}_{\infty}$ đến $\mathcal{H}_{\infty}$.

Cho một hàm lượng tử $g$ bảo toàn số chiều và bảo toàn chuẩn, nghịch đảo của $g$ là một hàm lượng tử duy nhất $f$ thỏa điều kiện rằng với mọi $|\phi\rangle \in \mathcal{H}_{\infty}, f \circ g(|\phi\rangle)=g \circ f(|\phi\rangle)=|\phi\rangle$. Hàm lượng tử đặc biệt này được ký hiệu là $g^{-1}$.

Mệnh đề tiếp theo đảm bảo rằng $\widehat{\square_{1}^{\mathrm{QP}}}$ đóng dưới phép "nghịch đảo" vì $\widehat{\square_{1}^{\mathrm{QP}}}$ không chứa MEAS.\\
Mệnh đề 3.5 Với mọi hàm lượng tử $g \in \widehat{\square_{1}^{\mathrm{QP}}}, g^{-1}$ tồn tại và thuộc $\widehat{\square_{1}^{\mathrm{QP}}}$.\\
Chứng minh. Ta chứng minh mệnh đề này bằng quy nạp theo độ phức tạp mô tả của $g$ đã nêu ở trên. Nếu $g$ là một trong các hàm lượng tử khởi đầu, ta định nghĩa nghịch đảo của nó $g^{-1}$ như sau: $I^{-1}=I$, $P H A S E_{\theta}^{-1}=P H A S E_{-\theta}, R O T_{\theta}^{-1}=R O T_{-\theta}, N O T^{-1}=N O T$, và $S W A P^{-1}=S W A P$. Nếu $g$ được tạo từ một hoặc nhiều hàm lượng tử khác bằng một trong các quy tắc xây dựng, thì nghịch đảo của nó được định nghĩa là: $\operatorname{Compo}[g, h]^{-1}=\operatorname{Compo}\left[h^{-1}, g^{-1}\right], \operatorname{Branch}[g, h]^{-1}=\operatorname{Branch}\left[g^{-1}, h^{-1}\right]$, và $k Q R e c_{t}\left[g, h, p \mid\left\{f_{s}\right\}_{s \in\{0,1\}^{k}}\right]^{-1}=k Q R e c_{t}\left[g^{-1}, p^{-1}, h^{-1} \mid\left\{f_{s}^{-1}\right\}_{s \in\{0,1\}^{k}}\right]$.

Việc chúng tôi chọn các Lược đồ I-IV được thúc đẩy bởi một tập cổng lượng tử phổ dụng cụ thể. Lưu ý rằng một lựa chọn khác về các hàm lượng tử khởi đầu và quy tắc xây dựng có thể dẫn đến một tập hàm $\square_{1}^{\mathrm{QP}}$ khác. Lược đồ I dùng góc "tùy ý" $\theta$ trong $[0,2 \pi) \cap \widetilde{\mathbb{C}}$ để đưa vào $P H A S E_{\theta}$ và $R O T_{\theta}$; tuy nhiên, ta có thể giới hạn $\theta$ về giá trị duy nhất $\frac{\pi}{4}$ vì $P H A S E_{\theta}$ và $R O T_{\theta}$ với giá trị tùy ý $\theta \in[0,2 \pi)$ có thể được xấp xỉ với độ chính xác mong muốn bằng $W H$ và $Z_{2, \frac{\pi}{4}}$. Phần cuối này xuất phát từ thực tế là mọi toán tử đơn vị đơn qubit đều có thể được xấp xỉ đến độ chính xác bất kỳ từ các cổng lượng tử $W H$ và $Z_{2, \frac{\pi}{4}}$ 6 (xem thêm 25). Để thảo luận thêm về lựa chọn các lược đồ, xem Mục 6.1.

\subsection*{3.3 Xây dựng các hàm lượng tử phức tạp hơn}
Trước khi trình bày định lý chính (Định lý 4.1), chúng tôi muốn chuẩn bị các hàm lượng tử hữu ích và các quy tắc xây dựng mới suy ra trực tiếp từ các Lược đồ I-IV. Các hàm lượng tử và quy tắc xây dựng này sẽ được dùng cho chứng minh bổ đề then chốt của chúng tôi (Bổ đề 4.3), bổ đề hỗ trợ định lý chính.

Với mỗi $k \in \mathbb{N}^{+}$, ta giả sử thứ tự từ điển chuẩn $<$ trên $\{0,1\}^{k}$ và mọi phần tử trong $\{0,1\}^{k}$ được liệt kê theo từ điển là $s_{1}<s_{2}<\cdots<s_{2^{k}}$. Ví dụ, khi $k=2$, ta có $00<01<10<11$. Với mỗi xâu $s \in\{0,1\}^{n}$, $s^{R}$ ký hiệu xâu đảo của $s$; tức là $s^{R}=s_{n} s_{n-1} \cdots s_{2} s_{1}$ nếu $s=s_{1} s_{2} \cdots s_{n-1} s_{n}$. Ta mở rộng khái niệm này cho trạng thái lượng tử như sau. Cho trạng thái lượng tử $|\phi\rangle \in \mathcal{H}_{2^{n}}$, xâu đảo của $|\phi\rangle$, ký hiệu $\left|\phi^{R}\right\rangle$, có dạng $\sum_{s:|s|=n}\langle s \mid \phi\rangle \otimes\left|s^{R}\right\rangle$, trong đó $\langle s \mid \phi\rangle$ chỉ là một vô hướng. Chẳng hạn, nếu $|\phi\rangle=\alpha|01\rangle+\beta|10\rangle$, thì $\left|\phi^{R}\right\rangle=\alpha|10\rangle+\beta|01\rangle$.

Bổ đề 3.6 Cho $k \in \mathbb{N}$ với $k \geq 2$, cho $g, h \in \widehat{\square_{1}^{\mathrm{QP}}}$, và cho $\mathcal{G}_{k}=\left\{g_{s} \mid s \in\{0,1\}^{k}\right\}$ là một tập các hàm $\widehat{\square_{1}^{\mathrm{QP}}}$. Các hàm lượng tử sau đều thuộc $\widehat{\square_{1}^{\mathrm{QP}}}$. Bổ đề cũng đúng nếu thay $\widehat{\square_{1}^{\mathrm{QP}}}$ bằng $\square_{1}^{\mathrm{QP}}$.

Gọi $|\phi\rangle$ là một trạng thái lượng tử bất kỳ trong $\mathcal{H}_{\infty}$.

\begin{enumerate}
  \item Compo $\left[\mathcal{G}_{k}\right](|\phi\rangle)=g_{s_{1}} \circ g_{s_{2}} \circ \cdots \circ g_{s_{2^{k}}}(|\phi\rangle)$. (hợp thành nhiều hàm)
  \item Switch $_{k}[g, h](|\phi\rangle)=g(|\phi\rangle)$ nếu $\ell(|\phi\rangle)<k$ và Switch ${ }_{k}[g, h](|\phi\rangle)=h(|\phi\rangle)$ ngược lại. (chuyển mạch)
  \item $L E N G T H_{k}[g](|\phi\rangle)=|\phi\rangle$ nếu $\ell(|\phi\rangle)<k$ và $L E N G T H_{k}[g](|\phi\rangle)=g(|\phi\rangle)$ ngược lại.
  \item $\operatorname{REMOVE} E_{k}(|\phi\rangle)=|\phi\rangle$ nếu $\ell(|\phi\rangle)<k$ và $\operatorname{REMOVE} E_{k}(|\phi\rangle)=\sum_{s:|s|=k}\langle s \mid \phi\rangle \otimes|s\rangle$ ngược lại.
  \item $R E P_{k}(|\phi\rangle)=|\phi\rangle$ nếu $\ell(|\phi\rangle)<k$ và $R E P_{k}(|\phi\rangle)=\sum_{s:|s|=n-k}\langle s \mid \phi\rangle \otimes|s\rangle$ ngược lại.
  \item $S W A P_{k}(|\phi\rangle)=|\phi\rangle$ nếu $\ell(|\phi\rangle)<2 k$ và $S W A P_{k}(|\phi\rangle)=\sum_{s:|s|=k} \sum_{t:|t|=k}|s t\rangle\langle t s \mid \phi\rangle$ ngược lại.
  \item $\operatorname{REVERSE}(|\phi\rangle)=\left|\phi^{R}\right\rangle$.
  \item $\operatorname{Branch}_{k}\left[\mathcal{G}_{k}\right](|\phi\rangle)=|\phi\rangle$ nếu $\ell(|\phi\rangle)<k$ và $\operatorname{Branch}_{k}\left[\mathcal{G}_{k}\right](|\phi\rangle)=\sum_{s:|s|=k}|s\rangle \otimes g_{s}(\langle s \mid \phi\rangle)$ ngược lại.
  \item $\operatorname{RevBranch}{ }_{k}\left[\mathcal{G}_{k}\right](|\phi\rangle) \quad=\quad|\phi\rangle \quad$ if $\quad \ell(|\phi\rangle) \quad<\quad k \quad$ and $\operatorname{RevBranch}{ }_{k}\left[\mathcal{G}_{k}\right](|\phi\rangle) \quad= \sum_{s:|s|=k} g_{s}\left(\sum_{u:|u|=n-k}\langle u s \mid \phi\rangle \otimes|u\rangle\right) \otimes|s\rangle$ otherwise, trong đó $n=\ell(|\phi\rangle)$.\\
Sự khác nhau giữa $R E M O V E_{k}$ và $R E P_{k}$ khá tinh tế nhưng $R E M O V E_{k}$ dời $k$ qubit đầu của đầu vào xuống cuối, còn $R E P_{k}$ dời $k$ qubit cuối lên đầu. Chẳng hạn với các qustring cơ bản độ dài 4, ta có $R E P_{1}\left(\left|a_{1} a_{2} a_{3} a_{4}\right\rangle\right)=\left|a_{4} a_{1} a_{2} a_{3}\right\rangle$ và $R E M O V E_{1}\left(\left|a_{1} a_{2} a_{3} a_{4}\right\rangle\right)=\left|a_{2} a_{3} a_{4} a_{1}\right\rangle$. Tương tự, $\operatorname{Branch}_{k}\left[\mathcal{G}_{k}\right]$ áp dụng mỗi $g_{s}$ trong $\mathcal{G}_{k}$ lên trạng thái lượng tử thu được từ $|\phi\rangle$ bằng cách loại bỏ $k$ qubit đầu, trong khi $\operatorname{RevBranch}{ }_{k}\left[\mathcal{G}_{k}\right]$ áp dụng $g_{s}$ lên trạng thái lượng tử thu được bằng cách loại bỏ $k$ qubit cuối, bất cứ khi nào $k \leq \ell(|\phi\rangle)$. Ví dụ, nếu $\mathcal{G}_{k}=\{h\}_{s \in\{0,1\}^{k}}$ cho một hàm lượng tử đơn $h$, thì với mọi xâu $s \in\{0,1\}^{k}$, ta có $\operatorname{Branch}_{k}\left[\mathcal{G}_{k}\right](|s\rangle|\phi\rangle)=|s\rangle \otimes h(|\phi\rangle)$ và $\operatorname{RevBranch}_{k}\left[\mathcal{G}_{k}\right](|\phi\rangle|s\rangle)=h(|\phi\rangle) \otimes|s\rangle$.
\end{enumerate}

Chứng minh Bổ đề 3.6, Cho $k \in \mathbb{N}^{+}, g \in \widehat{\square_{1}^{\mathrm{QP}}}$, và $\mathcal{G}_{k}=\left\{g_{s} \mid s \in\{0,1\}^{k}\right\} \subseteq \widehat{\square_{1}^{\mathrm{QP}}}$. Với mỗi chỉ số $i \in[k]$, gọi $s_{i}$ là phần tử thứ $i$ theo thứ tự từ điển của $\{0,1\}^{k}$.

\begin{enumerate}
  \item Trước hết đặt $f_{2^{k}}=g_{s_{2^{k}}}$ và định nghĩa quy nạp $f_{i-1}=\operatorname{Compo}\left[g_{s_{i-1}}, f_{i}\right]$ với mọi chỉ số $i \in\left[2,2^{k}\right]_{\mathbb{Z}}$ để thu được $f_{1}=\operatorname{Compo}\left[\mathcal{G}_{k}\right]$. Hàm lượng tử nhận được $f_{1}$ hiển nhiên thuộc $\widehat{\square_{1}^{\mathrm{QP}}}$ vì $k$ là hằng số cố định độc lập với đầu vào của $f_{1}$.
  \item Đây là một trường hợp đặc biệt của quy tắc đệ quy lượng tử đơn qubit (hay quy tắc đệ quy lượng tử 1-qubit), trong đó $t=k-1, p=I$, và $f_{0}=f_{1}=I$. Vì vậy, Switch $[g, h]$ thuộc $\widehat{\square_{1}^{\mathrm{QP}}}$.
  \item Vì $L E N G T H_{k}[g]=$ Switch $_{k}[I, g], L E N G T H_{k}$ thuộc $\widehat{\square_{1}^{\mathrm{QP}}}$.
  \item Ta bắt đầu với trường hợp $k=1$. Hàm lượng tử cần tìm $R E M O V E_{1}$ được định nghĩa là
\end{enumerate}

$$
\operatorname{REMOVE}_{1}(|\phi\rangle)= \begin{cases}|\phi\rangle & \text { if } \ell(|\phi\rangle) \leq 1, \\ \sum_{a \in\{0,1\}}|a\rangle \otimes \operatorname{REMOV} E_{1}\left(\left\langle a \mid \psi_{S W A P, \phi}\right\rangle\right) & \text { otherwise },\end{cases}
$$

Chính thức hơn, ta đặt $R E M O V E_{1}=Q R e c_{1}\left[I, I, S W A P \mid R E M O V E_{1}, R E M O V E_{1}\right]$. Với trường hợp $k=2$, ta định nghĩa $R E M O V E_{2}=R E M O V E_{1} \circ R E M O V E_{1}$. Với mỗi chỉ số $k \geq 3, R E M O V E_{k}$ được thu được như sau. Gọi $h_{k}^{\prime}$ là $k$ phép hợp thành của $R E M O V E_{1}$ và định nghĩa $R E M O V E_{k}$ là $L E N G T H_{k}\left[h_{k}^{\prime}\right]$. Ta lưu ý rằng $L E N G T H_{k}$ là cần thiết trong định nghĩa vì $h_{k}^{\prime}$ không được đảm bảo bằng $R E M O V E_{k}$ khi $\ell(|\phi\rangle) \leq k-1$.\\
5) Trước tiên, ta định nghĩa $R E P_{1}$ là $R E P_{1}=Q R e c_{1}\left[I, S W A P, I \mid R E P_{1}, R E P_{1}\right]$, tức là,

$$
R E P_{1}(|\phi\rangle)= \begin{cases}|\phi\rangle & \text { if } \ell(|\phi\rangle) \leq 1 \\ \sum_{a \in\{0,1\}} S W A P\left(|a\rangle \otimes R E P_{1}(\langle a \mid \phi\rangle)\right) & \text { otherwise. }\end{cases}
$$

Hiển nhiên, $R E P_{1}$ thuộc $\widehat{\square_{1}^{\mathrm{QP}}}$. Với chỉ số tổng quát $k>1$, ta định nghĩa $R E P_{k}$ theo cách sau. Trước hết đặt $h_{k}^{\prime}$ là $k$ phép hợp thành của $R E P_{1}$. Cuối cùng đặt $R E P_{k}=L E N G T H_{k}\left[h_{k}^{\prime}\right]$.\\
6) Trước hết ta hiện thực một hàm lượng tử $S W A P_{i, i+j}$, hoán đổi qubit thứ $i$ và qubit thứ $(i+j)$ của mọi đầu vào. Mục tiêu này đạt được bằng cách xây dựng quy nạp $S W A P_{i, i+j}$ như sau. Ban đầu đặt $S W A P_{i, i+1}=R E P_{i-1} \circ S W A P \circ R E M O V E_{i-1}$. Với mọi chỉ số $j \in[k-i]$, định nghĩa $S W A P_{i, i+j}= S W A P_{i+j-1, i+j} \circ S W A P_{i, i+j-1} \circ S W A P_{i+j-1, i+j}$. Sau đó đặt $g=S W A P_{k, 2 k} \circ S W A P_{k-1,2 k-1} \circ \cdots \circ S W A P_{2, k+2} \circ S W A P_{1, k+1}$. Cuối cùng, chỉ cần định nghĩa $S W A P_{k}$ là $L E N G T H_{2 k}[g]$.\\
7) Hàm lượng tử cần tìm REVERSE có thể được định nghĩa là REVERSE = QRec $_{1}\left[I\right.$, REMOV $E_{1}, I \mid$ REVERSE, REVERSE $]$, tức là,

$$
R E V E R S E(|\phi\rangle)= \begin{cases}|\phi\rangle & \text { if } \ell(|\phi\rangle) \leq 1 \\ R E M O V E_{1}\left(\sum_{a \in\{0,1\}}(|a\rangle \otimes R E V E R S E(\langle a \mid \phi\rangle))\right) & \text { otherwise. }\end{cases}
$$

\begin{enumerate}
  \setcounter{enumi}{7}
  \item Lưu ý rằng khi $k=2, \operatorname{Branch}_{k}\left[\mathcal{G}_{k}\right]$ trùng với Branch và do đó thuộc $\widehat{\square_{1}^{\mathrm{QP}}}$. Từ đây giả sử $k \geq 3$. Với mỗi xâu $s \in\{0,1\}^{k}$, đặt $g_{s}^{(0)}=g_{s}$. Với mỗi chỉ số $i \in \mathbb{N}$ thỏa $i<k$ và mỗi xâu $s \in\{0,1\}^{*}$ với $|s|=k-i-1$, ta định nghĩa quy nạp $g_{s}^{(i+1)}$ là $\operatorname{Branch}\left[g_{s 0}^{(i)}, g_{s 1}^{(i)}\right]$, tức là,
\end{enumerate}

$$
g_{s}^{(i+1)}(|\phi\rangle)= \begin{cases}|\phi\rangle & \text { if } \ell(|\phi\rangle \leq 1, \\ |0\rangle \otimes g_{s 0}^{(i)}(\langle 0 \mid \phi\rangle)+|1\rangle \otimes g_{s 1}^{(i)}(\langle 1 \mid \phi\rangle) & \text { otherwise } .\end{cases}
$$

Cuối cùng, ta đặt $\operatorname{Branch}_{k}\left[\mathcal{G}_{k}\right]=g_{\lambda}^{(k)}$, trong đó $\lambda$ là xâu rỗng.\\
9) Giả sử $k \geq 2$, đặt $\operatorname{RevBranch}_{k}\left[\left\{g_{s}\right\}_{s \in\{0,1\}^{k}}\right]=\operatorname{REMOVE} E_{k} \circ \operatorname{Branch}_{k}\left[\left\{g_{s}\right\}_{s \in\{0,1\}^{k}}\right] \circ R E P_{k}$.

Bổ đề kế tiếp cho thấy ta có thể mở rộng mọi song ánh cổ điển trên $\{0,1\}^{k}$ thành hàm $\widehat{\square_{1}^{\mathrm{QP}}}$ liên kết của nó, hàm này hành xử đúng hệt song ánh đó trên $k$ bit đầu của đầu vào.

Bổ đề 3.7 Cho $k$ là một hằng trong $\mathbb{N}^{+}$. Với mọi song ánh $f$ từ $\{0,1\}^{k}$ đến $\{0,1\}^{k}$, tồn tại một hàm $\square_{1}^{\mathrm{QP}}$ $g_{f}$ sao cho, với mọi trạng thái lượng tử $|\phi\rangle \in \mathcal{H}_{\infty}$,

$$
g_{f}(|\phi\rangle)= \begin{cases}|\phi\rangle & \text { if } \ell(|\phi\rangle) \leq k-1, \\ \sum_{s \in\{0,1\}^{k}}|f(s)\rangle\langle s \mid \phi\rangle & \text { otherwise } .\end{cases}
$$

Chứng minh. Cho một song ánh $f$ trên $\{0,1\}^{k}$, chỉ cần chỉ ra tồn tại hàm $\widehat{\square_{1}^{\mathrm{QP}}}$ $h$ thỏa $h(|s\rangle|\phi\rangle)=|f(s)\rangle|\phi\rangle$ với mọi xâu $s \in\{0,1\}^{k}$ và mọi trạng thái lượng tử $|\phi\rangle \in \mathcal{H}_{\infty}$, vì $g_{f}$ được thu từ $h$ chỉ bằng cách đặt $g_{f}=L E N G T H_{k}[h]$. Lưu ý rằng nếu $h \in \widehat{\square_{1}^{\mathrm{QP}}}, h(\mathbf{0})=\mathbf{0}$ khiến $g_{f}(\mathbf{0})$ bằng $\mathbf{0}$.

Một song ánh trên $\{0,1\}^{k}$ về bản chất là một hoán vị trên $\left\{s_{1}, s_{2}, \ldots, s_{2^{k}}\right\}$, trong đó mỗi $s_{i}$ là xâu thứ $i$ theo thứ tự từ điển trong $\{0,1\}^{k}$, và do đó có thể biểu diễn thành tích của một số hữu hạn phép hoán vị đổi chỗ, mỗi phép đổi chỗ hoán đổi hai số phân biệt. Việc này có thể thực hiện bằng cách áp dụng hợp thành nhiều lần của các $S W A P_{i, i+j}$ tương ứng, được định nghĩa trong chứng minh Bổ đề 3.6(6). Do đó, $h$ thuộc $\widehat{\square_{1}^{\mathrm{QP}}}$.

Cho một trạng thái lượng tử $|\phi\rangle$, ta có thể đồng thời áp dụng một hàm lượng tử $f$ lên $k$ qubit đầu của $|\phi\rangle$ và một hàm lượng tử khác $g$ lên phần còn lại.

Bổ đề 3.8 Với $f, g \in \widehat{\square_{1}^{\mathrm{QP}}}$ và $k \in \mathbb{N}^{+}$, hàm lượng tử $f^{\leq k} \otimes g$, được định nghĩa bởi

$$
\left(f^{\leq k} \otimes g\right)(|\phi\rangle)= \begin{cases}f(|\phi\rangle) & \text { if } \ell(|\phi\rangle) \leq k \\ \sum_{s \in\{0,1\}^{k}} f(|s\rangle) \otimes g(\langle s \mid \phi\rangle) & \text { otherwise }\end{cases}
$$

thuộc $\widehat{\square_{1}^{\mathrm{QP}}}$. Ta ký hiệu $g$ này là Skip $[f]$.\\
Chứng minh. Cho hàm lượng tử $f$, ta giới hạn việc áp dụng $f$ lên $k$ qubit cuối của mọi đầu vào $|\phi\rangle$ bằng cách đặt $f^{\prime}=Q R e c_{k}\left[f, I, I \mid f_{0}, f_{1}\right]$, trong đó $f_{0}$ và $f_{1}$ đều là $f^{\prime}$. Suy ra $f^{\prime}(|\phi\rangle|s\rangle)=|\phi\rangle \otimes f(|s\rangle)$ với mọi $|\phi\rangle \in \mathcal{H}_{\infty}$ và $s \in\{0,1\}^{k}$.

Đặt $h=R E P_{k} \circ f^{\prime} \circ R E M O V E_{k}$. Hàm lượng tử này $h$ thỏa $h(|s\rangle|\phi\rangle)=f(|s\rangle) \otimes|\phi\rangle$ với mọi $s \in\{0,1\}^{k}$. Sau đó ta định nghĩa $\mathcal{G}_{k}=\left\{g_{s}\right\}_{s \in\{0,1\}^{k}}$ với $g_{s}=g$ cho mọi xâu $s \in\{0,1\}^{k}$ và đặt $g^{\prime}=h \circ \operatorname{Branch}_{k}\left[\mathcal{G}_{k}\right]$. Khi đó suy ra hàm lượng tử cần tìm $f^{\leq k} \otimes g$ bằng $\operatorname{Switch}_{k+1}\left[f, g^{\prime}\right]$.

Tiếp theo, ta trình bày Bổ đề 3.9, hữu ích cho chứng minh bổ đề then chốt của chúng tôi (Bổ đề 4.3) ở Mục 5. Bổ đề này cho phép ta bỏ qua, trước khi áp dụng một hàm lượng tử cho trước, một số lượng tùy ý các số 0 cho đến khi đọc được một số lượng cố định các số 1.\\
Bổ đề 3.9 Cho $f$ là một hàm lượng tử trong $\widehat{\square_{1}^{\mathrm{QP}}}$ và cho $k$ là một hằng trong $\mathbb{N}^{+}$. Tồn tại một hàm lượng tử $g$ trong $\widehat{\square_{1}^{\mathrm{QP}}}$ sao cho $g\left(\left|0^{m} 1^{k}\right\rangle \otimes|\phi\rangle\right)=\left|0^{m} 1^{k}\right\rangle \otimes f(|\phi\rangle)$ và $g\left(\left|0^{m+1}\right\rangle\right)=\left|0^{m+1}\right\rangle$ với mọi số $m \in \mathbb{N}$ và mọi trạng thái lượng tử $|\phi\rangle \in \mathcal{H}_{\infty}$. Bổ đề cũng đúng khi thay $\widehat{\square_{1}^{\mathrm{QP}}}$ bằng $\square_{1}^{\mathrm{QP}}$.

Chứng minh. Cho $k \geq 2$. Với hàm lượng tử $f \in \widehat{\square_{1}^{\mathrm{QP}}}$, trước hết ta mở rộng $f$ thành $f^{\prime}$ sao cho $f^{\prime}\left(\left|1^{k-1}\right\rangle|\phi\rangle\right)= \left|1^{k-1}\right\rangle \otimes f(|\phi\rangle)$ với mọi $|\phi\rangle \in \mathcal{H}_{\infty}$ và $f^{\prime}\left(\left|0^{m+1}\right\rangle\right)=\left|0^{m+1}\right\rangle$ với mọi $m \in \mathbb{N}$. Hàm lượng tử $f^{\prime}$ này có thể thu được theo quy nạp như sau. Ta đặt $f_{k-1}=\operatorname{Branch}[I, f], f_{i}=\operatorname{Branch}\left[I, f_{i+1}\right]$ với mỗi $i \in[k-2]$, và\\
cuối cùng định nghĩa $f^{\prime}$ là $f_{1}$. Khi $k=1$, ta đơn giản đặt $f^{\prime}=f$. Hàm lượng tử cần tìm $g$ trong bổ đề phải thỏa

$$
g(|\phi\rangle)= \begin{cases}|\phi\rangle & \text { if } \ell(|\phi\rangle) \leq 1, \\ |0\rangle \otimes f^{\prime}(g(\langle 0 \mid \phi\rangle))+|1\rangle\langle 1 \mid \phi\rangle & \text { otherwise } .\end{cases}
$$

Hàm $g$ này được định nghĩa hình thức là $g=Q R e c_{1}\left[I, \operatorname{Branch}\left[f^{\prime}, I\right], I \mid g, I\right]$. Điều này hoàn tất chứng minh.\\
Trong khuôn khổ của chúng tôi, có thể xây dựng một dạng "hạn chế" của biến đổi Fourier lượng tử (QFT). Cho xâu nhị phân $s=s_{1} s_{2} \cdots s_{k}$ độ dài $k$ với $s_{i} \in\{0,1\}$, ta ký hiệu $n u m(s)$ là số nguyên có dạng $\sum_{i=1}^{k} s_{i} 2^{k-i}$. Chẳng hạn, $\operatorname{num}(011)=1 \cdot 2^{1}+1 \cdot 2^{0}=5$ và $\operatorname{num}(1010)=1 \cdot 2^{3}+1 \cdot 2^{1}=10$. Hơn nữa, gọi $\omega_{k}=e^{2 \pi i / 2^{k}}$, trong đó $i=\sqrt{-1}$.

Bổ đề 3.10 Cho $k$ là một hằng cố định bất kỳ trong $\mathbb{N}^{+}$. Biến đổi Fourier lượng tử $k$-qubit sau thuộc $\widehat{\square_{1}^{\mathrm{QP}}}$. Với mọi phần tử $|\phi\rangle$ trong $\mathcal{H}_{\infty}$, gọi

$$
F_{k}(|\phi\rangle)= \begin{cases}|\phi\rangle & \text { if } \ell(|\phi\rangle)<k \\ \frac{1}{2^{k / 2}} \sum_{t:|t|=k} \sum_{s:|s|=k} \omega_{k}^{\text {num }(s) \text { num }(t)}|s\rangle\langle t \mid \phi\rangle & \text { otherwise }\end{cases}
$$

Chứng minh. Khi $k=1, F_{1}$ trùng với $W H$ và do đó $F_{1}$ thuộc $\widehat{\square_{1}^{\mathrm{QP}}}$ theo Bổ đề 3.3. Tiếp theo, giả sử $k \geq 2$. Đã biết rằng, với mọi $x_{1}, x_{2}, \ldots, x_{k} \in\{0,1\}$,

\begin{equation*}
F_{k}\left(\left|x_{1} x_{2} \cdots x_{k}\right\rangle\right)=\frac{1}{2^{k / 2}}\left(|0\rangle+\omega_{1}^{x_{k}}|1\rangle\right)\left(|0\rangle+\omega_{1}^{x_{k-1}} \omega_{2}^{x_{k}}|1\rangle\right) \cdots\left(|0\rangle+\prod_{i=1}^{k} \omega_{i}^{x_{i}}|1\rangle\right) . \tag{1}
\end{equation*}

Về sự kiện này và chứng minh của nó, xem chẳng hạn 25.\\
Hãy nhắc lại hàm lượng tử đặc biệt $S W A P_{i, i+j}$ trong chứng minh Bổ đề 3.6(6), hàm này hoán đổi giữa qubit thứ $i$ và qubit thứ $j$. Dùng $C P H A S E_{\theta}$, với mọi cặp chỉ số $i, j$ thỏa $i<j$, ta định nghĩa $C P H A S E_{\theta}^{(i, j)}$ là $S W A P_{1, i} \circ S W A P_{2, j} \circ C P H A S E_{\theta} \circ S W A P_{2, j} \circ S W A P_{1, i}$, trong đó ta áp dụng $C P H A S E_{\theta}$ lên qubit thứ $i$ và qubit thứ $j$. Trước hết ta muốn xây dựng $G_{k}^{(1)}=F_{k} \circ R E V E R S E$, hàm này hoạt động tương tự $F_{k}$ nhưng nhận $\left|\phi^{R}\right\rangle$ làm đầu vào. Để đạt mục tiêu đó, ta định nghĩa quy nạp $\left\{G_{j}^{(i)}\right\}_{i, j \in \mathbb{N}^{+}}$ như sau. Ban đầu đặt $G_{0}^{(i)}=I$ với mọi $i \in \mathbb{N}^{+}$. Tiếp theo, định nghĩa $G_{1}^{(i)}$ là $G_{1}^{(i)}=H$ nếu $i=1$, và $G_{1}^{(i)}= R E P_{i-1} \circ H \circ R E M O V E_{i-1}$ ngược lại. Với mọi chỉ số $k \geq 2, G_{k}^{(i)}$ được định nghĩa là $G_{k-1}^{(i)} \circ C P H A S E_{\frac{\pi}{2^{k-1}}}^{(i, i+k-1)} \circ \left(G_{k-2}^{(i)}\right)^{-1} \circ G_{k-1}^{(i+1)}$. Không khó để chỉ ra rằng $G_{k}^{(1)}$ trùng với $F_{k} \circ R E V E R S E$ theo Eqn. (1).

Vì $G_{k}^{(1)}=F_{k} \circ R E V E R S E$, chỉ cần định nghĩa $F_{k}$ là $G_{k}^{(1)} \circ R E V E R S E$.\\
Định dạng QFT tổng quát, trong đó $k$ không bị giới hạn là một hằng số cụ thể, sẽ được bàn ở Mục 6.1 liên hệ với lựa chọn các Lược đồ I-IV tạo thành lớp hàm $\square_{1}^{\mathrm{QP}}$.

\section*{4 Các đóng góp chính}
Trong Mục 3.1, chúng tôi đã giới thiệu các hàm $\square_{1}^{\mathrm{QP}}$ và các hàm $\widehat{\square_{1}^{\mathrm{QP}}}$ ánh xạ từ $\mathcal{H}_{\infty}$ đến $\mathcal{H}_{\infty}$ bằng cách áp dụng các Lược đồ I-IV hữu hạn lần. Định lý chính của chúng tôi (Định lý 4.1) khẳng định rằng $\square_{1}^{\mathrm{QP}}$ có thể đặc trưng chính xác mọi hàm trong FBQP ánh xạ từ $\{0,1\}^{*}$ đến $\{0,1\}^{*}$, và do đó đặc trưng mọi ngôn ngữ trong BQP trên $\{0,1\}$ bằng cách đồng nhất ngôn ngữ với hàm đặc trưng tương ứng của nó. Định lý này sẽ được chứng minh bằng hai bổ đề then chốt, Bổ đề 4.244 .3 .

\subsection*{4.1 Một đặc trưng mới của FBQP}
Mục tiêu của chúng tôi là chứng minh sức mạnh của các hàm $\square_{1}^{\mathrm{QP}}$ (và do đó của các hàm $\widehat{\square_{1}^{\mathrm{QP}}}$) bằng cách chỉ ra trong Định lý 4.1 rằng các hàm $\square_{1}^{\mathrm{QP}}$ (cũng như các hàm $\widehat{\square_{1}^{\mathrm{QP}}}$) đặc trưng chính xác các hàm FBQP trên $\{0,1\}^{*}$. Để làm điều này, không may là có hai khó khăn lớn cần vượt qua.

Khó khăn thứ nhất nảy sinh khi xử lý các ký hiệu băng của QTM chỉ bằng qubit. Lưu ý rằng các QTM làm việc trên bảng chữ cái đầu vào không nhị phân đã được biết là có thể mô phỏng bằng các QTM nhận bảng chữ cái đầu vào nhị phân $\{0,1\}$. Ngay cả khi ta giảm thành công kích thước bảng chữ cái đầu vào xuống 2, trên thực tế, mô phỏng các QTM như vậy vẫn phải đòi hỏi xử lý đúng các ký hiệu băng không nhị phân, đặc biệt là ký hiệu phân biệt

% ---- Nội dung từ page_15_21_vi.tex ----
ký hiệu trống. Vì lý do thuận tiện sau này, chúng tôi sử dụng " $b$ " để chỉ ký hiệu trống, thay vì $\#$. Theo cách tương tự, giả sử rằng kết quả của một hàm $\square_{1}^{\mathrm{QP}}$ $f$ bao gồm một phần quan trọng, có ý nghĩa và phần "rác" còn lại, là tàn dư của quá trình tính toán. Vì chúng tôi chỉ sử dụng các qubit ( $|0\rangle$ và $|1\rangle$ ) để biểu diễn đầu vào và đầu ra, làm thế nào chúng tôi có thể phân biệt giữa phần có ý nghĩa và phần rác? Để mô phỏng các QTMs bằng các hàm $\square_{1}^{\text {QP }}$, do đó chúng tôi cần "mã hóa" tất cả các ký hiệu băng thành các qubit.

Bài viết này giới thiệu sơ đồ mã hóa đơn giản sau đây. Chúng tôi đặt $\hat{0}=00, \hat{1}=01$, và $\hat{b}=10$. Chúng tôi cũng đặt $\hat{2}=11$ và $\hat{3}=10$ để sử dụng sau này. Bảng chữ cái đầu vào $\Sigma=\{0,1\}$ do đó được dịch sang $\{\hat{0}, \hat{1}\}$ và bảng chữ cái băng $\Gamma=\Sigma \cup\{b\}$ được mã hóa thành $\{\hat{0}, \hat{1}, \hat{b}\}$. Cho mỗi chuỗi nhị phân $s=s_{1} s_{2} \cdots s_{n}$ với $s_{i} \in\{0,1\}$ cho mọi chỉ số $i \in[n]$, một mã $\tilde{s}$ của $s$ biểu thị chuỗi $\hat{s}_{1} \hat{s}_{2} \cdots \hat{s}_{n} \hat{2}$, trong đó mục cuối cùng $\hat{2}$ phục vụ như một dấu hiệu kết thúc, đánh dấu cuối mã. Khi đó ta có $|\tilde{s}|=2|s|+2$. Là những ví dụ nhanh, ta có được $\widetilde{0110}=\hat{0} \hat{1} \hat{1} \hat{0} \hat{2}=0001010011$ và $\widetilde{01} \cdot \widetilde{11}=\hat{0} \hat{1} \hat{2} \hat{1} \hat{2}=000111010111$. Sau này, chúng tôi sẽ chỉ ra trong bổ đề 5.1 rằng việc mã hóa các chuỗi và giải mã các chuỗi đã mã hóa có thể được thực hiện bởi các hàm $\widehat{\square_{1}^{\mathrm{QP}}}$ thích hợp.

Một khó khăn khác đến từ việc các hàm $\widehat{\square_{1}^{\mathrm{QP}}}$ không thể mở rộng các qubit của chúng. Lưu ý rằng một QTM được thiết kế để tự do sử dụng không gian lưu trữ bổ sung bằng cách di chuyển đầu băng của nó đơn giản đến các ô băng mới, trống rỗng vượt ra ngoài khu vực đầu vào, nơi mà một đầu vào ban đầu được viết. Để mô phỏng một QTM như vậy, chúng ta cần mô phỏng toàn bộ các hoạt động của nó được thực hiện trong không gian băng mở rộng tự do. Ngược lại, mọi hàm $\widehat{\square_{1}^{\mathrm{QP}}}$ đều bảo toàn chiều theo Bổ đề 3.4(4), do đó số lượng qubit đầu vào phải khớp với số lượng qubit đầu ra. Nếu chúng ta muốn mô phỏng không gian lưu trữ bổ sung của QTM, thì chúng ta cần cấp cho hàm $\widehat{\square_{1}^{\mathrm{QP}}}$ mục tiêu cùng lượng qubit bổ sung để sử dụng ngay từ đầu.

Chúng tôi giải quyết vấn đề thứ hai này bằng cách mở rộng mỗi đầu vào của các hàm lượng tử bằng cách thêm các số 0 bổ sung có độ dài liên quan đến thời gian chạy của QTM. Đối với bất kỳ đa thức $p$ nào và bất kỳ hàm lượng tử $g$ nào trên $\mathcal{H}_{\infty}$, chúng tôi định nghĩa $\left|\phi^{p}(x)\right\rangle=\left|0^{|x|} 1\right\rangle\left|0^{p(|x|)} 10^{11 p(|x|)+6} 1\right\rangle|x\rangle$ và $\left|\phi_{g}^{p}(x)\right\rangle=g\left(\left|\phi^{p}(x)\right\rangle\right)$ cho mọi chuỗi $x \in\{0,1\}^{*}$. Tương tự, đối với bất kỳ hàm $f$ nào trên $\{0,1\}^{*}$, chúng tôi đặt $\left.\left|\phi^{p, f}(x)\right\rangle=\left|\widetilde{\left.0^{|f(x)|}\right\rangle}\right| 0^{|f(x)|+1} 1\right\rangle\left|\phi^{p}(x)\right\rangle$ và $\left|\phi_{g}^{p, f}(x)\right\rangle=g\left(\left|\phi^{p, f}(x)\right\rangle\right)$.

Bất kỳ hàm FBQP nào cũng nhận các chuỗi đầu vào cổ điển và tạo ra các chuỗi cổ điển là kết quả của một máy Turing lượng tử chạy trong thời gian đa thức với xác suất cao. Trái lại, các hàm $\square_{1}^{\mathrm{QP}}$ của chúng tôi $f$ là để biến đổi mỗi trạng thái lượng tử trong $\mathcal{H}_{\infty}$ thành một trạng thái khác trong $\mathcal{H}_{\infty}$. Để có được các chuỗi đầu ra cổ điển, chúng ta cần quan sát kết quả của $f$.

Định lý chính (Định lý 4.1) khẳng định đại khái như sau: đối với bất kỳ hàm FBQP $f$ nào, luôn tồn tại một hàm lượng tử $g$ trong $\widehat{\square_{1}^{\mathrm{QP}}}$ sao cho, khi chúng ta quan sát $|f(x)|$ qubit đầu tiên của kết quả $g\left(\left|\phi^{p}(x)\right\rangle\right)(= \left|\phi_{g}^{p}(x)\right\rangle$ ) của $g$ trên đầu vào $\left|\phi^{p}(x)\right\rangle$, chúng ta nhận được $f(x)$ một cách chính xác với xác suất cao. Lưu ý rằng $g\left(\left|\phi^{p}(x)\right\rangle\right)$ cũng chứa các qubit bổ sung, gọi là các qubit "rác", mà không được quan sát trong quá trình tính toán $f(x)$. Những qubit rác này thực sự là tàn dư của quá trình tính toán $f$. Tuy nhiên, có thể loại bỏ những qubit đó bằng cách một phần đảo ngược toàn bộ quá trình tính toán, nếu chúng ta biết kích thước đầu ra $|f(x)|$ trước khi tính toán (bằng cách mở rộng $\left|\phi^{p}(x)\right\rangle$ thành $\left|\phi^{p, f}(x)\right\rangle$ và $\left|\phi_{g}^{p}(x)\right\rangle$ thành $\left|\phi_{g}^{p, f}(x)\right\rangle$ ).

Định lý 4.1 (Định lý chính) Cho $f$ là một hàm có giới hạn đa thức trên $\{0,1\}^{*}$. Ba phát biểu sau đây là tương đương về mặt logic.

\begin{enumerate}
  \item Hàm $f$ thuộc FBQP .
  \item Đối với bất kỳ hằng số $\varepsilon \in[0,1 / 2)$ nào, tồn tại một hàm lượng tử $g$ trong $\widehat{\square_{1}^{\mathrm{QP}}}$ và một đa thức $p$ sao cho $|f(x)| \leq p(|x|)$ và $\left\|\left\langle f(x) \mid \phi_{g}^{p}(x)\right\rangle\right\|^{2} \geq 1-\varepsilon$ cho tất cả $x \in\{0,1\}^{*}$.
  \item Đối với bất kỳ hằng số $\varepsilon \in[0,1 / 2)$ nào, tồn tại một hàm lượng tử $g$ trong $\widehat{\square_{1}^{\mathrm{QP}}}$ và một đa thức $p$ sao cho $|f(x)| \leq p(|x|)$ và $\left|\left\langle\Psi_{f(x)} \mid \phi_{g}^{p, f}(x)\right\rangle\right|^{2} \geq 1-\varepsilon$ cho tất cả $x \in\{0,1\}^{*}$, trong đó $\left|\Psi_{f(x)}\right\rangle= |f(x)\rangle\left|\left(0^{|f(x)|+1} 1\right)^{2}\right\rangle\left|\phi^{p}(x)\right\rangle$.
\end{enumerate}

Trong các phát biểu 2-3 của Định lý 4.1, $\left\langle f(x) \mid \phi_{g}^{p}(x)\right\rangle$ là một vector khác rỗng trong khi $\left\langle\Psi_{f(x)} \mid \phi_{g}^{p, f}(x)\right\rangle$ chỉ là một đại lượng vô hướng vì $\ell\left(|f(x)\rangle\left|\phi^{p}(x)\right\rangle\right)=\ell\left(\left|\phi_{g}^{p, f}(x)\right\rangle\right)$.

Ở đây, chúng tôi muốn chứng minh định lý chính, Định lý 4.1. Vì lý do chiến lược, chúng tôi chia định lý thành hai bổ đề kỹ thuật, Bổ đề 4.2 và 4.3 Để trình bày các bổ đề, chúng tôi cần giới thiệu thêm thuật ngữ cho QTMs. Việc thiết kế các QTMs nhiều băng dễ dàng hơn so với các QTM băng đơn. Tuy nhiên, vì các QTM nhiều băng có thể được mô phỏng bởi các QTM băng đơn bằng cách dịch nhiều băng thành nhiều dải của một băng đơn, nên để chứng minh Định lý 4.1, chỉ cần tập trung vào các QTM băng đơn.

Một QTM băng đơn được gọi là ở dạng chuẩn nếu $\delta\left(q_{f}, \sigma\right)=\left|q_{0}\right\rangle|\sigma\rangle|R\rangle$ đúng với bất kỳ ký hiệu băng nào\\
$\sigma \in \Gamma$. Nếu một QTM dừng lại trong một chồng chất các cấu hình cuối cùng mà đầu băng quay trở lại ô bắt đầu, thì ta gọi máy như vậy là tĩnh. Tham khảo [5] để biết các tính chất cơ bản của chúng. Vì tiện lợi, ta gọi một QTM là bảo thủ nếu nó được hình thành tốt, tĩnh và ở dạng chuẩn. Hơn nữa, một QTM được hình thành tốt được gọi là thuần túy nếu hàm chuyển đổi của nó thỏa mãn yêu cầu cụ thể sau: đối với mỗi cặp $(p, \sigma) \in Q \times \Gamma, \delta(p, \sigma)$ có dạng $\delta(p, \sigma)=e^{i \theta}|q, \tau, d\rangle$ hoặc $\delta(p, \sigma)=\cos \theta|q, \tau, d\rangle+\sin \theta\left|q^{\prime}, \tau^{\prime}, d^{\prime}\right\rangle$ cho một số $\theta \in[0,2 \pi)$ và hai bộ ( $q, \tau, d$ ) và ( $q^{\prime}, \tau^{\prime}, d^{\prime}$ ) khác biệt. Bernstein và Vazirani [5] khẳng định rằng bất kỳ QTM bảo thủ, thời gian đa thức, băng đơn nào $M$ cũng có thể được mô phỏng bởi một QTM bảo thủ, thời gian đa thức, băng đơn, thuần túy thích hợp $M^{\prime}$. Ngoài ra, khi $M$ có biên độ $\tilde{\mathbb{C}}$, thì $M^{\prime}$ cũng vậy.

Hãy nhớ rằng bất kỳ chuỗi đầu ra nào của một QTM đều bắt đầu từ ô bắt đầu và kéo dài về bên phải cho đến ký hiệu trống đầu tiên mặc dù có thể có các ký hiệu băng còn lại chưa được xóa ở các phần khác của băng. Vì tiện lợi, một QTM $M$ được gọi là có đầu ra sạch nếu, khi $M$ dừng lại, không có ký hiệu nào không trống xuất hiện ở vùng bên trái của chuỗi đầu ra, tức là vùng bao gồm tất cả các ô được đánh chỉ số bởi các số nguyên âm. Lấy một đa thức $p$ giới hạn thời gian chạy của $M$ trên mọi đầu vào. Vì $M$ dừng lại trong nhiều nhất $p(|x|)$ bước trên mọi trường hợp $x \in\{0,1\}^{*}$, nên chúng ta chỉ cần chú ý đến vùng băng thiết yếu của nó bao gồm tất cả các ô băng được đánh chỉ số từ $-p(|x|)$ đến $+p(|x|)$. Trong thực tế, chúng ta định nghĩa lại một cấu hình $\gamma$ của $M$ trên đầu vào $x$ có độ dài $n$ là một bộ ba ( $q, h, \sigma_{1} \cdots \sigma_{2 p(n)+1}$ ), trong đó $p \in Q, h \in \mathbb{Z}$ với $-p(n) \leq h \leq p(n)$, và $\sigma_{1}, \cdots, \sigma_{2 p(n)+1} \in\{0,1, b\}$ sao cho, với mọi chỉ số $i \in[2 p(n)+1], \sigma_{i}$ là một ký hiệu băng được viết tại ô được đánh chỉ số bởi $i-p(n)-1$. Lưu ý rằng ô bắt đầu nằm ở giữa $\sigma_{1} \cdots \sigma_{2 p(n)+1}$. Vì tiện lợi, ta tiếp tục sửa đổi khái niệm cấu hình bằng cách chia vùng băng thiết yếu thành hai phần ( $z_{1}, z_{2}$ ), trong đó $z_{1}=\sigma_{1} \sigma_{2} \cdots \sigma_{p(n)}$ đề cập đến vùng bên trái của ô bắt đầu, không bao gồm ô bắt đầu, và $z_{2}=\sigma_{p(n)+1} \sigma_{p(n)+2} \cdots \sigma_{2 p(n)+1}$ đề cập đến phần còn lại của vùng băng thiết yếu. Điều này cung cấp cho chúng ta một cấu hình đã được sửa đổi dưới dạng $\left(q, h, z_{1} z_{2}\right)$. Vì một lý do thực tiễn, ta tiếp tục thay đổi nó thành $\left(z_{2}, z_{1}, h, q\right)$, mà ta gọi là cấu hình lệch. Liên quan đến sự thay đổi này của các cấu hình, ta cũng sửa đổi toán tử tiến hóa thời gian ban đầu, $U_{\delta}$, của $M$ để nó hoạt động trên các cấu hình lệch. Cần lưu ý rằng toán tử mới này không thể được hiện thực hóa bởi mô hình QTM chuẩn. Để phân biệt nó với toán tử đánh giá thời gian ban đầu $U_{\delta}$ của $M$, ta gọi nó là toán tử tiến hóa thời gian lệch và viết nó là $\widehat{U}_{\delta}$.

Bổ đề 4.2 Cho $f$ là bất kỳ hàm lượng tử nào trong $\square_{1}^{\mathrm{QP}}$. Tồn tại một đa thức $p$ và một QTM băng đơn, bảo thủ, thuần túy $M$ tạo ra đầu ra sạch với biên độ $\tilde{\mathbb{C}}$ sao cho, đối với bất kỳ trạng thái lượng tử $|\phi\rangle$ nào trong $\mathcal{H}_{2^{n}}, M$ bắt đầu với một trạng thái lượng tử khác rỗng $|\phi\rangle$ được đưa ra trên băng đầu vào/công việc của nó và, khi nó dừng lại sau $p(n)$ bước, chồng chất $\left|\eta_{M, \phi}\right\rangle$ của các cấu hình lệch cuối cùng của $M$ trên $|\phi\rangle$ có dạng $f(|\phi\rangle) \otimes\left|0, q_{f}\right\rangle$, trong đó $f(|\phi\rangle)$ xuất hiện như nội dung của băng từ ô bắt đầu về bên phải cho đến ký hiệu trống đầu tiên và $q_{f}$ là trạng thái nội duy nhất của $M$.

Chứng minh. Trước tiên, chúng ta chứng minh rằng tất cả các hàm ban đầu trong Sơ đồ I có thể được tính chính xác trong thời gian đa thức trên các QTM băng đơn, có biên độ $\tilde{\mathbb{C}}$, bảo thủ, thuần túy có đầu ra sạch trên bảng chữ cái đầu vào/đầu ra $\{0,1\}$. Để rõ ràng, mục tiêu của chúng ta ở đây là chứng minh rằng, đối với mỗi hàm lượng tử $f$ trong Sơ đồ I, tồn tại một QTM với các tính chất được nêu trong bổ đề sao cho một chồng chất $\left|\eta_{M, \phi}\right\rangle$ của các cấu hình lệch cuối cùng của $M$ là dạng $f(|\phi\rangle) \otimes\left|0, q_{0}\right\rangle$, trong đó $f(|\phi\rangle)$ là nội dung của băng $M$ từ ô bắt đầu về bên phải cho đến ký hiệu trống đầu tiên.

Vì $I$ (đồng nhất) dễ mô phỏng, hãy xem xét $P H A S E_{\theta}$. Trong trường hợp này, ta lấy một QTM áp dụng một chuyển đổi có dạng $\delta\left(q_{0}, \sigma\right)=e^{i \theta \sigma}\left|q_{f}, \sigma, N\right\rangle$ cho bất kỳ bit $\sigma \in\{0,1\}$. Rõ ràng, nếu ta bắt đầu với một cấu hình lệch ban đầu $|\phi\rangle|0\rangle\left|q_{0}\right\rangle$, thì ta dừng lại với $P H A S E_{\theta}(|\phi\rangle) \otimes|0\rangle\left|q_{f}\right\rangle$; nói ngắn gọn, QTM này "tính chính xác" $P H A S E_{\theta}$. Đối với $R O T_{\theta}$, ta sử dụng chuyển đổi được định nghĩa bởi $\delta\left(q_{0}, \sigma\right)=\cos \theta\left|q_{f}, \sigma, N\right\rangle+(-1)^{\sigma} \sin \theta\left|q_{f}, \bar{\sigma}, N\right\rangle$ để tính chính xác $R O T_{\theta}$. Để mô phỏng $N O T$, ta định nghĩa một QTM có chuyển đổi $\delta\left(q_{0}, \sigma\right)=\left|q_{f}, \bar{\sigma}, N\right\rangle$. Trong trường hợp $S W A P$, ta cần chuẩn bị các trạng thái nội $p_{\sigma_{1}}$ và $r_{\sigma_{2}}$ cũng như các chuyển đổi sau: $\delta\left(q_{0}, \sigma_{1}\right)=\left|p_{\sigma_{1}}, \#, R\right\rangle, \delta\left(p_{\sigma_{1}}, \sigma_{2}\right)=\left|r_{\sigma_{2}}, \sigma_{1}, L\right\rangle, \delta\left(r_{\sigma_{2}}, \#\right)=\left|q_{f}, \sigma_{2}, N\right\rangle$ cho các bit $\sigma_{1}, \sigma_{2} \in\{0,1\}$. Đối với $M E A S[a](|\phi\rangle)$, ta bắt đầu bằng cách kiểm tra qubit đầu tiên của $|\phi\rangle$. Nếu nó không phải là $a$, thì ta cho QTM từ chối đầu vào; ngược lại, ta không làm gì cả. Một cách chính thức hơn, ta định nghĩa $\delta\left(q_{0}, \bar{a}\right)=\left|q_{\text {rej }}, \bar{a}, N\right\rangle$ và $\delta\left(q_{0}, a\right)=\left|q_{f}, a, N\right\rangle$. Không khó để thấy rằng $\left|\eta_{M, \phi}\right\rangle$ có dạng $f(|\phi\rangle) \otimes\left|0, q_{0}\right\rangle$.

Tiếp theo, chúng tôi muốn mô phỏng từng quy tắc xây dựng trên một QTM thích hợp. Theo giả thuyết quy nạp, tồn tại ba QTM băng đơn, thời gian đa thức, có biên độ $\tilde{\mathbb{C}}$, bảo thủ, thuần túy $M_{g}, M_{h}$, và $M_{p}$ thỏa mãn bổ đề cho $g, h$, và $p$, tương ứng. Bằng cách cài đặt một đồng hồ nội bộ theo cách nào đó với một đa thức $r$ nhất định, chúng tôi có thể làm cho $M_{g}, M_{h}$, và $M_{p}$ dừng lại chính xác trong $r(n)$ thời gian trên bất kỳ đầu vào nào có độ dài $n \in \mathbb{N}$. Trong phần tiếp theo, hãy để $|\phi\rangle$ biểu thị bất kỳ đầu vào nào trong $\mathcal{H}_{\infty}$.\\[0pt]
[composition] Xem xét trường hợp $f=C o m p o[g, h]$. Chúng tôi tính $f$ như sau. Trước tiên, chúng tôi chạy $M_{h}$ trên $|\phi\rangle$ và nhận được một chồng chất các cấu hình lệch cuối cùng, ví dụ, $h(|\phi\rangle) \otimes\left|0, q_{f}^{\prime}\right\rangle$ cho một trạng thái nội duy nhất $q_{f}^{\prime}$ của $M_{h}$. Vì $M_{h}$ là tĩnh và ở dạng chuẩn, chúng tôi có thể tiếp tục chạy $M_{g}$ trên trạng thái lượng tử kết quả\\
$h(|\phi\rangle)$, coi $q_{f}^{\prime}$ là trạng thái nội ban đầu mới và tạo ra $\left|\eta_{M, \phi}\right\rangle=g(h(|\phi\rangle)) \otimes\left|0, q_{f}\right\rangle$ cho một trạng thái nội duy nhất $q_{f}$ của $M_{g}$.\\[0pt]
[branching] Giả sử rằng $f=\operatorname{Branch}[g, h]$. Chúng tôi kiểm tra qubit đầu tiên của $|\phi\rangle$. Nếu nó là 0 , thì chúng tôi chạy $M_{g}$ trên $\langle 0 \mid \phi\rangle$; ngược lại, chúng tôi chạy $M_{h}$ trên $\langle 1 \mid \phi\rangle$. Điều này tạo ra $|0\rangle \otimes g(\langle 0 \mid \phi\rangle) \otimes\left|0, q_{h, f}\right\rangle+|1\rangle \otimes h(\langle 1 \mid \phi\rangle) \otimes\left|0, q_{g, f}\right\rangle$, trong đó $q_{g, f}$ và $q_{h, f}$ lần lượt là các trạng thái nội duy nhất của $M_{g}$ và $M_{h}$. Lưu ý rằng chúng tôi không cần kiểm tra xem $\ell(|\phi\rangle) \geq 2$ hay không. Một cách chính thức, chúng tôi thêm các chuyển đổi sau vào các chuyển đổi của $M_{g}$ và $M_{h}: \delta\left(q_{0}, 0\right)=\left|q_{0}, \#_{g}, R\right\rangle$ và $\delta\left(q_{0}, 1\right)=\left|q_{0}, \#_{h}, R\right\rangle$, trong đó $\#_{g}$ và $\#_{h}$ là các ký hiệu được chỉ định đánh dấu bên trái của các "ô bắt đầu" mới cho $M_{g}$ và $M_{h}$, tương ứng. Khi kết thúc, đầu băng di chuyển trở lại ô bắt đầu. Sau đó, chúng tôi thay thế $\#_{g}$ và $\#_{h}$ lần lượt bằng 0 và 1 bằng cách đi vào một trạng thái nội mới $q_{f}$; cụ thể, $\delta\left(q_{g, f}, \#_{g}\right)=\left|q_{f}, 0, N\right\rangle$ và $\delta\left(q_{h, f}, \#_{h}\right)=\left|q_{f}, 1, N\right\rangle$. Việc chuyển đổi này là có thể một cách đơn vị.\\[0pt]
[multi-qubit quantum recursion] Cuối cùng, hãy chứng minh một mô phỏng của phép đệ quy lượng tử đa qubit được giới thiệu bởi $f=k Q R e c_{t}\left[g, h, p \mid \mathcal{F}_{k}\right]$ với $\mathcal{F}_{k}=\left\{f_{s}\right\}_{s \in\{0,1\}^{k}}$. Để dễ đọc, hãy xem xét trường hợp đơn giản nhất khi $k=1, f_{0}=f$, và $f_{1}=I$. Các trường hợp khác có thể được xử lý tương tự. Xem xét QTM "nhiều băng" $M_{f}$ hoạt động đại khái như mô tả dưới đây. Hãy để $T$ là một hằng số dương đủ lớn.\\
(1) Trong giai đoạn ban đầu này, bắt đầu với đầu vào $|\phi\rangle$, chúng tôi chuẩn bị một bộ đếm trên một băng công việc mới và một đồng hồ nội bộ trên một băng công việc khác. Chúng tôi sử dụng đồng hồ để điều chỉnh thời điểm kết thúc của tất cả các đường tính toán. Đếm số lượng $\ell(|\phi\rangle)$ các qubit đơn giản bằng cách tăng bộ đếm khi di chuyển đầu băng đầu vào từ ô bắt đầu sang bên phải. Ban đầu, chúng tôi đặt trạng thái lượng tử hiện tại, ví dụ, $|\xi\rangle$ biểu thị trên băng là $|\phi\rangle$ và đặt bộ đếm hiện tại $k$ là $\ell(|\xi\rangle)(=n)$. Đi đến giai đoạn phân chia.\\
(2) Trong giai đoạn phân chia này, chúng tôi thực hiện theo cách quy nạp thủ tục sau bằng cách sử dụng đồng hồ. (\textit{) Giả sử rằng băng đầu vào hiện tại chứa một trạng thái lượng tử $|\xi\rangle$ và bộ đếm có $k=\ell(|\xi\rangle)$. Nếu $k \leq t$, thì chờ cho đến khi đồng hồ đạt $T$ và sau đó đi đến giai đoạn xử lý. Ngược lại, chạy $M_{p}$ trên $|\xi\rangle$ để tạo ra $\left|\psi_{p, \xi}\right\rangle$ và quan sát qubit đầu tiên của $\left|\psi_{p, \xi}\right\rangle$ trong cơ sở tính toán, thu được $\left\langle b \mid \psi_{p, \xi}\right\rangle$ cho mỗi $b \in\{0,1\}$. Nếu $b$ là 1 , thì chạy $M_{h}$ trên $|1\rangle\left\langle 1 \mid \psi_{p, \xi}\right\rangle$, thu được $h\left(|1\rangle\left\langle 1 \mid \psi_{p, \xi}\right\rangle\right)$, được coi là $f(|1\rangle\langle 1 \mid \xi\rangle)$, chờ cho đến khi đồng hồ đạt $T$, và sau đó bắt đầu giai đoạn xử lý. Ngược lại, khi $b$ là 0 , di chuyển bit 0 này sang một băng riêng biệt để ghi nhớ và sau đó cập nhật cả $|\xi\rangle$ và $k$ thành $\left\langle 0 \mid \psi_{p, \xi}\right\rangle$ và $k-1$, tương ứng. Tiếp tục (}).\\
(3) Trong giai đoạn xử lý này, chúng tôi bắt đầu với một trạng thái lượng tử $|\xi\rangle$, được tạo ra trong giai đoạn phân chia. Cho $k=\ell(|\xi\rangle)$. Chúng tôi thực hiện theo cách quy nạp thủ tục sau. (\textbf{) Nếu $k \leq t$, thì chúng tôi chạy $M_{g}$ trên đầu vào $|\xi\rangle$ và tạo ra $g(|\xi\rangle)$, được coi là $f(|\xi\rangle)$. Cập nhật $|\xi\rangle$ thành trạng thái lượng tử kết quả. Ngược lại, chúng tôi di chuyển lại bit 0 cuối cùng đã được lưu trữ từ băng riêng biệt, chạy $M_{h}$ trên $|0\rangle|\xi\rangle$, thu được $h(|0\rangle|\xi\rangle)$, và sau đó cập nhật $|\xi\rangle$ thành trạng thái lượng tử đã thu được và $k$ thành $k+1$ vì $\ell(|0\rangle|\xi\rangle)=k+1$. Nếu tất cả các $b$ đã được tiêu thụ (tương đương, $k=n$ ), thì chờ cho đến khi đồng hồ đạt $2 T$, xuất $|\xi\rangle$, và dừng. Ngược lại, tiếp tục (}).

Thời gian chạy của QTM trên được giới hạn trên bởi một đa thức nhất định theo độ dài $\ell(|\phi\rangle)$ vì mỗi thủ tục (*) và (**) được lặp lại nhiều nhất $\ell(|\phi\rangle)$ lần. Mặc dù $M_{f}$ lưu trữ các bit trên băng riêng biệt, nhưng tất cả các bit đó đều được di chuyển trở lại và sử dụng hết vào cuối quá trình tính toán. Thực tế này cho thấy một chồng chất các cấu hình lệch cuối cùng của $M_{f}$ có dạng $f(|\phi\rangle) \otimes\left|0, q_{f}\right\rangle$ cho một trạng thái nội duy nhất $q_{f}$. Cuối cùng, chúng tôi chuyển đổi QTM đa băng này thành một QTM băng đơn tương đương về mặt tính toán có các tính chất mong muốn được nêu trong bổ đề.

Điều này hoàn thành chứng minh của Bổ đề 4.2.\\
Vì lý do ký hiệu, chúng tôi viết $M[r]$ để chỉ chuỗi đầu ra được viết trong một cấu hình lệch cuối cùng $r$, bao gồm chỉ vùng băng thiết yếu của $M$. Chúng tôi viết $F S C_{M, n}$ để biểu thị tập hợp tất cả các cấu hình lệch cuối cùng có thể có của $M$ được tạo ra cho bất kỳ đầu vào nào có độ dài $n$.

Chìa khóa để chứng minh Định lý 4.1 là bổ đề sau đây, đảm bảo sự tồn tại của một hàm $\widehat{\square_{1}^{\mathrm{QP}}}$ $g$ mà kết quả của nó $\left.\left.g\left(\mid \phi^{p}(|\psi\rangle)\right\rangle\right)\left(=\mid \phi_{g}^{p}(|\psi\rangle)\right\rangle\right)$ "gần như" đặc trưng cho cấu hình lệch cuối cùng đã mã hóa $\sum_{x:|x|=n} \sum_{r \in F S C_{M, n}}\langle x \mid \psi\rangle|\widetilde{M[r]}\rangle\left|\xi_{x, r}\right\rangle$ của $M$, trong đó $\widetilde{M[r]}$ là mã hóa của $M[r]$. Ở đây, mã hóa $\widetilde{M[r]}$ là cần thiết để xây dựng $g$ trong chứng minh của bổ đề. Tuy nhiên, như cũng được chỉ ra trong chứng minh, việc thực hiện một thủ tục giải mã bổ sung cho $\widetilde{M[r]}$ cho phép chúng ta thay thế $\widetilde{M[r]}$ bằng $M[r]$.

Bổ đề 4.3 (Bổ đề chính) Cho $M$ là một QTM băng đơn, thời gian đa thức, có biên độ $\tilde{\mathbb{C}}$, bảo thủ, thuần túy có đầu ra sạch trên bảng chữ cái đầu vào/đầu ra $\Sigma=\{0,1\}$ và bảng chữ cái băng $\Gamma=\{0,1, b\}$. Giả sử rằng, khi $M$ dừng lại trên đầu vào $|\psi\rangle$ có độ dài $n$, một chồng chất các cấu hình lệch cuối cùng đã mã hóa có dạng\\
$\sum_{x:|x|=n} \sum_{r \in F S C_{M, n}}\langle x \mid \psi\rangle|\widetilde{M[r]}\rangle\left|\xi_{x, r}\right\rangle$, trong đó $\left|\xi_{x, r}\right\rangle$ biểu thị một trạng thái lượng tử thích hợp mô tả phần còn lại của các cấu hình lệch cuối cùng đã mã hóa ngoài $\widetilde{M[r]}$. Tồn tại một hàm lượng tử $g$ trong $\widehat{\square_{1}^{\mathrm{QP}}}$ và một đa thức $p$ sao cho, đối với bất kỳ số $n \in \mathbb{N}^{+}$và mọi trạng thái lượng tử $\left.|\psi\rangle \in \mathcal{H}_{2^{n}}, \mid \phi_{g}^{p}(|\psi\rangle)\right\rangle$ có dạng $\sum_{x:|x|=n} \sum_{r \in F S C_{M, n}}\langle x \mid \psi\rangle|\widetilde{M[r]}\rangle\left|\widehat{\xi}_{x, r}\right\rangle$ cho các trạng thái lượng tử nhất định $\left\{\left|\widehat{\xi}_{x, r}\right\rangle\right\}_{x, r}$ thỏa mãn rằng, đối với bất kỳ $x, x^{\prime} \in\{0,1\}^{n}$ và $r, r^{\prime} \in F S C_{M, n}$, (i) $\left\|\left\langle\xi_{x, r} \mid \xi_{x^{\prime}, r^{\prime}}\right\rangle\right\|=\left\|\left\langle\widehat{\xi}_{x, r} \mid \widehat{\xi}_{x^{\prime}, r^{\prime}}\right\rangle\right\|$ và (ii) $\left\langle\widehat{\xi}_{x, r} \mid \widehat{\xi}_{x, r^{\prime}}\right\rangle=0$ nếu $r \neq r^{\prime}$. Hơn nữa, có thể sửa đổi $g$ thành $g^{\prime}$, thỏa mãn $\left.\mid \phi_{g^{\prime}}^{p}(|\phi\rangle)\right\rangle=\sum_{x:|x|=n} \sum_{r \in F S C_{M, n}}\langle x \mid \psi\rangle|M[r]\rangle\left|\widehat{\xi}_{x, r}^{\prime}\right\rangle$ cho các trạng thái lượng tử thích hợp $\left\{\left|\widehat{\xi}_{x^{\prime}, r^{\prime}}^{\prime}\right\rangle\right\}_{x^{\prime}, r^{\prime}}$ thỏa mãn các điều kiện (i)-(ii).

Chứng minh của Bổ đề 4.3 khá dài và được hoãn lại cho đến Mục 5 . Trong lúc đó, chúng tôi quay lại Định lý 4.1 và trình bày chứng minh của nó bằng cách sử dụng các Bổ đề 4.2 và 4.3.

Chứng minh Định lý 4.1. Cho $\varepsilon \in[0,1 / 2)$ là bất kỳ hằng số nào và cho $f$ là bất kỳ hàm có giới hạn đa thức nào ánh xạ $\{0,1\}^{*}$ tới $\{0,1\}^{*}$.\\
$(1 \Rightarrow 2)$ Giả sử rằng $f$ thuộc FBQP. Lấy một QTM đa băng, thời gian đa thức, có biên độ $\tilde{\mathbb{C}}$, được hình thành tốt $N$ tính toán $f$ với xác suất lỗi bị chặn. Hãy chọn một đa thức $p$ giới hạn thời gian chạy của $N$ trên mọi đầu vào. Có thể "mô phỏng" $N$ với xác suất thành công cao bằng một QTM băng đơn nhất định, có biên độ $\tilde{\mathbb{C}}$, bảo thủ, thuần túy, có đầu ra sạch sao cho máy lấy đầu vào $x$ và kết thúc bằng việc tạo ra $f(x) b w$ trong vùng bên phải của ô bắt đầu của băng đầu vào/công việc với xác suất lỗi bị chặn, trong đó $b$ là một ký hiệu băng trống duy nhất. Vì bất kỳ QTM lỗi bị chặn nào cũng có thể khuếch đại tự do xác suất thành công của nó, chúng ta có thể giả sử mà không mất tính tổng quát rằng xác suất lỗi của $M$ là nhiều nhất $\varepsilon$. Lưu ý rằng cấu hình lệch cuối cùng đã mã hóa bắt đầu bằng một chuỗi đầu ra. Do đó, chồng chất các cấu hình lệch cuối cùng đã mã hóa của $M$ trên đầu vào $x$ có độ dài $n$ phải có dạng $\sum_{r \in F S C_{M, n}}|\widetilde{M[r]}\rangle\left|\xi_{x, r}\right\rangle$ cho các trạng thái lượng tử được chọn phù hợp $\left\{\left|\xi_{x, r}\right\rangle\right\}_{r \in F S C_{M, n}}$. Vì $M$ tính toán $f$ với xác suất lỗi nhiều nhất $\varepsilon$, chúng ta kết luận rằng $\sum_{r \in F S C_{M, n}} \|\langle\widetilde{f(x)} \mid \widetilde{M[r]}\rangle\left|\xi_{x, r}\right\rangle \|^{2} \geq 1-\varepsilon$ cho tất cả $x$. Bổ đề 4.3 tiếp tục cung cấp cho chúng ta một hàm lượng tử đặc biệt $g \in \widehat{\square_{1}^{\mathrm{QP}}}$ sao cho, đối với bất kỳ số $n \in \mathbb{N}$ và bất kỳ chuỗi $x \in\{0,1\}^{n},\left|\phi_{g}^{p}(x)\right\rangle$ có dạng $\sum_{r \in F S C_{M, n}}|M[r]\rangle\left|\widehat{\xi}_{x, r}\right\rangle$ cho các trạng thái lượng tử nhất định $\left\{\left|\widehat{\xi}_{x, r}\right\rangle\right\}_{r \in F S C_{M, n}}$ thỏa mãn rằng, đối với tất cả $r, r^{\prime} \in F S C_{M, n},\left\|\left\langle\xi_{x, r} \mid \xi_{x, r^{\prime}}\right\rangle\right\|=\left\|\left\langle\widehat{\xi}_{x, r} \mid \widehat{\xi}_{x, r^{\prime}}\right\rangle\right\|$ và $\left\langle\widehat{\xi}_{x, r} \mid \widehat{\xi}_{x, r^{\prime}}\right\rangle=0$ nếu $r \neq r^{\prime}$. Do đó, chúng ta kết luận rằng $\left\|\left\langle f(x) \mid \phi_{g}^{p}(x)\right\rangle\right\|^{2}=\sum_{r \in F S C_{M, n}} \|\langle f(x) \mid M[r]\rangle\left|\widehat{\xi}_{x, r}\right\rangle\left\|^{2}=\sum_{r \in F S C_{M, n}}|\langle\widetilde{f(x)} \mid \widetilde{M[r]}\rangle|^{2}\right\|\left|\widehat{\xi}_{x, r}\right\rangle \|^{2}$ vì và $\langle f(x) \mid M[r]\rangle=\langle\widetilde{f(x)} \mid \widetilde{M[r]}\rangle$ và $\left\langle\widehat{\xi}_{x, r} \mid \widehat{\xi}_{x, r^{\prime}}\right\rangle=0$ nếu $r \neq r^{\prime}$. Hơn nữa, từ $\|\left|\xi_{x, r}\right\rangle\|=\|\left|\widehat{\xi}_{x, r}\right\rangle \|$, ta có được số hạng $|\langle\widetilde{f(x)} \mid \widetilde{M[r]}\rangle|^{2} \|\left|\widehat{\xi}_{x, r}\right\rangle \|^{2}$ bằng $\|\langle\widetilde{f(x)} \mid \widetilde{M[r]}\rangle\left|\xi_{x, r}\right\rangle \|^{2}$. Do đó, $\left\|\left\langle f(x) \mid \phi_{g}^{p}(x)\right\rangle\right\|^{2}$ ít nhất là $1-\varepsilon$.\\
($2 \Rightarrow 3$) Cho $\varepsilon$ là bất kỳ hằng số nào trong $[0,1 / 2)$ và đặt $\varepsilon^{\prime}=1-\sqrt{1-\varepsilon}$. Hãy chọn một đa thức $p$ và một hàm $g \in \widehat{\square_{1}^{\mathrm{QP}}}$ sao cho $|f(x)| \leq p(|x|)$ và $\left\|\left\langle f(x) \mid \phi_{g}^{p}(x)\right\rangle\right\|^{2} \geq 1-\varepsilon^{\prime}$ cho tất cả các chuỗi $x \in\{0,1\}^{*}$. Bắt đầu với đầu vào $\left.\left|\phi^{p, f}(x)\right\rangle\left(=\left|\widetilde{\left.0^{|f(x)|}\right\rangle}\right| 0^{|f(x)|+1} 1\right\rangle\left|\phi^{p}(x)\right\rangle\right)$, trước tiên chúng ta áp dụng $g$ vào phần cuối $\left|\phi^{p}(x)\right\rangle$ của $\left|\phi^{p, f}(x)\right\rangle$ và thu được $\left|\widetilde{0^{|f(x)|}}\right\rangle\left|0^{|f(x)|+1} 1\right\rangle\left|\phi_{g}^{p}(x)\right\rangle$. Tiếp theo, chúng ta áp dụng Bổ đề 5.1 để mã hóa $|f(x)|$ qubit đầu tiên của $\left|\phi_{g}^{p}(x)\right\rangle$ với sự trợ giúp của $\left|0^{|f(x)|+1} 1\right\rangle$ và sau đó thu được $\left|\eta_{f, x}\right\rangle=\sum_{s:|s|=|f(x)|}\left|\widetilde{0^{|f(x)|}}\right\rangle|\tilde{s}\rangle \otimes\left\langle s \mid \phi_{g}^{p}(x)\right\rangle$. Sử dụng hàm lượng tử $C O P Y_{2}$ được đưa ra trong Bổ đề 5.2, chúng ta sau đó sao chép mỗi qustring $|\tilde{s}\rangle$ của $\left|\eta_{f, x}\right\rangle$ lên $\left|\widetilde{0^{|f(x)|}}\right\rangle$ và tạo ra một trạng thái lượng tử có dạng $\sum_{s:|s|=|f(x)|}\left(|\tilde{s}\rangle|\tilde{s}\rangle \otimes\left\langle s \mid \phi_{g}^{p}(x)\right\rangle\right)$. Sau đó, chúng ta giải mã lần xuất hiện đầu tiên và lần xuất hiện thứ hai của $|\tilde{s}\rangle$ thành $\left|0^{|f(x)|+1} 1\right\rangle|s\rangle$. Vì tiện lợi, chúng ta biến đổi lần xuất hiện đầu tiên của $\left|0^{|f(x)|+1} 1\right\rangle|s\rangle$ thành $|s\rangle\left|0^{|f(x)|+1} 1\right\rangle$. Trong phần tiếp theo, chúng ta viết tắt $|s\rangle\left\langle s \mid \phi_{g}^{p}(x)\right\rangle$ là $\left|\zeta_{g}^{p}[s]\right\rangle$. Cuối cùng, chúng ta áp dụng cục bộ $g^{-1}$ vào phần cuối $\left|\zeta_{g}^{p}[s]\right\rangle$, tạo ra $\left|\xi_{x}\right\rangle=\sum_{s:|s|=|f(x)|}\left(|s\rangle\left|\left(0^{|f(x)|+1} 1\right)^{2}\right\rangle \otimes g^{-1}\left(\left|\zeta_{g}^{p}[s]\right\rangle\right)\right)$.

Cho $\left|\Psi_{f(x)}\right\rangle=|f(x)\rangle\left|\left(0^{|f(x)|+1} 1\right)^{2}\right\rangle\left|\phi^{p}(x)\right\rangle$. Vì $\left|\phi_{g}^{p}(x)\right\rangle=\sum_{s:|s|=|f(x)|}\left|\zeta_{g}^{p}[s]\right\rangle$ và $g^{-1}\left(\left|\phi_{g}^{p}(x)\right\rangle\right)= \left|\phi^{p}(x)\right\rangle$, ta suy ra $\left|\phi^{p}(x)\right\rangle=\sum_{s:|s|=|f(x)|} g^{-1}\left(\left|\zeta_{g}^{p}[s]\right\rangle\right)$. Do đó, $\left|\Psi_{f(x)}\right\rangle$ trùng khớp với $\sum_{s:|s|=|f(x)|}\left(|f(x)\rangle\left|\left(0^{|f(x)|+1} 1\right)^{2}\right\rangle \otimes g^{-1}\left(\left|\zeta_{g}^{p}[s]\right\rangle\right)\right)$. Hãy xem xét tích vô hướng $\left\langle\Psi_{f(x)} \mid \xi_{x}\right\rangle$. Bằng một phép tính đơn giản, $\left\langle\Psi_{f(x)} \mid \xi_{x}\right\rangle$ bằng $\sum_{s:|s|=|f(x)|}\langle f(x) \mid s\rangle \otimes \tau_{x}(s)$, trong đó $\tau_{x}(s)$ là tích vô hướng giữa $\left|\phi^{p}(x)\right\rangle$ và $g^{-1}\left(\left|\zeta_{g}^{p}[s]\right\rangle\right)$. Vì $g^{-1}$ thuộc về $\widehat{\square_{1}^{\mathrm{QP}}}$ và bảo toàn chiều cũng như bảo toàn chuẩn theo Mệnh đề 3.5 và Bổ đề 3.4, $\tau_{x}(s)$ bằng $\left\langle\zeta_{g}^{p}[s] \mid \zeta_{g}^{p}[s]\right\rangle$. Vì $s$ bị ép phải lấy giá trị $f(x)$ trong $\left\langle\Psi_{f(x)} \mid \xi_{x}\right\rangle$, ta có được $\left\langle\Psi_{f(x)} \mid \xi_{x}\right\rangle=\left\langle\zeta_{g}^{p}[f(x)] \mid \zeta_{g}^{p}[f(x)]\right\rangle$, bằng $\left\|\left\langle f(x) \mid \phi_{g}^{p}(x)\right\rangle\right\|^{2}$. Do đó, ta kết luận rằng $\left|\left\langle\Psi_{f(x)} \mid \xi_{x}\right\rangle\right|^{2}=\left\|\left\langle f(x) \mid \phi_{g}^{p}(x)\right\rangle\right\|^{4} \geq\left(1-\varepsilon^{\prime}\right)^{2}=1-\varepsilon$, như yêu cầu.\\
($3 \Rightarrow 1$) Vì $f$ có giới hạn đa thức, hãy lấy một đa thức $p$ sao cho $|f(x)| \leq p(|x|)$ đúng với tất cả các chuỗi\\
$x$. Vì chúng ta không biết độ dài của $f(x)$, chúng ta muốn mở rộng $f$ bằng cách đặt $f_{1}(x)=f(x) 01^{p(|x|)+2-|f(x)|}$ để $\left|f_{1}(x)\right|=p(|x|)+2$ cho tất cả các chuỗi $x$. Cố định $\varepsilon \in[0,1 / 2)$. Hãy giả sử rằng tồn tại một hàm $g_{1} \in \widehat{\square_{1}^{\mathrm{QP}}}$ và một đa thức $p_{1}$ thỏa mãn $\left|f_{1}(x)\right| \leq p_{1}(|x|)$ và $\left|\left\langle\Psi_{f_{1}(x)} \mid \phi_{g}^{p_{1}, f_{1}}(x)\right\rangle\right|^{2} \geq 1-\varepsilon$ cho tất cả các trường hợp $x \in\{0,1\}^{*}$.

Sử dụng Bổ đề 4.2 cho hàm lượng tử $g_{1}$, chúng ta có thể lấy một QTM băng đơn, thời gian đa thức, bảo thủ, thuần túy có biên độ $\tilde{\mathbb{C}}$ cho $M$ mà $M$ trên đầu vào $|\phi\rangle$ tạo ra đầu ra sạch của $g_{1}(|\phi\rangle)$ trên băng của nó. Chúng ta xem xét máy sau. Trên đầu vào $x \in \Sigma^{*}$, trước tiên chúng ta tính giá trị $p(|x|)$ một cách xác định và tạo ra $\left|\phi^{p, f_{1}}\right\rangle=\left|0^{p(|x|)+2} 1\right\rangle\left|0^{p(|x|)} 10^{9 p(|x|)} 1\right\rangle \otimes|x\rangle$. Sau đó, chúng ta chạy $M$ trên $\left|\phi^{p, f_{1}}(x)\right\rangle$ để tạo ra $\left|\phi_{g_{1}}^{p, f_{1}}(x)\right\rangle$. Vì $\left|\left\langle\Psi_{f_{1}(x)} \mid \phi_{g_{1}}^{p, f_{1}}(x)\right\rangle\right|^{2} \geq 1-\varepsilon,\left|\phi_{g_{1}}^{p, f_{1}}(x)\right\rangle$ chứa $f_{1}(x)$ với xác suất ít nhất $1-\varepsilon$. Từ $f_{1}(x)$, chúng ta trích xuất $f(x)$ và xuất nó. Điều này kết luận rằng $f$ thuộc $F B Q P$.

\subsection*{4.2 Định lý dạng chuẩn lượng tử}
Các bổ đề chính của chúng tôi, Bổ đề 4.2 và 4.3, tiếp tục dẫn đến một phiên bản lượng tử của định lý dạng chuẩn của Kleene 19, 20, khẳng định sự tồn tại của một vị từ đệ quy nguyên thủy $T(e, x, y)$ và một hàm đệ quy nguyên thủy $U(y)$ sao cho, đối với bất kỳ hàm đệ quy $f(x)$ nào, một chỉ số thích hợp (gọi là số Gödel) $e \in \mathbb{N}$ thỏa mãn $f(x)=U(\mu y . T(e, x, y))$ cho tất cả các đầu vào $x \in \mathbb{N}$, trong đó $\mu$ là toán tử tối thiểu hóa. Phát biểu này về cơ bản tương đương với sự tồn tại của máy Turing phổ quát 32. Ở đây, chúng tôi muốn chứng minh một dạng yếu hơn của định lý dạng chuẩn lượng tử bằng cách sử dụng các Bổ đề 4.2, 4.3.\\
Định lý 4.4 (Định lý dạng chuẩn lượng tử) Tồn tại một hàm lượng tử $f$ trong $\square_{1}^{\mathrm{QP}}$ sao cho, đối với bất kỳ hàm lượng tử $g$ nào trong $\square_{1}^{\mathrm{QP}}$ và bất kỳ hằng số $\varepsilon \in(0,1 / 2)$ nào, tồn tại một chuỗi nhị phân $e$ và một đa thức $p$ thỏa mãn $\|\left|\psi_{g, x}\right\rangle\left\langle\psi_{g, x}\right|-\operatorname{tr}_{n}\left(\left|\eta_{f, x}\right\rangle\left\langle\eta_{f, x}\right|\right) \|_{\mathrm{tr}} \leq \varepsilon$ cho bất kỳ đầu vào nào $|x\rangle$ với $x \in\{0,1\}^{n}$, trong đó $\left|\psi_{g, x}\right\rangle=g(|x\rangle)$ và $\left|\eta_{f, x}\right\rangle=f\left(|\tilde{e}\rangle\left|0^{p(|x|)} 1\right\rangle|x\rangle\right)$. Hàm như vậy được gọi là phổ quát.

Thuật ngữ bổ sung $\left|0^{p(|x|)} 1\right\rangle$ trong $\left|\eta_{f, x}\right\rangle$ là cần thiết để cung cấp cho $g$ đủ không gian làm việc như trong trường hợp của Định lý 4.1. Để chứng minh định lý, chúng tôi sử dụng thực tế là có một QTM phổ quát, có thể mô phỏng tất cả các QTM băng đơn được hình thành tốt với độ trễ đa thức với bất kỳ độ chính xác mong muốn nào. Một máy phổ quát như vậy đã được xây dựng bởi Bernstein và Vazirani [5, Định lý 7.1] và bởi Nishimura và Ozawa [26, Định lý 4.1]. Chúng tôi nói rằng $M_{1}$ trên đầu vào $x_{1}$ mô phỏng $M_{2}$ trên đầu vào $x_{2}$ với độ chính xác nhiều nhất $\varepsilon$ nếu khoảng cách biến thiên toàn phần giữa hai phân bố xác suất $\left\{\left\|\langle y| U_{M_{1}}^{p_{1}\left(\left|x_{1}\right|\right)}\left|c_{0,1}^{\left(x_{1}\right)}\right\rangle\right\|^{2}\right\}_{y \in\{0,1\} \leq n}$ và $\left\{\left\|\langle y| U_{M_{2}}^{p_{2}\left(\left|x_{2}\right|\right)}\left|c_{0,2}^{\left(x_{2}\right)}\right\rangle\right\|^{2}\right\}_{y \in\{0,1\} \leq n}$ là nhiều nhất $\varepsilon$, trong đó $n=\max \left\{p_{1}\left(\left|x_{1}\right|\right), p_{2}\left(\left|x_{2}\right|\right)\right\}$ và, cho mỗi chỉ số $i \in\{1,2\}$, $U_{M_{i}}$ là toán tử tiến hóa thời gian của $M_{i}, p_{i}(\cdot)$ biểu thị thời gian chạy của $M_{i}, c_{0, i}^{\left(x_{i}\right)}$ là cấu hình lệch ban đầu của $M_{i}$ trên đầu vào $x_{i}$, và $y$ trải rộng trên tất cả các chuỗi đầu ra có thể có của $M_{i}$.

Mệnh đề 4.5 [5, 17, 25, 26] Tồn tại một QTM băng đơn, được hình thành tốt, tĩnh sao cho, đối với mỗi hằng số $\varepsilon \in(0,1)$, một số $t \in \mathbb{N}$, một QTM băng đơn được hình thành tốt $M$ có biên độ $\tilde{\mathbb{C}}$, $U$ trên đầu vào $\langle M, x, t, \varepsilon\rangle$ mô phỏng $M$ trên đầu vào $x$ trong $t$ bước với độ chính xác nhiều nhất $\varepsilon$ với độ trễ của một đa thức trong $t$ và $\log (1 / \varepsilon)$, trong đó $\langle M, x, t, \varepsilon\rangle$ đề cập đến một mã hóa cố định, hiệu quả của bộ tứ ( $M, x, t, \varepsilon$ ). QTM như vậy được gọi là phổ quát.

Hệ số cải tiến $\log (1 / \varepsilon)$ trong Mệnh đề 4.5 được gán cho Kitaev 17 và Solovay (trích dẫn trong 25 , Phụ lục 3]).

Đối với chứng minh Định lý 4.4, chúng ta cần mô phỏng một QTM phổ quát $U$ được cung cấp bởi Mệnh đề 4.5 trên một QTM bảo thủ nhất định ngay cả với độ chính xác thấp hơn. Về hình thức đầu vào được đưa cho $f$, chúng ta cần chia bộ tứ ( $M, x, t, \varepsilon$ ) trong mệnh đề thành ba phần ( $M, \varepsilon$ ), $t$, và $x$ và sau đó sửa đổi $U$ để $U$ có thể nhận bất kỳ đầu vào nào có dạng $|\tilde{e}\rangle\left|0^{t} 1\right\rangle|x\rangle$, trong đó $e=\langle M, \varepsilon\rangle$, và bắt chước $M$ trên $x$ trong thời gian $t$ với độ chính xác nhiều nhất $\varepsilon$. Hơn nữa, chúng ta cần buộc $U$ tạo ra đầu ra sạch bằng cách di chuyển tất cả các ký hiệu không trống xuất hiện trong vùng bên trái của bất kỳ chuỗi đầu ra nào đến nơi khác trong thời gian đa thức trong $t$.

Định lý 4.4 hiện nay được suy ra trực tiếp bằng cách kết hợp các Bổ đề 4.24 .3 và Mệnh đề 4.5.\\
Chứng minh Định lý 4.4. Như đã giải thích ở trên, cho $U$ biểu thị một QTM phổ quát đã được sửa đổi có thể nhận các đầu vào có dạng $|\tilde{e}\rangle\left|0^{t} 1\right\rangle|x\rangle$ cho bất kỳ số $e, t \in \mathbb{N}^{+}$và bất kỳ chuỗi $x \in\{0,1\}^{*}$ nào và tạo ra đầu ra sạch. Cho bất kỳ hàm lượng tử $g \in \square_{1}^{\mathrm{QP}}$ nào, Bổ đề 4.2 đảm bảo sự tồn tại của một đa thức $r$ và một QTM băng đơn, có biên độ $\tilde{\mathbb{C}}$, bảo thủ, thuần túy có đầu ra sạch cho $M$ mà, trên bất kỳ đầu vào nào $x \in\{0,1\}^{n}$, $M$ tạo ra trong $r(n)$ bước một chồng chất $g\left(\left|0^{p(n)} 1\right\rangle|x\rangle\right) \otimes\left|0, q_{f}\right\rangle$ của các cấu hình lệch cuối cùng bao gồm\\
của vùng băng thiết yếu của $M$ và trạng thái nội của $M$. Để chính xác hơn, chúng tôi ký hiệu $\widehat{U}_{M}$ là toán tử tiến hóa thời gian lệch của $M$ và $c_{0, M}^{(x)}$ là cấu hình lệch ban đầu của $M$ trên bất kỳ đầu vào nào $|x\rangle$. Chúng tôi tiếp tục đặt $\left|\zeta_{M, x}\right\rangle$ là $\widehat{U}_{M}^{r(n)}\left|c_{0, M}^{(x)}\right\rangle$ và $g\left(\left|0^{p(n)} 1\right\rangle|x\rangle\right)$ là $\left|\psi_{g, x}\right\rangle$. Lưu ý rằng $\left|\zeta_{M, x}\right\rangle$ bằng $\left|\psi_{g, x}\right\rangle \otimes\left|0, q_{f}\right\rangle$. Vì $\left|\psi_{g, x}\right\rangle\left\langle\psi_{g, x}\right|=\operatorname{tr}_{n}\left(\left|\psi_{g, x}\right\rangle\left\langle\psi_{g, x}\right| \otimes\left|0, q_{f}\right\rangle\left\langle 0, q_{f}\right|\right),\left|\psi_{g, x}\right\rangle\left\langle\psi_{g, x}\right|$ có thể được biểu thị là $\operatorname{tr}_{n}\left(\left|\zeta_{M, x}\right\rangle\left\langle\zeta_{M, x}\right|\right)$.

Ngược lại, chúng tôi ký hiệu $\widehat{U}_{\delta}$ là toán tử tiến hóa thời gian lệch của $U$ và $c_{0, U}^{\left(x_{e}\right)}$ là cấu hình lệch ban đầu của $U$ trên đầu vào $\left|x_{e}\right\rangle$ cho $e=\langle M, \varepsilon\rangle$ và $x_{e}=\tilde{e} 0^{r(|x|)} 1 x$. Cho $m=\left|x_{e}\right|$ và đặt $\left|\zeta_{U, x_{e}}\right\rangle$ là $\widehat{U}_{\delta}^{p(m)}\left|c_{0, U}^{\left(x_{e}\right)}\right\rangle$, được viết là $\sum_{r \in F S C_{U, m}}|U[r]\rangle\left|\xi_{x_{e}, r}\right\rangle$ cho một tập hợp thích hợp $\left\{\left|\xi_{x_{e}, r}\right\rangle\right\}_{r}$ của các trạng thái lượng tử trực giao. Mệnh đề 4.5 sau đó đảm bảo rằng, bằng cách lựa chọn một đa thức $s$ thích hợp, khoảng cách biến thiên toàn phần giữa $\left\{\left\|\langle y| \widehat{U}_{M}^{t}\left|c_{0, M}^{(x)}\right\rangle\right\|^{2}\right\}_{y \in\{0,1\}^{n}}$ và $\left\{\left\|\langle y| \widehat{U}_{\delta}^{s(t)}\left|c_{0, U}^{\left(x_{e}\right)}\right\rangle\right\|^{2}\right\}_{y \in\{0,1\}^{n}}$ là nhiều nhất $\varepsilon$ cho bất kỳ số $t \geq 0$ nào.

Chúng tôi áp dụng Bổ đề 4.3 và sau đó nhận được một hàm lượng tử $\square_{1}^{\mathrm{QP}}$ $f$ sao cho, đối với bất kỳ trạng thái lượng tử $|\phi\rangle$ nào, $f(|\phi\rangle)$ biểu thị kết quả của $\widehat{U}_{\delta}^{s(r(n))}$ được áp dụng cho $|\phi\rangle$. Vì tiện lợi, chúng tôi biểu thị $f\left(\left|x_{e}\right\rangle\right)$ là $\left|\eta_{f, x_{e}}\right\rangle$. Bổ đề 4.3 một lần nữa ngụ ý rằng $\left|\eta_{f, x_{e}}\right\rangle=\sum_{r \in F S C_{U, m}}|U[r]\rangle\left|\widehat{\xi}_{x_{e}, r}\right\rangle$ với $\left\langle\xi_{x_{e}, r} \mid \xi_{x_{e}, r^{\prime}}\right\rangle=\left\langle\widehat{\xi}_{x_{e}, r} \mid \widehat{\xi}_{x_{e}, r^{\prime}}\right\rangle$ và $\left\langle\xi_{x, r} \mid \xi_{x, r^{\prime}}\right\rangle=0$ nếu $r \neq r^{\prime}$ cho tất cả $r, r^{\prime} \in F S C_{U, m}$. Do đó, chúng tôi nhận được $\operatorname{tr}_{n}\left(\left|\eta_{f, x_{e}}\right\rangle\left\langle\eta_{f, x_{e}}\right|\right)=\sum_{r, r^{\prime} \in F S C_{U, m}}|U[r]\rangle\left\langle U\left[r^{\prime}\right]\right|$. $\operatorname{tr}\left(\left|\widehat{\xi}_{x_{e}, r}\right\rangle\left\langle\widehat{\xi}_{x_{e}, r^{\prime}}\right|\right)=\sum_{r, r^{\prime} \in F S C_{U, m}}\left\langle\widehat{\xi}_{x_{e}, r^{\prime}} \mid \widehat{\xi}_{x_{e}, r}\right\rangle|U[r]\rangle\left\langle U\left[r^{\prime}\right]\right|$, bằng $\sum_{r \in F S C_{U, m}} \|\left|\widehat{\xi}_{x_{e}, r}\right\rangle \|^{2}|U[r]\rangle\langle U[r]|$. Tương tự, chúng tôi nhận được $\operatorname{tr}_{n}\left(\left|\zeta_{U, x_{e}}\right\rangle\left\langle\zeta_{U, x_{e}}\right|\right)=\sum_{r \in F S C_{U, m}} \|\left|\xi_{x_{e}, r}\right\rangle \|^{2}|U[r]\rangle\langle U[r]|$. Từ những tính toán đó cùng với $\|\left|\xi_{x_{e}, r}\right\rangle\|=\|\left|\widehat{\xi}_{x_{e}, r}\right\rangle \|$, đẳng thức $\operatorname{tr}_{n}\left(\left|\eta_{f, x_{e}}\right\rangle\left\langle\eta_{f, x_{e}}\right|\right)=\operatorname{tr}_{n}\left(\left|\zeta_{U, x_{e}}\right\rangle\left\langle\zeta_{U, x_{e}}\right|\right)$ được suy ra ngay lập tức.

Chúng tôi do đó kết luận rằng $\|\left|\psi_{g, x}\right\rangle\left\langle\psi_{g, x}\right|-\operatorname{tr}_{n}\left(\left|\eta_{f, x}\right\rangle\left\langle\eta_{f, x}\right|\right)\left\|_{\text {tr }}=\right\| \operatorname{tr}_{n}\left(\left|\zeta_{M, x}\right\rangle\left\langle\zeta_{M, x}\right|\right)-\operatorname{tr}_{n}\left(\left|\zeta_{U, x_{e}}\right\rangle\left\langle\zeta_{U, x_{e}}\right|\right) \|_{\text {tr }}$. Số hạng cuối cùng được giới hạn trên bởi $\left.\sum_{y:|y|=n}| |\left|\langle y| \widehat{U}_{\delta}^{s(r(n))}\right| c_{0, U}^{(x)}\right\rangle\left\|^{2}-\right\|\langle y| \widehat{U}_{M}^{r(n)}\left|c_{0, M}^{(x)}\right\rangle \|^{2} \mid$, rõ ràng là nhiều nhất $\varepsilon$. Do đó, $f$ là phổ quát. \(\square\)

\section*{5 Chứng minh Bổ đề chính}
Để hoàn thành chứng minh Định lý 4.1, chúng ta cần chứng minh bổ đề chính, Bổ đề 4.3. Mục này nhằm cung cấp chứng minh mong muốn của bổ đề. Chứng minh của chúng tôi được lấy cảm hứng từ một kết quả của Yao 40, người đã chứng minh một mô phỏng mạch lượng tử của một QTM.

\subsection*{5.1 Mô phỏng chức năng của QTMs}
Cốt lõi của chứng minh Bổ đề 4.3 là một mô phỏng từng bước trực tiếp về hành vi của một QTM băng đơn, có biên độ $\tilde{\mathbb{C}}$, bảo thủ, thuần túy $M=\left(Q, \Sigma, \Gamma, \delta, q_{0}, Q_{f}\right)$, có đầu ra sạch. Để đơn giản, chúng tôi giả sử rằng $\Sigma=\{0,1\}$ và $\Gamma=\{0,1, b\}$, trong đó $b$ ở đây đại diện cho một ký hiệu băng trống duy nhất, thay vì $\#$ được sử dụng trong các phần trước. Chúng tôi cũng giả sử rằng $Q_{f}=\left\{q_{f}\right\}$ và $Q=\{0,1\}^{\ell}$ cho một số chẵn $\ell>0$ nhất định với $q_{0}=0^{\ell}$ và $q_{f}=1^{\ell}$. Hãy giả sử rằng, bắt đầu với chuỗi đầu vào nhị phân $x$ được viết trên băng đầu vào/công việc/đầu ra đơn, $M$ dừng lại trong nhiều nhất $p(|x|)$ bước, trong đó $p$ là một đa thức thích hợp liên quan chỉ với $M$. Chúng tôi cũng giả sử rằng tất cả các đường tính toán của $M$ trên mỗi đầu vào dừng lại đồng thời. Vì tiện lợi, chúng tôi cũng yêu cầu rằng $p(n)>\ell$ cho bất kỳ $n \in \mathbb{N}$ nào. Điều quan trọng là phải nhớ rằng $M$ cuối cùng dừng lại bằng cách đi vào một trạng thái nội duy nhất $q_{f}$ và làm cho đầu băng ở yên tại ô bắt đầu và không có ký hiệu nào không trống xuất hiện trong vùng bên trái của bất kỳ chuỗi đầu ra nào.

Cho $x=x_{1} x_{2} \cdots x_{n}$ là bất kỳ đầu vào nào được đưa cho $M$, trong đó $x_{i}$ là một bit trong $\{0,1\}$ cho mỗi chỉ số $i \in[n]$. Liên quan đến $p$, một vùng băng thiết yếu của $M$ trên $x$ bao gồm tất cả các ô băng được đánh chỉ số giữa $-p(n)$ và $+p(n)$. Chúng tôi biểu thị nội dung băng của vùng băng thiết yếu như vậy là một chuỗi có dạng $\sigma_{1} \sigma_{2} \cdots \sigma_{2 p(n)+1}$ có độ dài chính xác $2 p(|x|)+1$ trên bảng chữ cái băng $\Gamma=\{0,1, b\}$. Chúng tôi theo dõi sự thay đổi của các ký hiệu này khi $M$ thực hiện các bước di chuyển của nó, trong đó $\sigma_{i}$ là một ký hiệu băng được viết tại ô được đánh chỉ số $i-p(n)-1$ cho mỗi chỉ số $i \in[2 p(n)+1]$.

Vì tất cả các hàm $\square_{1}^{\text {QP }}$ được định nghĩa để xử lý các trạng thái lượng tử trong $\mathcal{H}_{\infty}$, chúng tôi cần mã hóa từng ký hiệu băng và do đó là nội dung băng. Chúng tôi do đó định nghĩa một qustring mới mã hóa đúng cấu hình $\gamma=\left(q, h, \sigma_{1} \sigma_{2} \cdots \sigma_{2 p(n)+1}\right)$ của $M$ là

$$
|q\rangle \otimes\left|s_{1}, \hat{\sigma}_{1}\right\rangle \otimes\left|s_{2}, \hat{\sigma}_{2}\right\rangle \otimes\left|s_{3}, \hat{\sigma}_{3}\right\rangle \otimes \cdots \otimes\left|s_{2 p(n)+1}, \hat{\sigma}_{2 p(n)+1}\right\rangle,
$$

trong đó mỗi $\hat{\sigma}_{i}$ nằm trong $\{\hat{0}, \hat{1}, \hat{b}\}$, mỗi $s_{i} \in\{\hat{2}, \hat{3}\}$ biểu thị sự hiện diện của đầu băng (trong đó $\hat{2}$ có nghĩa là "đầu ở đây" và $\hat{3}$ có nghĩa là "không có đầu ở đây") tại ô $i-p(n)-1$, và $q$ là một trạng thái nội trong $Q$. Trong các phần tiếp theo, chúng tôi gọi một qustring như vậy là mã của cấu hình $\gamma$ của $M$ và biểu thị nó là $|\hat{\gamma}\rangle$. Mã $|\hat{\gamma}\rangle$ này có độ dài $\ell(|\hat{\gamma}\rangle)=8 p(n)+\ell+4$, là số chẵn và lớn hơn $n$.

Cho bất kỳ đầu vào nhị phân $x=x_{1} x_{2} \cdots x_{n}$ có độ dài $n$, hãy nhớ lại từ Mục 4.1 rằng $\left|\phi^{p}(x)\right\rangle= \left|0^{|x|} 1\right\rangle\left|0^{p(|x|)} 1\right\rangle\left|0^{11 p(|x|)+6} 1\right\rangle|x\rangle$ và $\left|\phi^{p, f}(x)\right\rangle=\left|\widetilde{0^{|f(x)|}}\right\rangle\left|0^{|f(x)|+1} 1\right\rangle\left|\phi^{p}(x)\right\rangle$. Ngoại trừ Bước 1) trong Mục 5.2 cũng như tất cả các bước trong Mục 5.4, chúng tôi luôn bỏ qua các chuỗi tiền tố $\widetilde{0^{|f(x)|}} 0^{|f(x)|+1} 10^{|x|} 10^{p(|x|)} 1$ trong $\left|\phi^{p, f}(x)\right\rangle$ và $0^{|x|} 10^{p(|x|)} 1$ trong $\left|\phi^{p}(x)\right\rangle$, và chúng tôi chú ý đến các qubit còn lại. Hàm lượng tử mong muốn $g$ sẽ được xây dựng từng bước thông qua các Mục 5.2 5.5.

\subsection*{5.2 Xây dựng một cấu hình ban đầu đã mã hóa}
Cho một đầu vào nhị phân $x=x_{1} x_{2} \cdots x_{n}$, cấu hình ban đầu $\gamma_{0}$ của $M$ trên $x$ có dạng ( $q_{0}, 0, b \cdots b x b \cdots b$ ), và do đó mã $\left|\hat{\gamma}_{0}\right\rangle$ của $\gamma_{0}$ phải có dạng

$$
\left|q_{0}\right\rangle \otimes|\hat{3}, \hat{b}\rangle \otimes \cdots \otimes|\hat{3}, \hat{b}\rangle \otimes\left|\hat{2}, \hat{x}_{1}\right\rangle \otimes\left|\hat{3}, \hat{x}_{2}\right\rangle \otimes\left|\hat{3}, \hat{x}_{3}\right\rangle \otimes \cdots \otimes\left|\hat{3}, \hat{x}_{n}\right\rangle \otimes|\hat{3}, \hat{b}\rangle \otimes \cdots \otimes|\hat{3}, \hat{b}\rangle
$$

trong đó $\hat{x}_{i}$ là mã của $x_{i}, q_{0}\left(=0^{\ell}\right)$ là trạng thái nội ban đầu, và $x_{1}$ nằm ở ô 0 . Lưu ý rằng $\ell\left(\left|\hat{\gamma}_{0}\right\rangle\right)= 8 p(n)+\ell+4$. Trong phần tiếp theo, chúng tôi sẽ chỉ ra cách tạo mã cụ thể $\left|\hat{\gamma}_{0}\right\rangle$ này từ trạng thái lượng tử $\left|\phi^{p}(x)\right\rangle$. Để đơn giản, chúng tôi sẽ bỏ qua thuật ngữ $\left|\hat{q}_{0}\right\rangle$ trong các bước sau ngoại trừ Bước 8).

\begin{enumerate}
  \item Bắt đầu với đầu vào $\left|\phi^{p}(x)\right\rangle$, trước tiên chúng ta biến đổi nó thành $\left|0^{n} 1\right\rangle\left|0^{p(n)} 1\right\rangle\left|10^{11 p(n)+5} 1\right\rangle|x\rangle$ bằng một hàm lượng tử $h_{1}=\operatorname{Skip}[N O T]$, thỏa mãn cả hai $h_{1}\left(\left|0^{m} 1\right\rangle|\psi\rangle\right)=\left|0^{m} 1\right\rangle \otimes N O T(|\psi\rangle)$ và $h_{1}\left(\left|0^{m+1}\right\rangle\right)=\left|0^{m+1}\right\rangle$ cho bất kỳ số $m \in \mathbb{N}$ nào và bất kỳ trạng thái lượng tử $|\psi\rangle \in \mathcal{H}_{\infty}$. Bổ đề 3.9 đảm bảo rằng $h_{1}$ thực sự tồn tại trong $\widehat{\square_{1}^{\mathrm{QP}}}$.
\end{enumerate}

Từ $\left|0^{p(n)} 1\right\rangle\left|10^{11 p(n)+5} 1\right\rangle$, chúng ta muốn tạo ra $\left|0^{p(n)} 1\right\rangle\left|0^{p(n)} 1\right\rangle\left|0^{10 p(n)+5} 1\right\rangle$ bằng một hàm lượng tử $\widehat{\square_{1}^{\mathrm{QP}}}$ thích hợp $f_{1}$. Để định nghĩa hàm lượng tử mong muốn $f_{1}$, trước tiên chúng ta xây dựng một hàm lượng tử khác $g_{1}$ ánh xạ $\left|0^{p(n)} 1\right\rangle\left|0^{m} 10^{k} 1\right\rangle$ thành $\left|0^{p(n)} 1\right\rangle\left|0^{m+1} 10^{k-1} 1\right\rangle$ cho bất kỳ hai số nguyên $m \geq 0$ và $k \geq 1$ nào bằng cách đặt

$$
g_{1}(|\phi\rangle)= \begin{cases}|\phi\rangle & \text { if } \ell(|\phi\rangle) \leq 1 \\ |0\rangle \otimes g_{1}(\langle 0 \mid \phi\rangle)+|1\rangle \otimes g_{2}(\langle 1 \mid \phi\rangle) & \text { otherwise }\end{cases}
$$

trong đó $g_{2}$ được giới thiệu như

$$
g_{2}(|\phi\rangle)= \begin{cases}|\phi\rangle & \text { if } \ell(|\phi\rangle) \leq 1 \\ S W A P\left(|0\rangle \otimes g_{2}(\langle 0 \mid \phi\rangle)+|1\rangle\langle 1 \mid \phi\rangle\right) & \text { otherwise }\end{cases}
$$

Một cách chính thức hơn, ta đặt $g_{2}=Q R e c_{1}\left[I, S W A P, I \mid g_{2}, I\right]$ và đặt $\hat{g}_{2}=\operatorname{Branch}\left[I, g_{2}\right]$. Sau đó ta đặt $g_{1}= Q \operatorname{Rec}_{1}\left[I, I, \hat{g}_{2} \mid g_{1}, I\right]$. Hàm lượng tử $f_{1}$ cuối cùng được định nghĩa là $f_{1}=Q \operatorname{Rec}_{1}\left[I, g_{1}, I \mid f_{1}, I\right]$; cụ thể,

$$
f_{1}(|\phi\rangle)= \begin{cases}|\phi\rangle & \text { if } \ell(|\phi\rangle) \leq 1 \\ g_{1}\left(|0\rangle \otimes f_{1}(\langle 0 \mid \phi\rangle)+|1\rangle\langle 1 \mid \phi\rangle\right) & \text { otherwise }\end{cases}
$$

Theo cách tương tự, ta tiếp tục biến đổi $\left|0^{p(n)} 1\right\rangle\left|0^{10 p(n)+5} 1\right\rangle$ thành $\left|0^{p(n)} 1\right\rangle\left|0^{2 p(n)+2} 1\right\rangle\left|0^{8 p(n)+2} 1\right\rangle$. Sau đó ta đổi $\left|0^{n} 1\right\rangle\left|0^{2 p(n)+2} 1\right\rangle$ thành $\left|0^{n} 1\right\rangle\left|0^{n} 1\right\rangle\left|0^{2 p(n)-n+1} 1\right\rangle$ và $\left|0^{n} 1\right\rangle\left|0^{8 p(n)+2} 1\right\rangle$ thành $\left|0^{n} 1\right\rangle\left|0^{n} 1\right\rangle\left|0^{8 p(n)-n+1} 1\right\rangle$. Tổng thể, đầu vào $\left|\phi^{p}(x)\right\rangle$ được biến đổi thành $\left|0^{n} 1\right\rangle\left|\left(0^{p(n)} 1\right)^{3}\right\rangle\left|\left(0^{n} 1\right)^{2}\right\rangle\left|0^{2 p(n)-n+1} 1\right\rangle\left|0^{8 p(n)-n+1} 1\right\rangle|x\rangle$.\\
(*) Trong phần tiếp theo, bằng cách bỏ qua $\left|0^{n} 1\right\rangle\left|\left(0^{p(n)} 1\right)^{3}\right\rangle\left|\left(0^{n} 1\right)^{2}\right\rangle$, chúng ta tạm thời giả sử rằng đầu vào của chúng ta là $\left|0^{8 p(n)-n+1} 1\right\rangle|x\rangle$.\\
2) Để dễ đọc, chúng ta giải thích bước này bằng một ví dụ minh họa của $\left|0^{6} 1\right\rangle\left|x_{1} x_{2} x_{3}\right\rangle$, mà chúng ta muốn biến đổi thành $|00\rangle\left|\hat{3} \hat{x}_{1} \hat{x}_{2} \hat{x}_{3}\right\rangle$. Vì mục đích này, chúng ta bắt đầu bằng cách thay đổi $\left|0^{6} 1 x_{1} x_{2} x_{3}\right\rangle$ thành $\left|x_{3} x_{2} x_{1} 10^{6}\right\rangle$ bằng cách áp dụng REVERSE.

Để có được $\left|x_{3} 0 x_{2} 0 x_{1} 001\right\rangle|00\rangle$ từ $\left|x_{3} x_{2} x_{1}\right\rangle\left|10^{6}\right\rangle$, chúng ta tiếp tục áp dụng $g_{3}=S W A P \circ R E P_{1}$, biến đổi $\left|x_{3} x_{2} x_{1} 10^{6}\right\rangle$ thành $\left|x_{3} 0 x_{2} x_{1} 10^{5}\right\rangle$. Chúng ta lặp lại việc áp dụng $g_{3}$ và biến đổi $\left|x_{3} x_{2} x_{1}\right\rangle\left|10^{6}\right\rangle$ thành $\left|x_{3} 0 x_{2} 0 x_{1} 0\right\rangle\left|10^{3}\right\rangle$. Quá trình này có thể được thực hiện bằng hàm lượng tử $h_{3}=2 Q R e c_{2}\left[I, h^{\prime}, g_{3} \mid\left\{h_{s}^{\prime \prime}\right\}_{s \in\{0,1\}^{2}}\right]$ được định nghĩa bởi phép đệ quy lượng tử 2 qubit, trong đó $h^{\prime}=$ Branch $_{2}\left[\left\{h_{s}^{\prime}\right\}_{s \in\{0,1\}^{2}}\right]$ với $h_{a 1}^{\prime}=S W A P$ và $h_{a 0}^{\prime}=I$ cũng như $h_{a 1}^{\prime \prime}=I$ và $h_{a 0}^{\prime \prime}=h_{3}$ cho mỗi bit $a \in\{0,1\}$; cụ thể,

$$
h_{3}(|\phi\rangle)= \begin{cases}|\phi\rangle & \text { if } \ell(|\phi\rangle) \leq 2, \\ \sum_{a \in\{0,1\}}\left(S W A P\left(|a 1\rangle\left\langle a 1 \mid \psi_{g_{3}, \phi}\right\rangle\right)+|a 0\rangle \otimes h_{3}\left(\left\langle a 0 \mid \psi_{g_{3}, \phi}\right\rangle\right)\right) & \text { otherwise },\end{cases}
$$

trong đó $\left|\psi_{g_{3}, \phi}\right\rangle$ đại diện cho $g_{3}(|\phi\rangle)$. Chúng ta tiếp tục áp dụng $g_{4}=R E V E R S E \circ h_{3}$ để biến đổi $\left|x_{3} x_{2} x_{1}\right\rangle\left|10^{6}\right\rangle$ thành $\left|0^{3} 1\right\rangle\left|\hat{x}_{1} \hat{x}_{2} \hat{x}_{3}\right\rangle$. Trong trường hợp tổng quát, quá trình trên biến đổi $\left|0^{n} 1\right\rangle\left|x_{1} x_{2} \cdots x_{m}\right\rangle$ thành $|1\rangle\left|\hat{x}_{1} \hat{x}_{2} \cdots \hat{x}_{n}\right\rangle$.

Cuối cùng, chúng ta áp dụng $g_{5}$, biến đổi $\left|0^{3} 1\right\rangle\left|\hat{x}_{1} \hat{x}_{2} \hat{x}_{3}\right\rangle$ thành $|00\rangle\left|\hat{3} \hat{x}_{1} \hat{x}_{2} \hat{x}_{3}\right\rangle$, được định nghĩa bởi

$$
g_{5}(|\phi\rangle)= \begin{cases}|\phi\rangle & \text { if } \ell(|\phi\rangle) \leq 1 \\ S W A P\left(|00\rangle \otimes g_{5}(\langle 00 \mid \phi\rangle)+\sum_{y \in\{0,1\}^{2}-\{00\}}(|y\rangle\langle y \mid \phi\rangle)\right) & \text { otherwise }\end{cases}
$$

\subsection*{5.3 Mô phỏng một bước đơn lẻ}
Để mô phỏng toàn bộ quá trình tính toán của $M$ trên bất kỳ đầu vào $x$ nào, chúng ta cần mô phỏng từng bước của $M$ một cách tuần tự cho đến khi $M$ cuối cùng đi vào trạng thái nội $q_{f}$. Trong phần tiếp theo, chúng ta sẽ trình bày cách mô phỏng một bước đơn lẻ của $M$ bằng cách thay đổi vị trí đầu đọc, ký hiệu băng và trạng thái nội tại trong một cấu hình xiên đã cho.

Lưu ý rằng bước của $M$ chỉ liên quan đến ba ô liên tiếp, một trong số đó đang được đầu đọc quét. Để mô tả ba ô liên tiếp cùng với trạng thái nội tại của $M$, nói chung, chúng ta sử dụng biểu thức $r$ có dạng $p s_{1} \sigma_{1} s_{2} \sigma_{2} s_{3} \sigma_{3}$ sử dụng $p \in\{0,1\}^{\ell}, \sigma_{i} \in\{0,1, b\}$, và $s_{i} \in\{2,3\}$ cho mỗi chỉ số $i \in[3]$. Mỗi biểu thức với $s_{i}=2$ cho biết $M$ đang ở trạng thái $q$, quét ô thứ $i$ trong ba ô. Lưu ý rằng độ dài của $r$ là $\ell+6$. Cho $T$ là tập hợp tất cả các $r$ có thể có như vậy. Lưu ý rằng $T$ là một tập hợp hữu hạn. Mã $|\hat{r}\rangle$ của $r=p s_{1} \sigma_{1} s_{2} \sigma_{2} s_{3} \sigma_{3}$ là $|q\rangle\left|\hat{s}_{1}, \hat{\sigma}_{1}\right\rangle\left|\hat{s}_{2}, \hat{\sigma}_{2}\right\rangle\left|\hat{s}_{3}, \hat{\sigma}_{3}\right\rangle$, có độ dài $\ell+12$.

Để tiện lợi, chúng ta gọi $r$ là mục tiêu nếu $s_{1}=3, s_{2}=2, s_{3}=3$, và $\delta\left(q, \sigma_{2}\right)$ được xác định. Để sử dụng về sau, $s_{1} \sigma_{1} s_{2} \sigma_{2} s_{3} \sigma_{3}$ không có $q$ được gọi là tiền mục tiêu nếu $q s_{1} \sigma_{1} s_{2} \sigma_{2} s_{3} \sigma_{3}$ là một mục tiêu.

\begin{enumerate}
  \item Cho $r=q s_{1} \sigma_{1} s_{2} \sigma_{2} s_{3} \sigma_{3}$ là bất kỳ phần tử cố định nào trong $T$. Chúng ta chuẩn bị một qubit đánh dấu $|0\rangle$ ở cuối $|\phi\rangle$ để đánh dấu rằng việc mô phỏng một bước đơn lẻ của $M$ đang được thực hiện hoặc đã được hoàn thành. Hãy định nghĩa một hàm lượng tử $f_{8}$, biến đổi $|\tilde{r}\rangle|0\rangle$ thành $|\tilde{r}\rangle|0\rangle$ hoặc $|\tilde{s}\rangle|1\rangle$, nơi $s$ là biểu thức thu được từ $r$ bằng cách áp dụng $\delta$ một lần nếu $r$ là một mục tiêu. Cho $A$ là tập hợp tất cả các mục tiêu trong $T$ và đặt $A^{\diamond}=\{\tilde{r} \mid r \in A\}$. Tương tự, chúng ta định nghĩa $T^{\diamond}$ từ $T$.
\end{enumerate}

Vì mục đích này, trước tiên chúng ta định nghĩa một hàm lượng tử hỗ trợ $h_{r}(|b\rangle|\psi\rangle)=N O T(|b\rangle) \otimes|\tilde{r}\rangle\langle\tilde{r} \mid \psi\rangle+ \sum_{s \in\{0,1\}^{\ell+12}-A^{\circ}}(|b\rangle \otimes|s\rangle\langle s \mid \psi\rangle)$ cho bất kỳ mục tiêu $r \in T$ nào và bất kỳ $b \in\{0,1\}$. Tiếp theo, chúng ta định nghĩa $\left\{g_{r}\right\}_{r \in T}$ như sau. Nếu $r$ là một phi mục tiêu trong $T$, thì chúng ta đặt $g_{r}=I$. Trong phần tiếp theo, chúng ta giả định rằng $r$ là một mục tiêu. Cho $g_{r}(|0\rangle|s\rangle|\phi\rangle)=|0\rangle|s\rangle|\phi\rangle$ cho bất kỳ chuỗi $s \in\{0,1\}^{\ell+12}$. Vì $M$ là dạng chuẩn, $M$ chỉ có hai loại chuyển đổi, được thể hiện trong (i) và (ii) dưới đây.\\
(i) Xét trường hợp $\delta\left(q, \sigma_{2}\right)=e^{i \theta}\left|q^{\prime}, \tau, d\right\rangle$. Khi $d=L$, chúng ta đặt

$$
g_{r}\left(|1\rangle|q\rangle\left|\hat{3}, \hat{\sigma}_{1}\right\rangle\left|\hat{2}, \hat{\sigma}_{2}\right\rangle\left|\hat{3}, \hat{\sigma}_{3}\right\rangle|\phi\rangle\right)=e^{i \theta}|1\rangle\left|q^{\prime}\right\rangle\left|\hat{2}, \hat{\sigma}_{1}\right\rangle|\hat{3}, \hat{\tau}\rangle\left|\hat{3}, \hat{\sigma}_{3}\right\rangle|\phi\rangle .
$$

Khi $d=R$, ngược lại, chúng ta định nghĩa

$$
g_{r}\left(|1\rangle|q\rangle\left|\hat{3}, \hat{\sigma}_{1}\right\rangle\left|\hat{2}, \hat{\sigma}_{2}\right\rangle\left|\hat{3}, \hat{\sigma}_{3}\right\rangle|\phi\rangle\right)=e^{i \theta}|1\rangle\left|q^{\prime}\right\rangle\left|\hat{3}, \hat{\sigma}_{1}\right\rangle|\hat{3}, \hat{\tau}\rangle\left|\hat{2}, \hat{\sigma}_{3}\right\rangle|\phi\rangle .
$$

(ii) Xét trường hợp $\delta\left(q, \sigma_{2}\right)=\cos \theta\left|q_{1}, \tau_{1}, d_{1}\right\rangle+\sin \theta\left|q_{2}, \tau_{2}, d_{2}\right\rangle$. Nếu ( $d_{1}, d_{2}$ ) = ( $R, L$ ), thì chúng ta định nghĩa $g_{r}$ là

$$
g_{r}\left(|1\rangle|q\rangle\left|\hat{3}, \hat{\sigma}_{1}\right\rangle\left|\hat{2}, \hat{\sigma}_{2}\right\rangle\left|\hat{3}, \hat{\sigma}_{3}\right\rangle|\phi\rangle\right)=\cos \theta|1\rangle\left|q_{1}\right\rangle\left|\hat{3}, \hat{\sigma}_{1}\right\rangle\left|\hat{3}, \hat{\tau}_{1}\right\rangle\left|\hat{2}, \hat{\sigma}_{3}\right\rangle|\phi\rangle+\sin \theta|1\rangle\left|q_{2}\right\rangle\left|\hat{2}, \hat{\sigma}_{1}\right\rangle\left|\hat{3}, \hat{\tau}_{2}\right\rangle\left|\hat{3}, \hat{\sigma}_{3}\right\rangle|\phi\rangle .
$$

Các giá trị khác của ( $d_{1}, d_{2}$ ) được xử lý tương tự.\\
Lưu ý rằng $\left\{g_{r}(|1\rangle|\tilde{r}\rangle)\right\}_{r \in A}$ tạo thành một tập hợp trực chuẩn vì $M$ được hình thành tốt, và do đó $\delta$ thỏa mãn tất cả các điều kiện được nêu trong Mục 2.2. Một khi qubit đánh dấu trở thành $|1\rangle$, chúng ta không cần áp dụng $g_{r} \circ h_{r}$. Do đó, chúng ta tiếp tục đặt $g_{r}^{\prime}=\operatorname{Branch}\left[g_{r} \circ h_{r}, I\right]$. Bằng cách kết hợp tất cả các $g_{r}^{\prime}$, chúng ta định nghĩa $g$ là $g=\operatorname{Compo}\left[\left\{g_{r}^{\prime}\right\}_{r \in T}\right]$, ngụ ý rằng

$$
g(|0\rangle|\phi\rangle)=\sum_{r \in T}\left(g_{r}(|1\rangle|\tilde{r}\rangle\langle\tilde{r} \mid \phi\rangle)\right)+\sum_{s \in\{0,1\}^{\ell+12}-T^{\diamond}}(|0\rangle|s\rangle\langle s \mid \phi\rangle)
$$

cho bất kỳ trạng thái lượng tử $|\phi\rangle \in \mathcal{H}_{\infty}$. Chúng ta có thể khẳng định rằng $g$ bảo toàn chuẩn và có thể được xây dựng từ các hàm lượng tử ban đầu trong Định nghĩa 3.1 bằng cách áp dụng các quy tắc xây dựng. Từ khẳng định này, $g$ thuộc về $\widehat{\square_{1}^{\mathrm{QP}}}$. Cuối cùng, chúng ta định nghĩa $f_{8}=R E M O V E_{1} \circ\left(g^{\leq \ell+13} \otimes I\right)$, nơi $g^{\leq \ell+13} \otimes I$ được định nghĩa như trong Bổ đề 3.8. Về trực giác, $f_{8}$ thay đổi nội dung của ba ô băng liên tiếp mà ô ở giữa đang được quét bởi đầu đọc băng.\\
2) Chúng ta muốn áp dụng $f_{8}$ cho tất cả ba ô băng liên tiếp. Trước tiên, chúng ta tìm một mã của một tiền mục tiêu $\hat{s}_{1} \hat{\sigma}_{1} \hat{s}_{2} \hat{\sigma}_{2} \hat{s}_{3} \hat{\sigma}_{3}$ (sử dụng phép đệ quy lượng tử), di chuyển một trạng thái nội $q$ cũng như một dấu hiệu $|0\rangle$ về phía trước (bằng $\left.R E P_{\ell+1}\right)$ để tạo ra một khối có dạng $|0\rangle|q\rangle\left|\hat{s}_{1} \hat{\sigma}_{1} \hat{s}_{2} \hat{\sigma}_{2} \hat{s}_{3} \hat{\sigma}_{3}\right\rangle$ và áp dụng $f_{8}$ cho khối này. Điều này thay đổi $|0\rangle$ thành $|1\rangle$ để ghi lại việc thực hiện quy trình hiện tại. Sau đó, chúng ta di chuyển trạng thái nội và dấu hiệu đã thu được trở lại cuối (bằng $R E M O V E_{\ell+1}$ ). Toàn bộ quy trình này có thể được thực hiện bởi một hàm lượng tử phù hợp, ví dụ, $F_{2}$.

Một cách chính thức hơn, trước tiên chúng ta đặt một hàm lượng tử khác $p$ là $R E M O V E_{\ell+1} \circ L E N G T H_{\ell+13}\left[f_{8}\right] \circ R E P_{\ell+1}$. Hàm lượng tử $F_{2}^{\prime}$ sau đó được định nghĩa là

$$
F_{2}^{\prime}(|\phi\rangle|0\rangle|q\rangle)= \begin{cases}|\phi\rangle|0\rangle|q\rangle & \text { nếu } \ell(|\phi\rangle)<\ell+13 \\ p\left(\sum_{s:|s|=4}\left(|s\rangle \otimes F_{2}^{\prime}(\langle s \mid \phi\rangle|0\rangle|q\rangle)\right)\right) & \text { ngược lại. }\end{cases}
$$

Để hoàn thành biến đổi, chúng ta tiếp tục định nghĩa $F_{2}=R E P_{\ell+1} \circ F_{2}^{\prime} \circ R E M O V E_{\ell+1}$. Sau khi áp dụng $F_{2}$, qubit đầu tiên của $|0\rangle|q\rangle|\phi\rangle$ chuyển thành $|1\rangle$, đánh dấu việc thực hiện một bước đơn lẻ của $M$. Nếu chúng ta muốn lặp lại việc áp dụng $F_{2}$, chúng ta cần đặt lại $|1\rangle$ về $|0\rangle$.

\subsection*{5.4 Hoàn thành toàn bộ mô phỏng}
Chúng ta đã chỉ ra trong Mục 5.3 cách mô phỏng một bước đơn lẻ của $M$ trên $x$ bằng $F_{2}$. Ở đây, chúng ta muốn mô phỏng tất cả các bước của $M$ bằng cách áp dụng $F_{2}$ một cách quy nạp cho bất kỳ đầu vào nào có dạng $\left|0^{p(|x|)} 1\right\rangle \otimes|\psi\rangle$, nơi $|\psi\rangle$ biểu thị một chồng chất của các cấu hình xiên được mã hóa của $M$ trên $x$. Quá trình này được thực hiện bởi một hàm lượng tử mới $F_{3}$, lặp lại việc áp dụng $N O T \circ F_{2}$ cho $|0\rangle|\psi\rangle, p(|x|)$ lần, nơi NOT đặt lại $|1\rangle$ về $|0\rangle$.

Trước tiên, chúng ta lấy một hàm lượng tử $\hat{g}=\operatorname{Skip}\left[N O T \circ F_{2}\right]$ bằng cách sử dụng Bổ đề 3.9. Hàm lượng tử mong muốn $F_{3}$ phải thỏa mãn

$$
F_{3}(|\phi\rangle)= \begin{cases}|\phi\rangle & \text { nếu } \ell(|\phi\rangle) \leq 1 \\ |0\rangle \otimes \hat{g}\left(F_{3}(\langle 0 \mid \phi\rangle)\right)+|1\rangle \otimes I(\langle 1 \mid \phi\rangle) & \text { ngược lại. }\end{cases}
$$

Lưu ý rằng số lần áp dụng $\hat{g}$ chính xác là $p(|x|)$. Hàm $F_{3}$ này có thể được hiện thực hóa bằng cách sử dụng phép đệ quy lượng tử một qubit có dạng $F_{3}=Q R e c_{1}\left[I, \operatorname{Branch}[\hat{g}, I], I \mid F_{3}, I\right]$.

\subsection*{5.5 Chuẩn bị đầu ra}
Giả sử rằng một chồng chất của các cấu hình xiên được mã hóa cuối cùng của $M$ trên đầu vào đã cho $x$ có độ dài $n$ có dạng $\sum_{r \in F S C_{M}(x)}|\widetilde{M[r]}\rangle\left|\xi_{x, r}\right\rangle$ cho một chuỗi nhất định $\left\{\left|\xi_{x, r}\right\rangle\right\}_{r \in F S C_{M, n}}$ của các trạng thái lượng tử, nơi $M[r]$ biểu thị chuỗi đầu ra của $M$ xuất hiện trong một cấu hình xiên cuối cùng $r$ bao gồm chỉ vùng băng thiết yếu của $M$ trên $x$.

Để chỉ ra phần đầu tiên của Bổ đề 4.3, ở cuối mô phỏng, chúng ta cần tạo ra $\widetilde{M[r]}$ ở phần ngoài cùng bên trái của chuỗi lượng tử thu được bằng mô phỏng. Để đạt được mục tiêu này, trước tiên chúng ta cần di chuyển nội dung của vùng bên trái của ô bắt đầu đến cuối vùng băng thiết yếu.

Là một ví dụ minh họa, giả sử rằng chúng ta đã nhận được trạng thái lượng tử $|\hat{3} \hat{b} \hat{3} \hat{b} \hat{b} \hat{b}\rangle|\hat{2} \hat{0} \hat{3} \hat{0} \hat{3} \hat{1}\rangle|\hat{3} \hat{b}\rangle$ sau khi thực hiện các bước trong Mục 5.4. Trong trường hợp này, kết quả của $M$ trong cấu hình xiên cuối cùng này $r$ là 001, và do đó $\widetilde{M[r]}=\hat{0} \hat{0} \hat{1} \hat{2}$ giữ. Trong phần tiếp theo, chúng ta muốn biến đổi $|\hat{3} \hat{b} \hat{3} \hat{b} \hat{3} \hat{b}\rangle|\hat{2} \hat{0} \hat{3} \hat{0} \hat{3} \hat{1}\rangle|\hat{3} \hat{b}\rangle$ thành $|\hat{0} \hat{0} \hat{1} \hat{2}\rangle|\hat{3} \hat{b} \hat{b} \hat{b} \hat{3} \hat{b}\rangle|\hat{1} \hat{3} \hat{3} \hat{3}\rangle$, bằng với $|\widetilde{M[r]}\rangle|\hat{3} \hat{b} \hat{3} \hat{b} \hat{3} \hat{b}\rangle|\hat{1} \hat{3} \hat{3} \hat{3}\rangle$, bằng một hàm lượng tử $\widehat{\square_{1}^{\mathrm{QP}}}$ phù hợp, ví dụ, $F_{4}$.\\
(i) Bắt đầu với $|\hat{3} \hat{b} \hat{3} \hat{b} \hat{3} \hat{b}\rangle|\hat{2} \hat{0} \hat{3} \hat{0} \hat{3} \hat{1}\rangle|\hat{3} \hat{b}\rangle$, để đánh dấu phần cuối của ô băng, chúng ta thay đổi dấu hiệu cuối cùng $\hat{3}$ thành $\hat{1}$ bằng cách áp dụng NOT cho $\hat{3}$ và sau đó nhận được $|\hat{3} \hat{b} \hat{3} \hat{b} \hat{3} \hat{b}\rangle|\hat{2} \hat{0} \hat{3} \hat{0} \hat{3} \hat{1}\rangle|\hat{1} \hat{b}\rangle$. Quá trình này được gọi là $f_{9}$.\\
(ii) Bằng cách lặp lại di chuyển $|\hat{3} \hat{b}\rangle$ ngoài cùng bên trái trong $|\hat{3} \hat{b} \hat{3} \hat{b} \hat{3} \hat{b}\rangle|\hat{2} \hat{0} \hat{3} \hat{0} \hat{3} \hat{1}\rangle|\hat{1} \hat{b}\rangle$ đến cuối chuỗi lượng tử, cuối cùng chúng ta tạo ra $|\hat{2} \hat{0} \hat{3} \hat{0} \hat{3} \hat{1}\rangle|\hat{1} \hat{b} \hat{3} \hat{b} \hat{3} \hat{b} \hat{3} \hat{b}\rangle$. Để hiện thực hóa toàn bộ biến đổi này, chúng ta cần định nghĩa

$$
f_{10}(|\phi\rangle)= \begin{cases}|\phi\rangle & \text { nếu } \ell(|\phi\rangle)<4, \\ \sum_{a \in\{0,1\}}|\hat{2} \hat{a}\rangle\langle\hat{2} \hat{a} \mid \phi\rangle+\sum_{z \in B_{4}} R E M O V E_{4}\left(|z\rangle \otimes f_{10}(\langle z \mid \phi\rangle)\right) & \text { ngược lại },\end{cases}
$$

nơi $B_{4}=\{0,1\}^{4}-\{\hat{2} \hat{0}, \hat{2} \hat{1}\}$. Chính thức hơn, chúng ta định nghĩa $f_{11}$ là $f_{11}=2 Q R e c_{4}\left[I, h^{\prime \prime}, I \mid\left\{f_{z}\right\}_{z \in\{0,1\}^{4}}\right]$, nơi $h^{\prime \prime}=\operatorname{Branch}\left[\left\{h_{z}^{\prime \prime}\right\}_{z \in\{0,1\}^{4}}\right]$ với $h_{\hat{2} \hat{0}}^{\prime \prime}=h_{\hat{2} \hat{1}}^{\prime \prime}=I$ và $h_{z}^{\prime \prime}=R E M O V E_{4}$ cũng như $f_{\hat{2} \hat{0}}=f_{\hat{2} \hat{1}}=I$ và $f_{z}=f_{10}$ cho bất kỳ $z \in B_{4}$.\\
(iii) Chuỗi lượng tử hiện tại có dạng $|\hat{2} \hat{0} \hat{3} \hat{0} \hat{3} \hat{1}\rangle|\hat{1} \hat{b}\rangle|\hat{3} \hat{b} \hat{3} \hat{b} \hat{3} \hat{b}\rangle$. Chúng ta lần lượt loại bỏ từng dấu hiệu $\{\hat{1}, \hat{2}, \hat{3}\}$ trong $|\hat{2} \hat{0} \hat{3} \hat{0} \hat{3} \hat{1}\rangle|\hat{1} \hat{b}\rangle$ đến cuối toàn bộ chuỗi lượng tử và tạo ra $|\hat{0} \hat{0} \hat{1} \hat{b}\rangle|\hat{3} \hat{b} \hat{b} \hat{b} \hat{3} \hat{b}\rangle|\hat{1} \hat{3} \hat{3} \hat{2}\rangle$. Để thực hiện quá trình này, chúng ta áp dụng $h_{8}$ được định nghĩa bởi $h_{8}=2 Q R E C_{3}\left[I, R E M O V E_{2}, I \mid\left\{h_{y}\right\}_{y \in\{0,1\}^{4}}\right]$, nơi $h_{\hat{1} \hat{b}}=I$ và $h_{y}=h_{8}$ cho tất cả $y \in\{0,1\}^{4}-\{\hat{1} \hat{b}\}$; nói cách khác,

$$
h_{8}(|\phi\rangle)= \begin{cases}|\phi\rangle & \text { nếu } \ell(|\phi\rangle)<4, \\ \operatorname{REMOVE} E_{2}\left(|\hat{1} \hat{b}\rangle\langle\hat{1} \hat{b} \mid \phi\rangle+\sum_{y \in\{0,1\}^{4}-\{\hat{1} \hat{b}\}}|y\rangle \otimes h_{8}(\langle y \mid \phi\rangle)\right) & \text { ngược lại. }\end{cases}
$$

(iv) Cuối cùng, chúng ta thay đổi $\hat{b}$ ngoài cùng bên phải trong $|\hat{0} \hat{0} \hat{1} \hat{b}\rangle$ thành $\hat{2}$ và sau đó nhận được $|\hat{0} \hat{0} \hat{1} \hat{2}\rangle(=|\widetilde{001}\rangle)$. Để tạo ra một chuỗi lượng tử như vậy, chúng ta áp dụng hàm lượng tử $h_{6}$ được định nghĩa là

$$
h_{9}(|\phi\rangle)= \begin{cases}|\phi\rangle & \text { nếu } \ell(|\phi\rangle)<2 \\ \sum_{y \in\{\hat{0}, \hat{1}\}}|y\rangle \otimes h_{9}(\langle y \mid \phi\rangle)+|\hat{b}\rangle\langle\hat{2} \mid \phi\rangle+|\hat{2}\rangle\langle\hat{b} \mid \phi\rangle & \text { ngược lại }\end{cases}
$$

Chính thức, chúng ta đặt $h_{9}=2 Q \operatorname{Rec}_{1}\left[I, h_{9}^{\prime}, I \mid\left\{h_{y}\right\}_{y \in\{0,1\}^{2}}\right]$, nơi $h_{9}^{\prime}=\operatorname{Branch}\left[\left\{h_{y}^{\prime \prime}\right\}_{y \in\{0,1\}^{2}}\right]$ với $h_{\hat{2}}^{\prime \prime}=h_{\hat{b}}^{\prime \prime}= S W A P \circ N O T \circ S W A P$ và $h_{y}^{\prime \prime}=I$ cùng với $h_{\hat{2}}=h_{\hat{b}}=I$ và $h_{y}=h_{9}$ cho bất kỳ $y \in\{\hat{0}, \hat{1}\}$.\\
(vi) Để thực hiện các Bước (i)-(v) cùng một lúc, chúng ta kết hợp tất cả các hàm lượng tử được sử dụng trong các Bước (i)-(v) và định nghĩa một hàm lượng tử đơn $F_{4}=h_{9} \circ h_{8} \circ f_{10} \circ f_{9}$. Tổng thể, chuỗi lượng tử kết quả là $|\hat{0} \hat{0} \hat{1} \hat{2}\rangle|\hat{3} \hat{b} \hat{3} \hat{b} \hat{3} \hat{b}\rangle|\hat{1} \hat{3} \hat{3} \hat{2}\rangle$.

Trạng thái lượng tử $\left|\phi_{F_{4}}^{p}(x)\right\rangle$ có thể được biểu thị là $\sum_{r \in F S C_{M, n}}|\widetilde{M[r]}\rangle\left|\widehat{\xi}_{x, r}\right\rangle$ cho các trạng thái lượng tử nhất định $\left\{\left|\widehat{\xi}_{x, r}\right\rangle\right\}_{r}$. Vì chúng ta xử lý một vùng băng thiết yếu của $M$, nó ngay lập tức theo sau rằng, cho mọi $x, x^{\prime} \in\{0,1\}^{n}$ và $r, r^{\prime} \in F S C_{M, n},\left\|\left\langle\xi_{x, r} \mid \xi_{x^{\prime}, r^{\prime}}\right\rangle\right\|=\left\|\left\langle\widehat{\xi}_{x, r} \mid \widehat{\xi}_{x^{\prime}, r^{\prime}}\right\rangle\right\|$ và $\left\langle\hat{\xi}_{x, r} \mid \hat{\xi}_{x, r^{\prime}}\right\rangle=0$ nếu $r \neq r^{\prime}$. Do đó, $F_{4}$ thỏa mãn điều kiện của phần đầu tiên của bổ đề.

Đối với phần thứ hai của bổ đề, chúng ta cần thêm nữa truy xuất $|M[r]\rangle$ từ chuỗi lượng tử được mã hóa $\mid \widetilde{M[r]\rangle}$. Từ ví dụ minh họa trước đó $|\hat{0} \hat{0} \hat{1} \hat{2}\rangle|\hat{3} \hat{b} \hat{3} \hat{b} \hat{3} \hat{b}\rangle|\hat{1} \hat{3} \hat{3} \hat{2}\rangle$, chúng ta cần tạo ra $|001\rangle|1\rangle|\hat{3} \hat{b} \hat{b} \hat{b} \hat{3} \hat{b}\rangle|\hat{1} \hat{3} \hat{3} \hat{2}\rangle|1000\rangle$. Vì mục đích này, chúng ta định nghĩa hàm lượng tử $g_{7}$ bằng cách đặt $g_{7}=2 Q R e c_{1}\left[I, R E M O V E_{1}, I \mid\left\{g_{z}^{\prime}\right\}_{z \in\{0,1\}^{2}}\right]$; cụ thể là,

$$
g_{7}(|\phi\rangle)= \begin{cases}|\phi\rangle & \text { nếu } \ell(|\phi\rangle)<2 \\ R E M O V E_{1}\left(\sum_{y \in\{0,1\}}\left(|0 y\rangle \otimes g_{7}(\langle 0 y \mid \phi\rangle)+|1 y\rangle\langle 1 y \mid \phi\rangle\right)\right) & \text { ngược lại }\end{cases}
$$

Sử dụng $g_{7}$, cuối cùng chúng ta đặt $F_{5}=g_{7} \circ F_{4}$ để có được phần thứ hai của bổ đề.\\
Điều này hoàn thành chứng minh của Bổ đề 4.3.

\subsection*{5.6 Các ứng dụng đơn giản của các quy trình mô phỏng}
Chứng minh của Bổ đề 4.3 cung cấp các quy trình hữu ích không chỉ để xây dựng hàm lượng tử mong muốn $g$ mà còn cho các hàm lượng tử chuyên dụng khác. Ở đây, như các ứng dụng đơn giản của các quy trình mô phỏng được đưa ra trong các Bước 1)-6) của Mục 5.2, chúng ta sẽ giải thích cách mã hóa/giải mã các chuỗi cổ điển và cách nhân bản thông tin cổ điển bằng các hàm lượng tử $\widehat{\square_{1}^{\mathrm{QP}}}$.

Các Bước 1)-8) trong Mục 5.2 mô tả một biến đổi giữa các chuỗi cổ điển và mã hóa của chúng. Một sự sửa đổi nhỏ của Bước 2) giới thiệu một bộ mã hóa Encode, mã hóa đúng các chuỗi nhị phân $s$ thành $\tilde{s}$.

Bổ đề 5.1 Tồn tại một hàm lượng tử Encode trong $\widehat{\square_{1}^{\mathrm{QP}}}$ thỏa mãn Encode $\left(\left|0^{k+1} 1\right\rangle \otimes|\phi\rangle\right)= \sum_{s \in\{0,1\}^{k}}(|\tilde{s}\rangle \otimes|s\rangle\langle s \mid \phi\rangle)$ cho bất kỳ số $k \in \mathbb{N}$ và bất kỳ trạng thái lượng tử $|\phi\rangle \in \mathcal{H}_{\infty}$. Hơn nữa, hàm lượng tử Decode $=$ Encode $^{-1}$ cũng nằm trong $\widehat{\square_{1}^{\mathrm{QP}}}$.

Cho các bit bổ sung $0^{k} 1$, $k$ qubit đầu tiên của $|\phi\rangle$ được mã hóa đúng bởi Encode bằng cách tiêu thụ $0^{k} 1$. Ví dụ, Encode biến đổi $\left|0^{3} 1\right\rangle\left|a_{1} a_{2}\right\rangle$ thành $\left|\hat{a}_{1} \hat{a}_{2} \hat{2}\right\rangle$ và Decode trả lại $\left|\hat{a}_{1} \hat{a}_{2} \hat{2}\right\rangle$ về $\left|0^{3} 1\right\rangle\left|a_{1} a_{2}\right\rangle$. Lưu ý rằng Mệnh đề 3.5 đảm bảo Decode $\in \widehat{\square_{1}^{\mathrm{QP}}}$ từ Encode $\in \widehat{\square_{1}^{\mathrm{QP}}}$.

Cơ học lượng tử nói chung cấm chúng ta sao chép các trạng thái lượng tử chưa biết; tuy nhiên, có thể sao chép mỗi chuỗi cổ điển lượng tử. Do đó, chúng ta có hàm lượng tử $C O P Y_{2}$ sau, sao chép nội dung của $k$ qubit đầu tiên của bất kỳ đầu vào nào.

Bổ đề 5.2 Tồn tại một hàm lượng tử $C O P Y_{2}$ trong $\widehat{\square_{1}^{\mathrm{QP}}}$ thỏa mãn điều kiện sau: cho bất kỳ $k \in \mathbb{N}^{+}$và bất kỳ $|\phi\rangle \in \mathcal{H}_{\infty}$,

$$
C O P Y_{2}\left(\left|\widetilde{0^{k}}\right\rangle \otimes|\phi\rangle\right)=\sum_{s \in\{0,1\}^{k}}(|\tilde{s}\rangle \otimes|\tilde{s}\rangle\langle\tilde{s} \mid \phi\rangle)
$$

Việc mã hóa $\widetilde{0^{k}}$ của $0^{k}$ là cần thiết để phân biệt $0^{k}$ với bất kỳ phần nào của $|\phi\rangle$ vì $k$ không phải là một hằng số cố định. Khi $|\phi\rangle$ có dạng $|x\rangle|\psi\rangle$, hàm lượng tử $C O P Y_{2}$ hoạt động như $C O P Y_{2}\left(\left|\widetilde{0^{k}}\right\rangle \otimes|\tilde{x}\rangle|\psi\rangle\right)=|\tilde{x}\rangle \otimes|\tilde{x}\rangle|\psi\rangle$.

Chứng minh Bổ đề 5.2. Chúng ta muốn xây dựng hàm lượng tử mong muốn $C O P Y_{2}$ như sau. Để làm cho quá trình xây dựng của chúng ta dễ đọc, chúng ta sử dụng một ví dụ minh họa của $\left|\widetilde{0^{k}}\right\rangle|\phi\rangle=\left|\widetilde{0^{2}}\right\rangle\left|\widetilde{a_{1} a_{2}}\right\rangle\left(=|\hat{0} \hat{0} \hat{2}\rangle\left|\hat{a}_{1} \hat{a}_{2} \hat{2}\right\rangle\right)$ để chỉ ra cách mỗi hàm lượng tử được xây dựng hoạt động.

\begin{enumerate}
  \item Các Bước 1)-6) của Mục 5.2 biến đổi $|\hat{0} \hat{0} \hat{1}\rangle\left|\hat{a}_{1} \hat{a}_{2}\right\rangle$ thành $\left|\hat{2} \hat{a}_{1} \hat{3} \hat{a}_{2}\right\rangle|\hat{3}\rangle$. Bằng một sự sửa đổi nhỏ của các bước này, có thể biến đổi $|\hat{0} \hat{0} \hat{2}\rangle\left|\hat{a}_{1} \hat{a}_{2} \hat{2}\right\rangle$ thành $\left|\hat{3} \hat{a}_{1} \hat{3} \hat{a}_{2} \hat{2}\right\rangle|\hat{2}\rangle$. Chúng ta ký hiệu $f_{1}$ là một hàm lượng tử thực hiện biến đổi này.
  \item Bằng cách sao chép $\hat{a}_{i}$ trong $\hat{3} \hat{a}_{i}$ lên $\hat{3}$ cho mỗi $i \in\{1,2\}$, chúng ta thay đổi $\hat{3} \hat{a}_{i}$ thành $\hat{a}_{i} \hat{a}_{i}$ và sau đó nhận được $\left|\hat{a}_{1} \hat{a}_{1} \hat{a}_{2} \hat{a}_{2} \hat{2}\right\rangle|\hat{2}\rangle$. Bước này có thể được thực hiện chính thức theo cách sau. Trước tiên, chúng ta định nghĩa một hàm lượng tử $h_{2}$ thỏa mãn $h_{2}\left(\left|b_{1} b_{2}\right\rangle\left|b_{3} b_{4}\right\rangle \otimes|\phi\rangle\right)=\left|b_{4} b_{2} b_{1} b_{3}\right\rangle \otimes|\phi\rangle$ cho bất kỳ $\phi \in \mathcal{H}_{\infty}$. Một hàm lượng tử như vậy thực sự tồn tại theo Bổ đề 3.7. Thứ hai, chúng ta đặt $D U P$ là $h_{2}^{-1} \circ(S W A P \circ N O T \circ S W A P \circ C N O T) \circ h_{2} \circ N O T$. Khi đó, theo sau rằng\\
$D U P(|\hat{3}\rangle|\hat{a}\rangle \otimes|\phi\rangle)=|\hat{a}\rangle|\hat{a}\rangle \otimes|\phi\rangle$ cho bất kỳ bit $a \in\{0,1\}$ và bất kỳ trạng thái lượng tử $|\phi\rangle \in \mathcal{H}_{\infty}$. Với DUP này, chúng ta tiếp tục định nghĩa $f_{2}$ bằng phép đệ quy lượng tử 4 qubit như

$$
f_{2}(|\phi\rangle)= \begin{cases}|\phi\rangle & \text { nếu } \ell(|\phi\rangle)<4, \\ D U P\left(\sum_{a \in\{0,1\}}|\hat{3} \hat{a}\rangle \otimes f_{2}(\langle\hat{3} \hat{a} \mid \phi\rangle)+\sum_{y \in B_{4}^{\prime}}|y\rangle\langle y \mid \phi\rangle\right. & \text { ngược lại },\end{cases}
$$

nơi $B_{4}^{\prime}=\{0,1\}^{4}-\{\hat{3} \hat{0}, \hat{3} \hat{1}\}$.\\
  \item Tiếp theo, chúng ta biến đổi $\left|\hat{a}_{1} \hat{a}_{1} \hat{a}_{2} \hat{a}_{2} \hat{2}\right\rangle|\hat{2}\rangle$ thành $\left|\hat{a}_{1} \hat{a}_{2} \hat{2}\right\rangle\left|\hat{2} \hat{a}_{2} \hat{a}_{1}\right\rangle$. Biến đổi này có thể được thực hiện bởi $f_{3}$ được định nghĩa là

$$
f_{3}(|\phi\rangle)= \begin{cases}|\phi\rangle & \text { nếu } \ell(|\phi\rangle)<4, \\ R E M O V E_{2}\left(\sum_{y \in\{0,1\}^{4}-\{\hat{2} \hat{2}\}}|y\rangle \otimes f_{3}(\langle y \mid \phi\rangle)+|\hat{2} \hat{2}\rangle\langle\hat{2} \hat{2} \mid \phi\rangle\right) & \text { ngược lại. }\end{cases}
$$

  \item Sau đó, chúng ta thay đổi $\left|\hat{a}_{1} \hat{a}_{2} \hat{2}\right\rangle\left|\hat{2} \hat{a}_{2} \hat{a}_{1}\right\rangle$ thành $\left|\hat{2} \hat{a}_{2} \hat{a}_{1}\right\rangle\left|\hat{2} \hat{a}_{2} \hat{a}_{1}\right\rangle$ bằng cách loại bỏ từng hai qubit trong phần đầu tiên đến cuối. Quá trình này được hiện thực hóa chính xác bởi $f_{4}$ được định nghĩa là

$$
f_{4}(|\phi\rangle)= \begin{cases}|\phi\rangle & \text { nếu } \ell(|\phi\rangle)<4, \\ \operatorname{REMOV} E_{2}\left(\sum_{y \in\{0,1\}^{2}-\{\hat{2}\}}\left(|y\rangle \otimes f_{5}(\langle y \mid \phi\rangle)\right)+|\hat{2}\rangle\langle\hat{2} \mid \phi\rangle\right) & \text { ngược lại } .\end{cases}
$$

  \item Cuối cùng, chúng ta đảo ngược toàn bộ chuỗi lượng tử để nhận được $\left|\hat{a}_{1} \hat{a}_{2} \hat{2}\right\rangle\left|\hat{a}_{1} \hat{a}_{2} \hat{2}\right\rangle$ bằng cách áp dụng REVERSE.
\end{enumerate}

Điều này hoàn thành chứng minh của bổ đề.

\section*{6 Những thách thức trong tương lai}
Trong Định nghĩa 3.1, chúng ta đã định nghĩa các hàm lượng tử $\square_{1}^{\mathrm{QP}}$ trên $\mathcal{H}_{\infty}$ và chúng ta đã đưa ra trong Định lý 4.1 một đặc trưng mới của các hàm FBQP theo các hàm lượng tử $\square_{1}^{\mathrm{QP}}$ này. Để chỉ ra các hướng nghiên cứu trong tương lai, chúng ta muốn nêu một câu hỏi mở thách thức trong Mục 6.1 và trình bày trong các Mục 6.26 .3 ba hệ quả có thể có của định nghĩa sơ đồ của chúng ta đối với các chủ đề về độ phức tạp mô tả, lý thuyết bậc nhất và các hàm bậc cao. Trong Mục 6.4, chúng ta sẽ đề cập một ứng dụng thực tiễn cho việc thiết kế các ngôn ngữ lập trình lượng tử.

\subsection*{6.1 Tìm kiếm một định nghĩa sơ đồ hợp lý hơn}
Định nghĩa sơ đồ của chúng ta (Định nghĩa 3.1) được tạo thành từ các hàm lượng tử ban đầu, được suy ra từ các cổng lượng tử tự nhiên, đơn giản, và các quy tắc xây dựng, bao gồm phép đệ quy lượng tử nhiều qubit, làm phong phú đáng kể phạm vi các hàm lượng tử được xây dựng. Việc lựa chọn các hàm lượng tử ban đầu và các quy tắc xây dựng mà chúng ta đã sử dụng trong bài viết này trực tiếp ảnh hưởng đến sự phong phú của các hàm lượng tử $\square_{1}^{\text {QP }}$. Mặc dù $\square_{1}^{\mathrm{QP}}$ hiện tại của chúng ta là đủ để đặc trưng hóa FBQP, nếu chúng ta tìm cách làm phong phú thêm các hàm lượng tử $\square_{1}^{\mathrm{QP}}$, một cách là bổ sung thêm các hàm lượng tử ban đầu. Là một ví dụ cụ thể, hãy xem xét phép biến đổi Fourier lượng tử (QFT), đóng vai trò quan trọng trong, ví dụ, thuật toán lượng tử phân tích thừa số của Shor 30. Chúng ta đã minh họa trong Bổ đề 3.10 cách thực hiện một dạng QFT bị hạn chế hoạt động trên một số lượng qubit cố định, và do đó nó thuộc về $\widehat{\square_{1}^{\mathrm{QP}}}$. Tuy nhiên, một dạng QFT tổng quát hơn, hoạt động trên một số lượng "bất kỳ" của các qubit, có thể không được hiện thực hóa chính xác bởi các hàm lượng tử $\widehat{\square_{1}^{\mathrm{QP}}}$ mặc dù nó có thể được xấp xỉ đến bất kỳ độ chính xác mong muốn nào bởi các hàm lượng tử $\widehat{\square_{1}^{\mathrm{QP}}}$. Để khắc phục việc loại trừ QFT khỏi lớp hàm lượng tử $\square_{1}^{\mathrm{QP}}$ của chúng ta, ví dụ, chúng ta có thể mở rộng $\widehat{\square_{1}^{\mathrm{QP}}}$ hiện tại bằng cách bổ sung vào các hàm lượng tử ban đầu hàm lượng tử được định nghĩa là

$$
C R O T\left(|\phi\rangle\left|0^{j}\right\rangle\right)=|0\rangle\langle 0 \mid \phi\rangle\left|0^{j}\right\rangle+\omega_{j}|1\rangle\langle 1 \mid \phi\rangle\left|0^{j}\right\rangle \text { (phép quay có điều khiển), }
$$

nơi $\omega_{j}=e^{2 \pi i / 2^{j}}$, với một thuật ngữ bổ sung $0^{j}$. Không khó để xây dựng QFT từ các hàm lượng tử trong $\widehat{\square_{1}^{\mathrm{QP}}}$ đã được mở rộng này bằng cách bổ sung CROT. Như ví dụ này cho thấy, vẫn còn quan trọng để tìm kiếm một định nghĩa sơ đồ đơn giản, hợp lý hơn cho các hàm lượng tử, có khả năng đặc trưng chính xác cả BQP và FBQP và cũng đơn giản hóa chứng minh Định lý 4.1.

Từ quan điểm của chủ nghĩa tối giản, ngược lại, chúng ta có thể loại bỏ một số sơ đồ nhất định hoặc thay thế chúng bằng các sơ đồ đơn giản hơn nhưng vẫn đảm bảo kết quả đặc trưng của BQP và FBQP theo các hàm lượng tử $\square_{1}^{\mathrm{QP}}$. Là một ví dụ cụ thể, chúng ta có thể hỏi liệu phép đệ quy lượng tử nhiều qubit của chúng ta có thể được thay thế bằng phép đệ quy lượng tử 1 qubit với chi phí bổ sung các hàm lượng tử ban đầu.

\subsection*{6.2 Giới thiệu độ phức tạp mô tả và các lý thuyết bậc nhất}
Như đã lưu ý trong Mục 3.1, định nghĩa sơ đồ của chúng ta về các hàm lượng tử $\square_{1}^{\mathrm{QP}}$ cung cấp cho chúng ta một phương tiện tự nhiên để gán độ phức tạp mô tả - một thước đo độ phức tạp mới - cho từng hàm lượng tử trong $\square_{1}^{\mathrm{QP}}$. Thước đo độ phức tạp này đã được sử dụng để chứng minh, ví dụ, Bổ đề 3.4. Là một hệ quả của định lý chính của chúng ta, thước đo độ phức tạp này cũng chuyển sang "độ phức tạp mô tả" của các ngôn ngữ và hàm trong BQP và FBQP, và do đó nó tự nhiên giúp chúng ta giới thiệu khái niệm về độ phức tạp mô tả của các ngôn ngữ và hàm đó.

Hơn nữa, có thể mở rộng thước đo độ phức tạp này cho "bất kỳ" ngôn ngữ và hàm nào trên $\{0,1\}^{*}$, không nhất thiết bị giới hạn trong FBQP và BQP, và thảo luận về "độ phức tạp tương đối" của chúng đối với $\square_{1}^{\mathrm{QP}}$. Chính thức hơn, cho một hàm $f$ trên $\{0,1\}^{*}$, độ phức tạp mô tả $\square_{1}^{\mathrm{QP}}$ của $f$ tại độ dài $n$ là số lượng tối thiểu các lần chúng ta sử dụng các hàm lượng tử ban đầu và các quy tắc xây dựng để xây dựng một hàm lượng tử $\square_{1}^{\text {QP }}$ $g$ mà $\ell\left(\left|\phi_{g}^{p}(x)\right\rangle\right)=|f(x)|$ và $\left|\left\langle f(x) \mid \phi_{g}^{p}(x)\right\rangle\right|^{2} \geq 2 / 3$ giữ cho một đa thức $p$ nhất định với $|f(x)| \leq p(|x|)$ cho tất cả các chuỗi $x$ có độ dài chính xác $n$. Chúng ta viết $d c(f)[n]$ để biểu thị độ phức tạp mô tả $\square_{1}^{\mathrm{QP}}$ của $f$ tại độ dài $n$. Rõ ràng, mọi hàm lượng tử $\square_{1}^{\mathrm{QP}}$ đều có độ phức tạp mô tả $\square_{1}^{\mathrm{QP}}$ không đổi tại mọi độ dài. Theo tinh thần tương tự nhưng dựa trên các máy hữu hạn lượng tử, Villagra và Yamakami 33 đã thảo luận về độ phức tạp trạng thái lượng tử được giới hạn ở các đầu vào có độ dài chính xác $n$ (cũng như độ dài tối đa $n$). Hóa ra thước đo độ phức tạp như vậy rất hữu ích. Tham khảo 33 để biết các định nghĩa chi tiết. Thước đo độ phức tạp mới của chúng ta $d c(f)[n]$ cũng được kỳ vọng là một công cụ hữu ích trong việc phân loại các ngôn ngữ và hàm theo sức mạnh mô tả theo cách khá khác biệt so với những gì QTMs và các mạch lượng tử làm.

Trong một quan điểm rộng hơn, định nghĩa sơ đồ của chúng ta về khả năng tính toán lượng tử trong thời gian đa thức có thể dẫn đến sự phát triển trong tương lai của một dạng lý thuyết bậc nhất phù hợp trên các trạng thái lượng tử trong không gian Hilbert hoặc các lý thuyết lượng tử bậc nhất, nói ngắn gọn. Trong tài liệu, các lý thuyết bậc nhất và các lý thuyết con tự nhiên của chúng đã trở thành một chủ đề nghiên cứu phong phú trong logic toán học và lý thuyết đệ quy và chúng cũng đã tìm thấy nhiều ứng dụng trong các lĩnh vực khác. Một dạng yếu của các lý thuyết con của chúng đã được thảo luận trong lý thuyết độ phức tạp lượng tử. Ví dụ, sử dụng các lượng tử bị chặn trên các trạng thái lượng tử trong không gian Hilbert, các dạng lượng tử của NP và phân cấp đa thức (thời gian) Meyer-Stockmeyer đã được thảo luận trong 37. Thật không may, chúng ta vẫn còn xa mới có được các lý thuyết bậc nhất được chấp nhận rộng rãi và các lý thuyết con hữu ích cho máy tính lượng tử.

\subsection*{6.3 Mở rộng sang các hàm bậc 2}
Thông thường, các hàm ánh xạ $\Sigma^{*}$ tới $\Sigma^{*}$ được phân loại là các hàm bậc 1, trong khi các hàm bậc 2 là các hàm nhận đầu vào từ $\Sigma^{*}$ cùng với các hàm bậc 1. Trong lý thuyết độ phức tạp tính toán, các hàm bậc 2 như vậy đã được thảo luận rộng rãi trong, ví dụ, $8,9,24,31,35$.

Theo cách tương tự như trường hợp cổ điển, chúng ta gọi các hàm lượng tử $\square_{1}^{\mathrm{QP}}$ trên $\mathcal{H}_{\infty}$ là các hàm bậc 1. Để giới thiệu các hàm bậc 2, trước tiên chúng ta bắt đầu với một hàm lượng tử tùy ý $O$ ánh xạ $\mathcal{H}_{\infty}$ tới $\mathcal{H}_{\infty}$, được xử lý như một hàm oracle (một oracle hàm hoặc đơn giản là một oracle) trong công thức sau. Từ một oracle như vậy, chúng ta định nghĩa một toán tử tuyến tính mới $\tilde{O}$ là $\tilde{O}(|\tilde{x}\rangle)=\sum_{s:|s|=|x|} \alpha_{s}|\tilde{s}\rangle$ nếu $O(|x\rangle)$ có dạng $\sum_{s:|s|=|x|} \alpha_{s}|s\rangle$ cho bất kỳ chuỗi $x \in\{0,1\}^{*}$, nơi $\tilde{s}$ là mã của $s$, được định nghĩa trong Mục 4.1. Lưu ý rằng $\tilde{O}(|\tilde{x}\rangle)= \widetilde{O(|x\rangle) \text { và } \ell(\tilde{O}(|\tilde{x}\rangle))=2 \ell(O(|x\rangle))+2 \text {. }}$

Trước tiên, chúng ta mở rộng các hàm lượng tử ban đầu và các quy tắc xây dựng được đưa ra trong Định nghĩa 3.1 bằng cách thay thế bất kỳ hàm lượng tử nào, ví dụ, $f(|\phi\rangle)$ trong mỗi sơ đồ của định nghĩa bằng $f(|\phi\rangle, O)$. Thứ hai, chúng ta giới thiệu một hàm lượng tử ban đầu khác, gọi là hàm truy vấn QUERY. Cho bất kỳ hai chuỗi lượng tử $|\phi\rangle$ và $|\psi\rangle$ có độ dài $n$, cho $|\phi\rangle \oplus|\psi\rangle=\sum_{s:|s|=n} \sum_{t:|t|=n}\langle s \mid \phi\rangle\langle t \mid \psi\rangle|s \oplus t\rangle$, nơi $s \oplus t$ biểu thị XOR bit-wise của $s$ và $t$. Là một trường hợp đặc biệt, theo sau rằng $\tilde{O}(|\tilde{x}\rangle) \oplus \tilde{O}(|\tilde{x}\rangle)=\left|0^{2|x|+2}\right\rangle$ cho bất kỳ $x \in\{0,1\}^{*}$. Hàm truy vấn sau đó được định nghĩa là

$$
\begin{aligned}
& Q U E R Y(|\phi\rangle, O) \\
& \qquad=\sum_{n \in[t]} \sum_{x \in \Sigma^{n}} \sum_{s \in \Sigma^{n}}\left(|\tilde{x}\rangle \otimes(\tilde{O}(|\tilde{x}\rangle) \oplus|\tilde{s}\rangle) \otimes\langle\tilde{x} \tilde{s} \mid \phi\rangle+\sum_{y \in \Sigma^{4 n+4} \wedge y \neq \tilde{x} \tilde{s}}|y\rangle\langle y \mid \phi\rangle\right)
\end{aligned}
$$

cho tất cả $|\phi\rangle \in \mathcal{H}_{\infty}$, nơi $t=\lfloor(\ell(|\phi\rangle)-4) / 4\rfloor$. Đặc biệt, $Q U E R Y(|\tilde{x}\rangle|\tilde{s}\rangle|\phi\rangle)$ bằng $|\tilde{x}\rangle \otimes(\tilde{O}(|\tilde{x}\rangle \oplus|\tilde{s}\rangle) \otimes|\phi\rangle)$. Nó cũng theo sau rằng $Q U E R Y \circ Q U E R Y(|\phi\rangle, O)=I(|\phi\rangle)$ vì $\tilde{O}(|\tilde{x}\rangle) \oplus \tilde{O}(|\tilde{x}\rangle)=\left|0^{2|x|+2}\right\rangle$.

Lưu ý rằng nếu $O$ nằm trong $\square_{1}^{\mathrm{QP}}$ thì hàm $Q_{O}(|\phi\rangle)={ }_{\text {def }} \operatorname{QUERY}(|\phi\rangle, O)$ cũng thuộc về $\square_{1}^{\mathrm{QP}}$. Cũng có thể chứng minh các kết quả cơ bản tương tự được thảo luận trong các phần trước. Những kết quả cơ bản này có thể mở ra một lĩnh vực phong phú về khả năng tính toán lượng tử bậc cao và chúng ta mong đợi những kết quả phong phú sẽ được khám phá trong lĩnh vực mới này.

\subsection*{6.4 Ứng dụng cho các ngôn ngữ lập trình lượng tử}
Một ứng dụng thực tiễn của định nghĩa sơ đồ của chúng ta có thể được tìm thấy trong lĩnh vực các ngôn ngữ lập trình lượng tử. Kể từ những ngày đầu của nghiên cứu máy tính lượng tử, một nỗ lực đáng kể đã được thực hiện bởi các nhà vật lý, nhà khoa học máy tính và kỹ sư máy tính để vẽ ra một bản đồ nghiên cứu thực tiễn cho một máy tính lượng tử thực sự.

Hướng tới việc hiện thực hóa một máy tính lượng tử "đa năng" như vậy, tuy nhiên, điều quan trọng hơn là làm cho chúng "có thể lập trình" theo cách mà việc chạy một "chương trình lượng tử" phù hợp tự do thay đổi các quá trình tính toán cho các bài toán mục tiêu khác nhau mà không cần mô hình lại phần cứng của chúng mỗi lần. Một chương trình lượng tử ở đây đề cập đến một chuỗi hữu hạn các hướng dẫn về cách vận hành máy tính lượng tử từng bước. Để viết một chương trình lượng tử như vậy, tuy nhiên, chúng ta cần phát triển các ngôn ngữ lập trình được cấu trúc tốt cho máy tính lượng tử (gọi là các ngôn ngữ lập trình lượng tử). Các ngôn ngữ lập trình lượng tử khác nhau đã được thảo luận trong hai thập kỷ qua trong quá trình phát triển các máy tính lượng tử thực tế. Tham khảo các khảo sát, ví dụ, 13 để biết nền tảng cần thiết.

Định nghĩa sơ đồ của chúng ta cung cấp một mô tả về cách định nghĩa một hàm lượng tử $\square_{1}^{\mathrm{QP}}$ đã cho. Mô tả này có thể được xem như một tập hợp các hướng dẫn, mỗi hướng dẫn chỉ dẫn cách áp dụng từng sơ đồ để xây dựng hàm lượng tử mong muốn và do đó nó giống như một chương trình chỉ đạo cách xây dựng hàm lượng tử. Do đó, mô tả sơ đồ của chúng ta về quá trình xây dựng các hàm lượng tử có thể giúp chúng ta thiết kế các ngôn ngữ lập trình lượng tử phù hợp trong tương lai.

\section*{References}
[1] L. M. Adleman, J. DeMarrais, and M. A. Huang, Quantum computability, SIAM J. Comput. 26 (1997) 1524-1540.\\[0pt]
[2] A. Barenco, C. H. Bennett, R. Cleve, D. DiVicenzo, N. Margolus, P. Shor, T. Sleator, J. Smolin, and H. Weinfurter, Elementary gates for quantum computation, Phys. Rev., A, 52 (1995) 3457-3467.\\[0pt]
[3] P. Benioff, The computer as a physical system: a microscopic quantum mechanical Hamiltonian model of computers represented by Turing machines. J. Statist. Phys. 22 (1980) 563-591.\\[0pt]
[4] C. H. Bennett, E. Bernstein, G. Brassard, and U. Vazirani, Strengths and weaknesses of quantum computing, SIAM J. Comput. 26 (1997) 1510-1523.\\[0pt]
[5] E. Bernstein and U. Vazirani, Quantum complexity theory, SIAM J. Comput. 26 (1997) 1411-1473.\\[0pt]
[6] P. O. Boykin, T. Mor, M. Pulver, V. Roychowdhury, and F. Vatan, On universal and fault-tolerant quantum computing. arXiv:quant-ph/9906054, 1999.\\[0pt]
[7] A. Cobham, The intrinsic computational difficulty of functions, in the Proceedings of the 1964 Congress for Logic, Mathematics, and Philosophy of Science, pp. 24-30, North-Holland, 1964.\\[0pt]
[8] R. L. Constable, Type two computational complexity, in the Proceedings of the 5th ACM Symposium on Theory of Computing (STOC'73), pp. 108-121, 1973.\\[0pt]
[9] S. Cook and B. M. Kapron, Characterizations of the basic feasible functionals of finite type, in the Proceedings of the 30th IEEE Conference on Foundations of Computer Science (STOC'89), pp. 154-159, 1989.\\[0pt]
[10] M. Davis, (ed.) The Unidecidable. Basic Papers on Undecidable Propositions, Unsolvable Problems, and Computable Functions, Raven Press, Hewlett, New York, 1965.\\[0pt]
[11] D. Deutsch, Quantum theory, the Church-Turing principle, and the universal quantum computer, Proceedings Royal Society London, Ser. A, 400 (1985) 97-117.\\[0pt]
[12] D. Deutsch, Quantum computational networks, Proceedings Royal Society London, Ser. A, 425 (1989) 73-90.\\[0pt]
[13] S. J. Gay. Quantum programming languages: survey and bibliography. Math. Struct. in Comp. Science 16 (2006) 581-600.\\[0pt]
[14] L. K. Grover, A fast quantum mechanical algorithm for database search, in the Proceedings of the 28th Annual ACM Symposium on the Theory of Computing (STOC'96), pp. 212-219, 1996.\\[0pt]
[15] L. Grover. Quantum mechanics in searching for a needle in a haystack. Physical Review Letters 79:325 (1997).\\[0pt]
[16] J. Hopcroft and J. Ullman, An Introduction to Automata Theory, Languages and Computation, AddisonWesley, Reading MA, 1979.\\[0pt]
[17] A. Kitaev. Quantum computations: algorithms and error correction. Russian Math. Surveys 52 (1997) 1191-1249.\\[0pt]
[18] A. Y. Kitaev, A. H. Shen, and M. N. Vyalyi. Classical and Quantum Computation (Graduate Studies in Mathematics), American Mathematical Society, 2002.\\[0pt]
[19] S. C. Kleene, General recursive functions of natural numbers, Math. Ann. 112 (1936) 727-742.\\[0pt]
[20] S. C. Kleene, Recursive predicates and quantifiers, Trans. A. M. S. 53 (1943) 41-73.\\[0pt]
[21] S. C. Kleene, Recursive functionals and quantifiers of finite types I, Trans. Amer. Math. Soc. 91 (1959) 1-52.\\[0pt]
[22] S. C. Kleene, Recursive functionals and quantifiers of finite types II, Trans. Amer. Math. Soc. 108 (1963) 106-142.\\[0pt]
[23] K. Ko and H. Friedman, Computational complexity of real functions, Theor. Comput. Sci. 20 (1982) 323352.\\[0pt]
[24] K. Mehlhorn, Polynomial and abstract subrecursive classes, J. Comput. System Sci. 12 (1976) 147-178.\\[0pt]
[25] M. A. Nielsen and I. L. Chuang, Quantum Computation and Quantum Information, Cambridge University Press, 2000.\\[0pt]
[26] H. Nishimura and M. Ozawa, Computational complexity of uniform quantum circuit families and quantum Turing machines. Theor. Comput. Sci. 276 (2002) 147-181.\\[0pt]
[27] M. Ozawa and H. Nishimura, Local transition functions of quantum Turing machines, RAIRO - Theoretical Informatics and Applications 276 (2000) 379-402.\\[0pt]
[28] G. Peano, Sul concetto di numero, Rivista di Matematica 1 (1891) 87-102, 256-267.\\[0pt]
[29] R. I. Soare, Computability and recursion. The Bulletin of Symbolic Logic, vol. 2, No. 3 (1996), pp.284-321.\\[0pt]
[30] P. W. Shor, Polynomial-time algorithms for prime factorization and discrete logarithms on a quantum computer, SIAM J. Comput. 26 (1997) 1484-1509.\\[0pt]
[31] M. Townsend, Complexity for type-2 relations, Notre Dame J. Formal Logic 31 (1990) 241-262.\\[0pt]
[32] A. M. Turing, On computable numbers with an application to the Entscheidungsproblem. Proc. London Math. Soc. ver. 2, 42 (1936) 230-265. Errutum ibid 43 (1937) 544-546.\\[0pt]
[33] M. Villagra and T. Yamakami. Quantum state complexity of formal languages, Proceedings of the 17th International Workshop on Descriptional Complexity of Formal Systems (DFCS 2015), Springer, Lecture Notes in Computer Science, vol. 9118, pp. 280-291, 2015.\\[0pt]
[34] T. Yamakami. Structural properties for feasiblely computable classes of type two. Mathematical Systems Theory 25 (1992) 177-201.\\[0pt]
[35] T. Yamakami, Feasible computability and resource bounded topology, Inf. Comput. 116 (1995) 214-230.\\[0pt]
[36] T. Yamakami, A foundation of programming a multi-tape quantum Turing machine, in the Proc. 24th Mathematical Foundations of Computer Science (MFCS'99), Lecture Notes in Computer Science, vol. 1672, pp. 430-441, 1999. See also arXiv:quant-ph/9906084.\\[0pt]
[37] T. Yamakami. Quantum NP and a quantum hierarchy. In the Proc. of the 2nd IFIP International Conference on Theoretical Computer Science (TCS 2002), Kluwer Academic Press (under the name Foundations of Information Technology in the Era of Network and Mobile Computing), the series IFIPThe International Federation for Information Processing, vol. 96 (Track 1), pp. 323-336, 2002.\\[0pt]
[38] T. Yamakami, Analysis of quantum functions, Int. J. Found. Comput. Sci. 14 (2003) 815-852. A preliminary version appeared in the Proc. 19th International Conference on Foundations of Software Technology and Theoretical Computer Science, Lecture Notes in Computer Science, Vol. 1738, pp. 407-419, 1999.\\[0pt]
[39] T. Yamakami, A recursive definition of quantum polynomial time computability (extended abstract), in the Proc. of the 9th Workshop on Non-Classical Models of Automata and Applications (NCMA 2017), Österreichische Computer Gesellschaft 2017, the Austrian Computer Society, pp. 243-258, 2017.\\[0pt]
[40] A. C. Yao, Quantum circuit complexity, in the Proc. of the 34th Annual IEEE Symposium on Foundations of Computer Science (FOCS'93), pp. 80-91, 1993.

\end{document}

